% note.tex — 標準模型の NLO の物理量：実の放射を含む崩壊幅と電弱の精密測定（MS-bar、FJ のタドポール、一般の R_xi ゲージ、d = 4）
% コンパイル: lualatex note.tex（目次と参照の解決のため2〜3回）
\documentclass[a4paper,10pt]{ltjsarticle}
\usepackage[margin=20mm]{geometry}
\usepackage[haranoaji,deluxe]{luatexja-preset}
\ltjsetparameter{jacharrange={-3}}
\usepackage{amsmath,amssymb,slashed}
\usepackage{booktabs,tabularx,array}
\usepackage[table]{xcolor}
\usepackage{tikz,ifthen}
\usetikzlibrary{arrows.meta,calc,positioning,decorations.markings,decorations.pathmorphing}
\usepackage[font=small,labelfont=bf]{caption}
\usepackage{graphicx}
\usepackage{placeins}
\usepackage{xurl}
\usepackage{hyperref}
\hypersetup{colorlinks=true,linkcolor=blue!50!black,citecolor=green!40!black,urlcolor=blue!60!black,
  pdftitle={NLO observables of the Standard Model: decay widths with real emission and electroweak precision observables in the MS-bar scheme with Fleischer-Jegerlehner tadpoles},pdfauthor={}}

\setlength{\emergencystretch}{3em}
\renewcommand{\jsParagraphMark}{\textcolor{black!65}{\rule[-0.02\zw]{0.62\zw}{0.62\zw}}\hspace{0.35\zw}}
\newcommand{\dd}{\mathrm{d}}
\newcommand{\MSbar}{\overline{\mathrm{MS}}}
\newcommand{\DUV}{\Delta_{\mathrm{UV}}}
\newcommand{\tr}{\mathrm{tr}}
\newcommand{\sw}{s_W}
\newcommand{\cw}{c_W}
\newcommand{\re}{\mathrm{Re}\,}
\newcommand{\im}{\mathrm{Im}\,}

\tikzset{
  photon/.style={thick,decorate,decoration={snake,amplitude=1.1pt,segment length=4.6pt}},
  gluon/.style={thick,decorate,decoration={coil,amplitude=2.3pt,segment length=3.6pt,aspect=0.75}},
  fermion/.style={thick,postaction={decorate},
    decoration={markings,mark=at position 0.58 with {\arrow{Stealth[length=1.9mm]}}}},
  scalar/.style={thick,dashed},
  ghost/.style={thick,densely dotted,postaction={decorate},
    decoration={markings,mark=at position 0.58 with {\arrow{Stealth[length=1.9mm]}}}},
  vtx/.style={circle,fill=black,inner sep=1.3pt},
  lab/.style={font=\small,align=center},
  slab/.style={font=\scriptsize,align=center},
  mom/.style={-{Stealth[length=1.6mm]},thin},
  ctv/.style={circle,draw=black,fill=white,inner sep=0pt,minimum size=7pt,path picture={
      \draw (path picture bounding box.north west) -- (path picture bounding box.south east)
            (path picture bounding box.south west) -- (path picture bounding box.north east);}},
  blob/.style={circle,draw=black,thick,fill=black!12,inner sep=0pt},
}

\title{標準模型の NLO の物理量：実の放射を含む崩壊幅と電弱の精密測定\\（$\MSbar$、FJ のタドポール、一般の $R_\xi$ ゲージ、$d=4$）}
\author{}
\date{2026年10月19日}

\begin{document}
\maketitle

\begin{abstract}
標準模型のノート~\cite{smnote}（以下 SM ノート）と CKM 行列のノート~\cite{ckmnote}（以下 CKM ノート）は、$\MSbar$ のパラメータと Fleischer--Jegerlehner（FJ）のタドポールの処方で、一般の $R_\xi$ ゲージの1ループの振幅が紫外で有限でゲージによらないことを確かめ、MadLoop の仮想の補正と赤外の極まで突き合わせた。ただし、どちらも実の光子やグルーオンの放射を足していないので、赤外で有限な物理量（崩壊幅、断面積）はまだ一つも求めていない。本ノートでは、CKM ノートの模型ファイル（\texttt{SM\_CKM.fr}）と同じ処方、道具で、実の放射を含む NLO の崩壊幅と、$Z$ 極の電弱の精密測定の量を求める。対象は $Z\to f\bar f$、$W^+\to f\bar f'$、$t\to bW^+$ の幅、$Z$ の有効結合と $\sin^2\theta_{\mathrm{eff}}$、$R_\ell$、$R_b$、$\rho$ パラメータと $M_W$ の予言である。手計算では、$\MSbar$ の結合と極質量の運動学で書いた木の幅、仮想と実の放射の赤外の極の構造（質量0の放射体の二重極と、質量のある二つの放射体の軟い極）、質量0のフェルミオン対への崩壊の包括的な補正 $\frac34C\alpha/\pi$（QCD では $\alpha_s/\pi$）、トップとボトムの二重項の $\Delta\rho=3y_t^2/(2K)$（$K=(4\pi)^2$。FJ のタドポールは打ち消す）を述べ、八つの手順の計画を定めた。手順 (1) では、CKM ノートの模型ファイルから書き出した FeynArts の模型と UFO が CKM ノートのものと同一であること、2体の崩壊の振幅の2乗が19個とも手計算の式と記号として一致すること、$\MSbar$ の結合と極質量の運動学をそろえた MadGraph の27個の部分幅が式と相対 $2\times10^{-15}$ で一致することを確かめた。木の $\Gamma_Z=2.428$~GeV は測定値より $2.7\%$ 小さく、差はほぼハドロンの幅の QCD の補正の大きさである。手順 (2) では、質量0のフェルミオン対の実の放射と仮想の補正を $\epsilon$ について厳密に求め、極が記号として消えて $\frac34C\alpha/\pi$ が残ることを確かめた。放射する粒子に質量を与えた $d=4$ の独立な計算でも、和は同じ値に $10^{-14}$ まで近づいた。MadGraph の木の3体の振幅の2乗も、式と相対 $6\times10^{-14}$ で一致した。計画の式の有限部分は規格化を定めておらず、$\pm\pi^2$ が抜けていたので直した。この補正を木の幅に掛けると、$\Gamma_Z$ と測定値の差は $-0.16\%$ に、ハドロンの幅の差は $-0.29\%$ に縮む。手順 (3) では、質量のある放射体（$Z\to b\bar b$、$c\bar c$、$\tau^+\tau^-$、$\mu^+\mu^-$ と $t\to bW^+$）の仮想の補正を Package-X で解析的に、軟らかい放射を $d$ 次元の eikonal 因子で、硬い放射を数値で求めた。仮想の補正の単極は式の軟らかい極と一致し、和は切断 $\delta$ によらなかった（$O(\delta)$ の誤差を除く）。質量0、閾値、$m_b=M_W=0$ の極限は既知の値に戻った。MadLoop の仮想の補正とは $7\times10^{-15}$、MadGraph の木の3体とは $8\times10^{-12}$ で一致した。$b$ の質量を入れると $\Gamma_Z$ と測定値の差は $-0.07\%$ になり、$R_b$ は $+1.1\%$ 離れる。トップの幅の QCD の補正は $-8.5\%$ で、$\Gamma_t=1.330$~GeV である。手順 (4a) では、$Z\to f\bar f$ の弱い補正（光子とグルーオンを除くループ、LSZ の因子、$Z$--光子の混合、$\MSbar$ の対抗項）を一般の $R_\xi$ ゲージで求めた。11個のフェルミオンで紫外の極は記号として消え、結果はゲージの母数によらなかった。MadLoop の電弱の仮想の補正とは $4\times10^{-13}$ で一致した。$Zb\bar b$ の頂点のトップの補正は $-1.44\%$ で、これで $R_b$ は測定値と $0.3\sigma$ で合う。$\mu$ への依存は、電弱の部分で $\pm1\%$ から $0.04\%$ 以下になった。手順 (4b) では、$W^+\to f\bar f'$ の電弱の補正を、光子が $W$ にも付くので光子を含めて求め、実の光子の放射と足した。仮想の補正の単極は $W$ の電荷を入れた軟らかい光子の極と打ち消し、結果はゲージの母数によらなかった。MadLoop の仮想の補正とは $6\times10^{-12}$、MadGraph の木の3体とは $8\times10^{-10}$ で一致した（MadGraph が $d$ 型のクォークの光子の結合を7桁で書き出す分を直したもの）。$e$ と $\mu$ の補正は $10^{-7}$ で一致し、レプトンの質量によらない（KLN）。$\mu=M_W$ の $\MSbar$ の結合では電弱の補正は $+0.01\%$（レプトン）、$-0.04\%$（軽いクォーク）と小さく、$\Gamma_W=2.086$~GeV、$\mathrm{BR}(W\to\ell\nu)=10.84\%$ で、測定値と $-1.1\sigma$、$-0.2\sigma$ である。手順 (4c) では、光子が崩壊するトップ、$b$、$W$ の三つに付く $t\to bW^+$ を同じ方法で扱った。極の打ち消しとゲージによらないことを確かめ、MadLoop とは $4\times10^{-14}$、MadGraph の木の3体とは $2\times10^{-10}$ で一致した。$m_b$ を $1/10$ にしても補正は $4\times10^{-5}$ しか変わらない（KLN）。電弱の補正は $\mu=m_t$ で $+3.28\%$（$G_F$ で書いた木に対しては $+1.52\%$）で、QCD の $-8.51\%$ と合わせて $\Gamma_t=1.378$~GeV になり、測定値と $-0.3\sigma$ で合う。手順 (5) では、(4a) の形状因子から $Z$ の有効結合と $\sin^2\theta_{\mathrm{eff}}$ を求めた。値はゲージによらず、$\mu$ を $M_Z/2$ から $2M_Z$ まで変えても $\pm3\times10^{-5}$ しか変わらない。$\sin^2\theta^{\mathrm{lept}}_{\mathrm{eff}}$ は本ノートの入力（$M_W$ の処方）で $0.23242$（測定値と $+5.6\sigma$）、$\alpha(0)$ を入力にした結合で $0.23167$（$+0.9\sigma$）で、処方の差 $7.5\times10^{-4}$ は $\Delta\rho$ の $O(\alpha\alpha_s)$ の補正の大きさである。この量は1ループの精度では測定値に届かない。手順 (6) では、トップとボトムの $\Delta\rho$ が自己エネルギーから $3y_t^2/(2K)$（$m_b$ を含めた式）に $10^{-16}$ で戻り、タドポールが消えることを確かめた。$\alpha(0)$、$G_F$、$M_Z$ からの $M_W$ の予言は $80.368$~GeV（極質量。Breit--Wigner で $80.395$~GeV）で、ゲージによらず、$\mu$ による変化は木の $528$~MeV から $34$~MeV に縮んだ。測定値とは $+1.9\sigma$（PDG 2024）で、$\Delta\rho$ の $O(\alpha\alpha_s)$ の補正（$-51$~MeV）より小さい差である。手順 (7) では、同じ1ループの幅を on-shell の $G_\mu$ の処方（$G_F$、$M_W$、$M_Z$ で書いた木）で組み直した。換算はゲージによらず、$G_\mu$ の結果の $\mu$ による変化は $0.02$〜$0.2\%$ で、$\MSbar$ の結果との差（$Z\to\nu\bar\nu$、$W\to e\nu$ で $10^{-4}$、$t\to bW^+$ で $5\times10^{-4}$、$Z\to e^+e^-$ で $1.2\times10^{-3}$）は $O(\alpha^2)$ の項で説明できた。手順 (8) では、$\mu$ の依存、処方の差、既知の高次の項から理論の不定性を見積もり、23個の予言を測定値と比べた。20個は $1.5\sigma$ 以内で合い、残る見えない幅、$\mathrm{BR}(W\to\tau\nu)$、$A_\ell$ のずれは測定のよく知られた性質である。$\sin^2\theta_{\mathrm{eff}}$ と $M_W$ では理論の不定性が実験の誤差の4〜7倍あり、2ループの補正が要る。
\end{abstract}

\tableofcontents

\section{目的と問い}\label{sec:intro}

本ノートの問いは三つである。\textbf{第一に、$\MSbar$ と FJ のタドポールの処方の1ループの仮想の補正の赤外の発散は、実の放射と足すと打ち消し、ゲージによらない有限の崩壊幅になるか。第二に、SM ノートで測定値から決めた $\MSbar$ のパラメータで、$Z$、$W$、$t$ の NLO の崩壊幅はいくらになり、電弱と QCD の補正はそれぞれどれだけか。第三に、$Z$ 極の電弱の精密測定の量（有効結合、$\sin^2\theta_{\mathrm{eff}}$、部分幅の比）と $M_W$ は、本ノートの1ループの計算でいくらになるか。}

第一の問いを立てる理由は、SM ノートと CKM ノートが仮想の補正の振幅だけを扱ったからである。紫外の極が消えること、$\xi$ によらないこと、赤外の単極の係数が軟い光子の式に一致することは確かめたが、赤外の極は実の放射と足して初めて消える（Bloch--Nordsieck、Kinoshita--Lee--Nauenberg~\cite{bn37,kinoshita62,ln64}）。質量0のフェルミオンでは軟らかい極と共線の極が重なった二重極になり、質量のあるフェルミオンでは軟らかい単極だけが残る。どちらの場合も、仮想と実の放射の極が打ち消すことを、本ノートの模型と道具で具体的に確かめる。

第二の問いを立てる理由は、両ノートの検算が物理量の予言にまだつながっていないからである。SM ノートの手順 (7) で、測定値（$G_F$、$M_Z$、$M_W$、$M_H$、$m_t$、$\alpha(0)$、$\alpha_s$）から1ループの換算で $\MSbar$ のパラメータを決めた。そのパラメータで、測られていない量（幅、比）を予言して測定値と比べることが、繰り込みの処方の一貫性の最終的な確かめになる。$\MSbar$ の結合と、崩壊する粒子の極質量の運動学がどう組み合わさるか（外線の LSZ の因子が何を吸収するか）も、ここで具体的にする。

第三の問いは、LEP と SLD の $Z$ 極の測定~\cite{lepewwg06}を本ノートの模型で読み直すものである。$\sin^2\theta_{\mathrm{eff}}$ は $\MSbar$ の $\hat s^2(\mu)$ に1ループの補正 $\hat\kappa$ を掛けたもので、その補正は $Z$ と光子の混合と頂点の補正からなる。$R_b$ には $Z b\bar b$ の頂点のトップの補正が入り、CKM ノートの $V_{tb}$ が現れる。$\rho$ パラメータのトップの寄与は、SM ノートが扱わなかった custodial 対称性の破れの量である。

求める量は五つである。(i) $\MSbar$ の結合と極質量の運動学で書いた木の崩壊幅。(ii) 仮想と実の放射の赤外の極の打ち消しと、包括的な QED と QCD の補正（質量0と質量のあるフェルミオン）。(iii) $Z\to f\bar f$、$W^+\to f\bar f'$、$t\to bW^+$ の電弱と QCD の NLO の幅。(iv) $Z$ の有効結合、$\sin^2\theta_{\mathrm{eff}}^\ell$、$\Gamma_Z$、$R_\ell$、$R_b$、$\sigma^0_{\mathrm{had}}$。(v) $\Delta\rho$ と、$G_F$ からの $M_W$ の予言。SM ノートと CKM ノートの結果（自己エネルギー、頂点、対抗項、$\MSbar$ のパラメータ、UFO）は求め直さず、引用して使う。一致の基準は SM ノートと同じである。

本ノートの構成は次の通りである。第~\ref{sec:setup} 節で模型、処方、入力と量の定義を定める。第~\ref{sec:results} 節で手計算の結果を述べる。第~\ref{sec:plan} 節でそれらをツールと突き合わせる計画を定める。第~\ref{sec:checks} 節に各手順の結果を、第~\ref{sec:discussion} 節にまとめを書く。導出は\ref{app:lo}、\ref{app:ir} にある。

\section{設定}\label{sec:setup}

\paragraph{模型と処方。}模型は CKM ノートの模型ファイル \texttt{SM\_CKM.fr} を \texttt{Src/} に写したものである（中身は同じで、頭の注記と読み込む場所だけを替えた）。SM ノートの \texttt{SM.fr} の湯川結合を世代の $3\times3$ の行列にした最新の模型で、CKM 行列 $V$ が $W$ と $G^\pm$ のクォークの結合に入り（$W^+\to u_i\bar d_j$、$R_b$ の $V_{tb}$）、$V=1$ では \texttt{SM.fr} の頂点に戻る（CKM ノートの手順 (1)）。パラメータは破れていないラグランジアンの $\MSbar$ のもの（$g'$、$g$、$g_s$、$\lambda$、$\mu^2$、湯川結合）で、タドポールは FJ の処方で扱う。ゲージは一般の $R_\xi$（$\xi_A,\xi_Z,\xi_W,\xi_G$）、次元正則化は $d=4-2\epsilon$ で、紫外と赤外の発散をともに次元正則化で扱う（検算には光子とグルーオンの質量 $\lambda$ も使う）。$\gamma_5$ は SM ノートと同じく反可換に扱う。$K=(4\pi)^2$ と書く。

\paragraph{入力。}パラメータは SM ノートの手順 (7) の値（PDG 2024~\cite{pdg24}、$M_W$ の処方、$\mu=M_t=172.57$~GeV で $g'=0.35936$、$g=0.64724$、$g_s=1.16535$、$\lambda=0.12785$、$v_{\mathrm{FJ}}=261.886$~GeV、純粋な QCD の補正を3ループまで入れた $y_t=0.93298$）である。$y_b$、$y_\tau$ は SM ノートの手順 (8) と同じく Martin--Robertson~\cite{mr19} の $M_t$ の値、$V$ は CKM ノートの手順 (6) の $\MSbar$ の値を使う。これらを SM ノートと CKM ノートの1ループの方程式で、崩壊する粒子の質量の近くの $\mu$ に走らせて使う（既定は $\mu=M$。$\mu$ を変えた変化は NLO の不定性の目安にする）。崩壊する粒子と、運動学に要る終状態の粒子の質量は極質量（PDG 2024 の $M_Z=91.1876$、$M_W=80.3692$、$M_H=125.20$、$m_t=172.57$~GeV など）とする。

\paragraph{崩壊幅の定義。}崩壊する粒子は安定とみなし（幅の効果 $\Gamma/M$ は扱わない）、その運動量を $p^2=M^2$（極質量）に置く。幅は
\begin{equation}
\Gamma=\frac1{2M}\int\dd\Phi\,\overline{|\mathcal M|^2},\qquad \Gamma_{\mathrm{NLO}}=\Gamma_0+\Gamma_V+\Gamma_R
\label{eq:gammadef}
\end{equation}
で、$\Gamma_0$ は木、$\Gamma_V=\frac1{2M}\int\dd\Phi_2\,2\re(\mathcal M_0^*\mathcal M_1)$ は繰り込んだ1ループの仮想の補正、$\Gamma_R$ は光子またはグルーオンを1個放射する3体の崩壊である。$\mathcal M_1$ は SM ノートの頂点の補正と対抗項に、外線の LSZ の因子（繰り込んだ自己エネルギーの極での留数）を掛けたものである。$\MSbar$ の結合で書いた頂点に、極質量で置いた外線の留数を掛けるので、極質量と $\MSbar$ の質量の差は留数に吸収され、別の質量の補正は要らない。

\section{結果}\label{sec:results}

\subsection{木の崩壊幅}\label{sec:lo}

$Z$ の $f\bar f$ への結合を $-i\frac{g_Z}2\gamma^\mu(v_f-a_f\gamma_5)$、$g_Z=\sqrt{g^2+g'^2}$、$v_f=T^3_f-2Q_f\hat s^2$、$a_f=T^3_f$、$\hat s^2=g'^2/(g^2+g'^2)$ と書くと、木の幅は（\ref{app:lo}）
\begin{equation}
\begin{aligned}
\Gamma_0(Z\to f\bar f)&=N_c\frac{g_Z^2M_Z}{48\pi}\beta_f\Bigl[\frac{3-\beta_f^2}2v_f^2+\beta_f^2a_f^2\Bigr],\qquad \beta_f=\sqrt{1-4m_f^2/M_Z^2},\\
\Gamma_0(W^+\to f\bar f')&=N_c|V_{ff'}|^2\frac{g^2M_W}{48\pi}\lambda^{1/2}(1,x,x')\Bigl[1-\frac{x+x'}2-\frac{(x-x')^2}2\Bigr],\qquad x=\frac{m_f^2}{M_W^2},\ x'=\frac{m_{f'}^2}{M_W^2},\\
\Gamma_0(t\to bW^+)&=|V_{tb}|^2\frac{g^2m_t^3}{64\pi M_W^2}\lambda^{1/2}(1,x_b,x_W)\bigl[(1-x_b)^2+x_W(1+x_b)-2x_W^2\bigr],\qquad x_b=\frac{m_b^2}{m_t^2},\ x_W=\frac{M_W^2}{m_t^2},\\
\Gamma_0(H\to f\bar f)&=N_c\frac{y_f^2M_H}{16\pi}\beta_f^3
\end{aligned}
\label{eq:lo}
\end{equation}
である（$\lambda(a,b,c)=a^2+b^2+c^2-2ab-2bc-2ca$、レプトンでは $N_c=1$）。結合 $g$、$g'$、$y_f$ は $\mu$ での $\MSbar$ の値、$M_Z$、$M_W$、$m_t$、$M_H$ と終状態の質量は極質量である。$t\to bW^+$ の $M_W^{-2}$ は外線の $W$ の偏極の和から来るので、極質量である。質量0の $f$ では $\Gamma_0(Z\to f\bar f)=N_cg_Z^2M_Z(v_f^2+a_f^2)/(48\pi)$ で、木の関係 $G_F/\sqrt2=g_Z^2/(8M_Z^2)$ を使うと、よく知られた $N_cG_FM_Z^3(v_f^2+a_f^2)/(6\sqrt2\pi)$ になる（$Z\to\nu\bar\nu$ で $G_FM_Z^3/(12\sqrt2\pi)$）。CKM ノートの式 (6) と SM ノートの式 (11) は、ここで結合を $\MSbar$、質量を極質量と読んだものである。木の図を図~\ref{fig:lo} に示す。

\begin{figure}[tbp]
\centering
\begin{tikzpicture}[scale=0.9]
% (a) Z -> f fbar
\begin{scope}[shift={(0,0)}]
  \draw[photon] (-1.2,0) -- (0,0); \draw[fermion] (0,0) -- (1.0,0.7); \draw[fermion] (1.0,-0.7) -- (0,0);
  \node[vtx] at (0,0) {};
  \node[font=\scriptsize] at (-1.0,0.25) {$Z$}; \node[font=\scriptsize] at (1.2,0.8) {$f$}; \node[font=\scriptsize] at (1.2,-0.8) {$\bar f$};
  \node[slab] at (0,-1.35) {(a) $-i\frac{g_Z}2\gamma^\mu(v_f-a_f\gamma_5)$};
\end{scope}
% (b) W+ -> f fbar'
\begin{scope}[shift={(3.9,0)}]
  \draw[photon] (-1.2,0) -- (0,0); \draw[fermion] (0,0) -- (1.0,0.7); \draw[fermion] (1.0,-0.7) -- (0,0);
  \node[vtx] at (0,0) {};
  \node[font=\scriptsize] at (-1.0,0.25) {$W^+$}; \node[font=\scriptsize] at (1.2,0.8) {$u_i$}; \node[font=\scriptsize] at (1.2,-0.8) {$\bar d_j$};
  \node[slab] at (0,-1.35) {(b) $-i\frac g{\sqrt2}V_{ij}\gamma^\mu P_L$\\（レプトンでは $V\to1$）};
\end{scope}
% (c) t -> b W+
\begin{scope}[shift={(7.8,0)}]
  \draw[fermion] (-1.2,0) -- (0,0); \draw[fermion] (0,0) -- (1.0,0.7); \draw[photon] (0,0) -- (1.0,-0.7);
  \node[vtx] at (0,0) {};
  \node[font=\scriptsize] at (-1.0,0.25) {$t$}; \node[font=\scriptsize] at (1.2,0.8) {$b$}; \node[font=\scriptsize] at (1.25,-0.8) {$W^+$};
  \node[slab] at (0,-1.35) {(c) $-i\frac g{\sqrt2}V_{tb}\gamma^\mu P_L$};
\end{scope}
% (d) H -> f fbar
\begin{scope}[shift={(11.7,0)}]
  \draw[scalar] (-1.2,0) -- (0,0); \draw[fermion] (0,0) -- (1.0,0.7); \draw[fermion] (1.0,-0.7) -- (0,0);
  \node[vtx] at (0,0) {};
  \node[font=\scriptsize] at (-1.0,0.25) {$H$}; \node[font=\scriptsize] at (1.2,0.8) {$f$}; \node[font=\scriptsize] at (1.2,-0.8) {$\bar f$};
  \node[slab] at (0,-1.35) {(d) $-i\frac{y_f}{\sqrt2}$};
\end{scope}
\end{tikzpicture}
\caption{手順 (1) の木の崩壊（式~\eqref{eq:lo}）と頂点の因子（\texttt{SM\_CKM.fr} の Feynman 則。$g_Z=\sqrt{g^2+g'^2}$、$v_f=T_3-2Q_f\hat s^2$、$a_f=T_3$、$P_L=(1-\gamma_5)/2$）。矢印はフェルミオンの流れ、波線はゲージ場、破線はヒッグス場。結合は $\mu=M$ の $\MSbar$ の値、運動学は極質量である。}
\label{fig:lo}
\end{figure}

\subsection{赤外の発散の構造}\label{sec:ir}

\paragraph{質量のある放射体。}色を持たない粒子が質量 $m_1$、$m_2$ のクォーク対に崩壊するとき（QED では電荷を持つフェルミオン対）、仮想の補正の赤外の極は単極だけで、CKM ノートの手順 (9)(e) で確かめた軟い極
\begin{equation}
\frac{\Gamma_V}{\Gamma_0}\bigg|_{1/\epsilon}=\frac{\alpha_sC_F}\pi\Bigl[-1-\frac1{2v}\ln\frac{1-v}{1+v}\Bigr],\qquad v=\sqrt{1-\frac{m_1^2m_2^2}{(p_1\cdot p_2)^2}}
\label{eq:softpole}
\end{equation}
になる（Catani--Dittmaier--Trocs\'anyi~\cite{cdt01}。QED では $\alpha_sC_F\to\alpha Q_fQ_{\bar f}$ の符号を合わせた形）。実の放射の振幅は、放射の運動量 $k\to0$ で木の振幅に eikonal 因子
\begin{equation}
\overline{|\mathcal M_R|^2}\to g_s^2C_F\Bigl[\frac{2p_1\cdot p_2}{(p_1\cdot k)(p_2\cdot k)}-\frac{m_1^2}{(p_1\cdot k)^2}-\frac{m_2^2}{(p_2\cdot k)^2}\Bigr]\overline{|\mathcal M_0|^2}
\label{eq:eikonal}
\end{equation}
を掛けたものになり、その位相空間の積分の赤外の極が式~\eqref{eq:softpole} と符号が逆になって、$\Gamma_V+\Gamma_R$ の極が消える。共線の発散は質量が切るので、$\ln(m^2/M^2)$ が有限部分に残るが、包括的な幅では仮想と実の放射の間で打ち消し、$m\to0$ で有限になる（KLN）。

\paragraph{質量0の放射体。}質量0のフェルミオン対では、軟らかい極と共線の極が重なり、二重極になる。$\Gamma_0$ を $d$ 次元の木の幅とし、共通の因子を
$N_\epsilon=\frac{\alpha C}{2\pi}\bigl(\frac{4\pi\mu^2}{M^2}\bigr)^\epsilon\frac1{\Gamma(1-\epsilon)}$ とすると~\cite{esw96}
\begin{equation}
\Gamma_V=\Gamma_0N_\epsilon\Bigl(-\frac2{\epsilon^2}-\frac3\epsilon-8+\pi^2\Bigr),\qquad
\Gamma_R=\Gamma_0N_\epsilon\Bigl(\frac2{\epsilon^2}+\frac3\epsilon+\frac{19}2-\pi^2\Bigr)
\label{eq:masslessVR}
\end{equation}
（$O(\epsilon)$ を除く。$\epsilon$ について厳密な形は \ref{app:ir}）で、和は
\begin{equation}
\Gamma=\Gamma_0\Bigl(1+\frac34\frac{\alpha C}\pi\Bigr),\qquad
C=C_F=\tfrac43\ (\text{グルーオン}：1+\tfrac{\alpha_s}\pi),\qquad C=Q_f^2\ (\text{$Z\to f\bar f$ の光子})
\label{eq:inclusive}
\end{equation}
になる（\ref{app:ir}）。有限部分の $\pm\pi^2$ は $N_\epsilon$ の $\epsilon^2$ の項の取り方で変わるが、二つで符号が逆なので和にはよらない。ベクトルと軸性ベクトルの結合は質量0では同じ補正を受ける。NLO の図を図~\ref{fig:ir} に示す。$W^+\to\ell^+\nu$ では光子が $W$ にも付くので、光子の部分だけを取り出した補正は単独ではゲージによらず、弱い補正と合わせて初めて意味を持つ。

\begin{figure}[tbp]
\centering
\begin{tikzpicture}[scale=0.85]
% ---- Z -> f fbar ----
% (a) emission from f
\begin{scope}[shift={(0,0)}]
  \draw[photon] (-1.1,0) -- (0,0);
  \draw[fermion] (0,0) -- (0.5,0.35); \draw[fermion] (0.5,0.35) -- (1.2,0.55); \draw[fermion] (1.15,-0.75) -- (0,0);
  \draw[gluon] (0.5,0.35) -- (0.95,1.1);
  \node[vtx] at (0,0) {}; \node[vtx] at (0.5,0.35) {};
  \node[font=\scriptsize] at (-0.9,0.25) {$Z$}; \node[font=\scriptsize] at (1.4,0.55) {$f$}; \node[font=\scriptsize] at (1.35,-0.8) {$\bar f$};
  \node[font=\scriptsize] at (1.2,1.15) {$g$};
  \node[slab] at (0.1,-1.6) {(a) 実の放射};
\end{scope}
% (b) emission from fbar
\begin{scope}[shift={(3.4,0)}]
  \draw[photon] (-1.1,0) -- (0,0);
  \draw[fermion] (0,0) -- (1.15,0.75); \draw[fermion] (1.2,-0.55) -- (0.5,-0.35); \draw[fermion] (0.5,-0.35) -- (0,0);
  \draw[gluon] (0.5,-0.35) -- (0.95,-1.1);
  \node[vtx] at (0,0) {}; \node[vtx] at (0.5,-0.35) {};
  \node[font=\scriptsize] at (-0.9,0.25) {$Z$}; \node[font=\scriptsize] at (1.35,0.8) {$f$}; \node[font=\scriptsize] at (1.4,-0.55) {$\bar f$};
  \node[font=\scriptsize] at (1.2,-1.15) {$g$};
  \node[slab] at (0.1,-1.6) {(b) 実の放射};
\end{scope}
% (c) vertex
\begin{scope}[shift={(6.8,0)}]
  \draw[photon] (-1.1,0) -- (0,0);
  \draw[fermion] (0,0) -- (0.6,0.42); \draw[fermion] (0.6,0.42) -- (1.15,0.8);
  \draw[fermion] (1.15,-0.8) -- (0.6,-0.42); \draw[fermion] (0.6,-0.42) -- (0,0);
  \draw[gluon] (0.6,0.42) -- (0.6,-0.42);
  \node[vtx] at (0,0) {}; \node[vtx] at (0.6,0.42) {}; \node[vtx] at (0.6,-0.42) {};
  \node[font=\scriptsize] at (-0.9,0.25) {$Z$}; \node[font=\scriptsize] at (1.35,0.85) {$f$}; \node[font=\scriptsize] at (1.35,-0.85) {$\bar f$};
  \node[font=\scriptsize] at (0.95,0) {$g$};
  \node[slab] at (0.1,-1.6) {(c) 頂点の補正};
\end{scope}
% (d) external field constant
\begin{scope}[shift={(10.2,0)}]
  \draw[photon] (-1.1,0) -- (0,0);
  \draw[fermion] (0,0) -- (0.35,0.245); \draw[fermion] (0.35,0.245) -- (0.9,0.63); \draw[fermion] (0.9,0.63) -- (1.2,0.84);
  \draw[fermion] (1.2,-0.84) -- (0,0);
  \draw[gluon] (0.35,0.245) .. controls (0.3,0.85) and (0.75,1.15) .. (0.9,0.63);
  \node[vtx] at (0,0) {}; \node[vtx] at (0.35,0.245) {}; \node[vtx] at (0.9,0.63) {};
  \node[font=\scriptsize] at (-0.9,0.25) {$Z$}; \node[font=\scriptsize] at (1.4,0.9) {$f$}; \node[font=\scriptsize] at (1.4,-0.9) {$\bar f$};
  \node[slab] at (0.1,-1.6) {(d) 場の定数 $\delta Z$};
\end{scope}
% ---- t -> b W+ ----
% (e) emission from t
\begin{scope}[shift={(0,-3.4)}]
  \draw[fermion] (-1.15,0) -- (-0.55,0); \draw[fermion] (-0.55,0) -- (0,0);
  \draw[fermion] (0,0) -- (1.15,0.75); \draw[photon] (0,0) -- (1.15,-0.75);
  \draw[gluon] (-0.55,0) -- (-0.1,0.95);
  \node[vtx] at (0,0) {}; \node[vtx] at (-0.55,0) {};
  \node[font=\scriptsize] at (-1.0,0.22) {$t$}; \node[font=\scriptsize] at (1.35,0.8) {$b$}; \node[font=\scriptsize] at (1.4,-0.8) {$W^+$};
  \node[font=\scriptsize] at (0.15,1.05) {$g$};
  \node[slab] at (0.1,-1.3) {(e) 実の放射};
\end{scope}
% (f) emission from b
\begin{scope}[shift={(3.4,-3.4)}]
  \draw[fermion] (-1.15,0) -- (0,0);
  \draw[fermion] (0,0) -- (0.5,0.35); \draw[fermion] (0.5,0.35) -- (1.2,0.55); \draw[photon] (0,0) -- (1.15,-0.75);
  \draw[gluon] (0.5,0.35) -- (0.95,1.1);
  \node[vtx] at (0,0) {}; \node[vtx] at (0.5,0.35) {};
  \node[font=\scriptsize] at (-1.0,0.22) {$t$}; \node[font=\scriptsize] at (1.4,0.55) {$b$}; \node[font=\scriptsize] at (1.4,-0.8) {$W^+$};
  \node[font=\scriptsize] at (1.2,1.15) {$g$};
  \node[slab] at (0.1,-1.3) {(f) 実の放射};
\end{scope}
% (g) vertex
\begin{scope}[shift={(6.8,-3.4)}]
  \draw[fermion] (-1.15,0) -- (-0.55,0); \draw[fermion] (-0.55,0) -- (0,0);
  \draw[fermion] (0,0) -- (0.6,0.42); \draw[fermion] (0.6,0.42) -- (1.15,0.8); \draw[photon] (0,0) -- (1.15,-0.75);
  \draw[gluon] (-0.55,0) .. controls (-0.5,0.75) and (0.2,0.95) .. (0.6,0.42);
  \node[vtx] at (0,0) {}; \node[vtx] at (-0.55,0) {}; \node[vtx] at (0.6,0.42) {};
  \node[font=\scriptsize] at (-1.0,0.22) {$t$}; \node[font=\scriptsize] at (1.35,0.85) {$b$}; \node[font=\scriptsize] at (1.4,-0.8) {$W^+$};
  \node[font=\scriptsize] at (-0.2,0.98) {$g$};
  \node[slab] at (0.1,-1.3) {(g) 頂点の補正};
\end{scope}
% (h) field constant of t
\begin{scope}[shift={(10.2,-3.4)}]
  \draw[fermion] (-1.55,0) -- (-1.1,0); \draw[fermion] (-1.1,0) -- (-0.35,0); \draw[fermion] (-0.35,0) -- (0,0);
  \draw[fermion] (0,0) -- (1.15,0.75); \draw[photon] (0,0) -- (1.15,-0.75);
  \draw[gluon] (-1.1,0) arc (180:0:0.375);
  \node[vtx] at (0,0) {}; \node[vtx] at (-1.1,0) {}; \node[vtx] at (-0.35,0) {};
  \node[font=\scriptsize] at (-1.5,-0.22) {$t$}; \node[font=\scriptsize] at (1.35,0.8) {$b$}; \node[font=\scriptsize] at (1.4,-0.8) {$W^+$};
  \node[slab] at (0.0,-1.3) {(h) $\delta Z_t$};
\end{scope}
% (i) field constant of b
\begin{scope}[shift={(13.6,-3.4)}]
  \draw[fermion] (-1.15,0) -- (0,0);
  \draw[fermion] (0,0) -- (0.35,0.245); \draw[fermion] (0.35,0.245) -- (0.9,0.63); \draw[fermion] (0.9,0.63) -- (1.2,0.84);
  \draw[photon] (0,0) -- (1.15,-0.75);
  \draw[gluon] (0.35,0.245) .. controls (0.3,0.85) and (0.75,1.15) .. (0.9,0.63);
  \node[vtx] at (0,0) {}; \node[vtx] at (0.35,0.245) {}; \node[vtx] at (0.9,0.63) {};
  \node[font=\scriptsize] at (-1.0,0.22) {$t$}; \node[font=\scriptsize] at (1.4,0.9) {$b$}; \node[font=\scriptsize] at (1.4,-0.8) {$W^+$};
  \node[slab] at (0.1,-1.3) {(i) $\delta Z_b$};
\end{scope}
\end{tikzpicture}
\caption{手順 (2)、(3) の NLO の図。上の段は $Z\to f\bar f$、下の段は $t\to bW^+$ で、巻いた線はグルーオンである（$Z\to\ell^+\ell^-$ と $Z\to\tau^+\tau^-$ では光子で、$C_F\to Q_f^2$）。(a)(b)(e)(f) は実の放射、(c)(g) は頂点の補正、(d)(h)(i) は外線の自己エネルギー（on-shell の場の定数。(d) は $\bar f$ の脚にも同じものがある）。質量0のフェルミオンでは (d) はスケールのない積分で0になる。仮想の補正の赤外の極は、実の放射の軟らかい領域（式~\eqref{eq:eikonal}）の極と打ち消す。手順 (2) の独立な計算では、放射する粒子に質量 $\lambda$ を与える。}
\label{fig:ir}
\end{figure}

\subsection{$\rho$ パラメータ}\label{sec:rho}

ゲージ場の自己エネルギーの横の部分を $\Sigma_T$ と書き、$\Delta\rho\equiv\hat\Sigma^{ZZ}_T(0)/M_Z^2-\hat\Sigma^{WW}_T(0)/M_W^2$ とする。トップとボトムの二重項の寄与は、$m_b\to0$ で
\begin{equation}
\Delta\rho=\frac{N_cy_t^2}{2K}=\frac{3y_t^2}{32\pi^2}\qquad\Bigl(=\frac{3G_Fm_t^2}{8\sqrt2\pi^2}\Bigr)
\label{eq:drho}
\end{equation}
である~\cite{veltman77}。紫外の極は custodial 対称性のために $Z$ と $W$ の間で打ち消す。FJ の処方のタドポールは $\Sigma^{ZZ}_T(0)$、$\Sigma^{WW}_T(0)$ にそれぞれ $M_Z^2$、$M_W^2$ に比例する同じ係数で入るので、$\Delta\rho$ からは消える。$G_F$ と $\MSbar$ の $v$ の関係（SM ノートの手順 (6) の $\mu$ 崩壊）と $\Delta\rho$ から、$M_W$ の予言（手順 (6)）を組む。

\subsection{$Z$ 極の有効結合}\label{sec:zpole}

$Z\to f\bar f$ の繰り込んだ1ループの振幅を、$p^2=M_Z^2$ で
\begin{equation}
\mathcal M=-i\frac{g_Z}2\sqrt{\hat Z_Z}\,\bar u(p_1)\gamma^\mu\bigl(\mathcal V_f-\mathcal A_f\gamma_5\bigr)v(p_2)\,\varepsilon_\mu
\label{eq:effcoup}
\end{equation}
と書き（$\hat Z_Z$ は $Z$ の留数、$\mathcal V_f$、$\mathcal A_f$ は複素の有効結合）、
\begin{equation}
\sin^2\theta^f_{\mathrm{eff}}\equiv\frac1{4|Q_f|}\Bigl(1-\re\frac{\mathcal V_f}{\mathcal A_f}\Bigr)=\hat\kappa_f(\mu)\,\hat s^2(\mu)
\label{eq:s2eff}
\end{equation}
と定める。木では $\hat\kappa_f=1$ で、1ループの $\hat\kappa_f-1$ は、$Z$ と光子の混合（$\re\hat\Sigma^{AZ}_T(M_Z^2)/M_Z^2$。光子の頂点 $-ieQ_f\gamma^\mu$ を通じて $\mathcal V_f$ だけをずらす）と、頂点と外線の補正からなる。部分幅は $|\mathcal V_f|^2$、$|\mathcal A_f|^2$ に式~\eqref{eq:lo} の形で入り、QED と QCD の補正 \eqref{eq:inclusive} がほぼ掛け算で付く。$R_\ell=\Gamma_{\mathrm{had}}/\Gamma_\ell$、$R_b=\Gamma_{b\bar b}/\Gamma_{\mathrm{had}}$、$\sigma^0_{\mathrm{had}}=12\pi\Gamma_e\Gamma_{\mathrm{had}}/(M_Z^2\Gamma_Z^2)$ はこれらの比と積である。$Z b\bar b$ の頂点には $W$ と $G^\pm$ と内線のトップのループが入り、$|V_{tb}|^2$ と $y_t^2$ に比例する非普遍的な補正になる。符号と細部は手順 (5) で SM ノートの規約に合わせて求める。

\section{検算の計画}\label{sec:plan}

第~\ref{sec:results} 節の各量を、どの手段で計算して突き合わせるかを表~\ref{tab:plan} に示す。手計算（付録）を基準とし、ツールの出力がそれと一致するかを確かめる。ツールの版と配置は計算環境の記録~\cite{env} にあり、SM ノート、CKM ノートと同じである。仮想の補正は本ノートの模型ファイル \texttt{SM\_CKM.fr} から書き出した FeynArts の模型（\texttt{GaugeXi} を残す）と FeynCalc で求め、実の放射は木の3体の振幅を FeynCalc で記号として求めて、位相空間を解析的に（質量0）または数値で（質量あり）積分する。MadGraph では、同じ模型ファイルから作った CKM ノートの UFO（電弱と QCD の対抗項を点の数値で持つもの）で MadLoop の仮想の補正を、木の3体の過程で実の放射を求める。測定値と文献の値は参照として並べ、比べる量は本ノートの模型で求めたものにする。

\begin{table}[tbp]
\centering\small
\caption{各量と、それを計算して突き合わせる手段。「—」はその手段を使わないことを表す。}
\label{tab:plan}
\begin{tabularx}{\linewidth}{>{\raggedright\arraybackslash}p{3.6cm}>{\raggedright\arraybackslash}p{1.3cm}>{\raggedright\arraybackslash}X>{\raggedright\arraybackslash}p{2.0cm}>{\raggedright\arraybackslash}p{2.2cm}}
\toprule
\textbf{量} & \textbf{手計算} & \textbf{FeynCalc + FeynArts、FormCalc} & \textbf{数値} & \textbf{MadGraph} \\
\midrule
木の幅 \eqref{eq:lo} & \ref{app:lo} & 振幅の2乗を記号として & 入力の走り & 木の幅 \\
赤外の極の打ち消し \eqref{eq:softpole}--\eqref{eq:inclusive} & \ref{app:ir} & $d$ 次元の3体の振幅、仮想の補正 & 位相空間の積分 & MadLoop と木の3体 \\
NLO の幅（$Z$、$W$、$t$） & 手順 (4) & 一般の $R_\xi$。FormCalc & $\mu$ の依存 & MadLoop と木の3体 \\
$Z$ 極の量 \eqref{eq:s2eff} & 第~\ref{sec:zpole} 節 & 一般の $R_\xi$ & 測定値との比較 & — \\
$\Delta\rho$ \eqref{eq:drho}、$M_W$ & 第~\ref{sec:rho} 節 & 自己エネルギー（SM ノート） & 測定値との比較 & — \\
\bottomrule
\end{tabularx}
\end{table}

\paragraph{一致の基準。}SM ノートと同じである。発散部分は $1/\epsilon$ の係数が記号として一致すること（仮想と実の放射の極の打ち消しは記号として0になること）を、有限部分は倍精度の丸め誤差の範囲（相対 $10^{-12}$ 程度）で一致することを求める。数値の位相空間の積分では、積分の精度を明示してその範囲で一致を求める。$\xi$ によらないはずの量は、記号として $\xi$ を含まないこと、または $\xi=0,1,3$ で同じ値になることを求める。測定値とは、本ノートで計算しない高次の補正の大きさ（$\mu$ を変えた変化と、文献の高次の値との差）の範囲で比べる。

\paragraph{手順。}次の順に進める。(1) 木の幅 \eqref{eq:lo} を、SM ノートの $\MSbar$ のパラメータを走らせた値と極質量で求め、FeynCalc の振幅の2乗と MadGraph の木の幅と比べる。(2) 質量0のフェルミオン対への崩壊（$Z\to\mu^+\mu^-$ の光子、$Z\to u\bar u$ のグルーオン）で、$d$ 次元の3体の振幅を FeynCalc で求めて位相空間を解析的に積分し、式~\eqref{eq:masslessVR} の $\Gamma_R$ と、仮想の補正との極の打ち消しと、式~\eqref{eq:inclusive} を確かめる。光子の質量で正則化した独立な計算とも比べる。(3) 質量のあるフェルミオン（$Z\to b\bar b$、$Z\to\tau^+\tau^-$、$t\to bW^+$）で、軟らかい部分を式~\eqref{eq:eikonal} の解析的な積分に、硬い部分を数値の積分にして（位相空間の切断 $\delta$ によらないことを確かめる）、仮想の補正の極 \eqref{eq:softpole} との打ち消しと、包括的な QCD の補正を求める。$t\to bW^+$ は文献の $O(\alpha_s)$ の結果~\cite{jk89}と参照として比べる。(4) $Z\to f\bar f$、$W^+\to f\bar f'$、$t\to bW^+$ の電弱と QCD の NLO の幅を、一般の $R_\xi$ ゲージの仮想の補正と実の放射で求め、$\xi$ によらないこと、$\mu$ への依存が走りと打ち消すことを確かめる。MadLoop の仮想の補正と木の3体の過程で求めた値と比べる。(5) $Z$ の有効結合 \eqref{eq:effcoup}、$\sin^2\theta^\ell_{\mathrm{eff}}$、$\Gamma_Z$、$R_\ell$、$R_b$、$\sigma^0_{\mathrm{had}}$ を求め、LEP と SLD の測定値~\cite{lepewwg06,pdg24}と比べる。(6) $\Delta\rho$ \eqref{eq:drho} を自己エネルギーから求め、$G_F$、$M_Z$、$\alpha$ からの $M_W$ の予言を SM ノートの手順 (7) と測定値と比べる。(7) $\MSbar$ と on-shell の処方の結果の差（SM ノートの手順 (5)〜(7) の換算で結ぶ）が高次の大きさであることを確かめる。(8) 予言をまとめ、$\mu$ の依存と文献の高次の値との差から不定性を見積もって、測定値と比べる。

\section{検算}\label{sec:checks}

第~\ref{sec:plan} 節の手順を順に行う。スクリプトは \texttt{Src/} にあり、手順 (1) は \texttt{run\_lo\_check.sh}（約1分）で、手順 (2) はその後に \texttt{run\_ir\_check.sh}（約4分）で、手順 (3) はさらにその後に \texttt{run\_massive\_check.sh}（約6分。ほとんどは MadLoop の模型の読み込み）で、手順 (4a) はその後に \texttt{run\_zweak\_check.sh}（約15分）で、手順 (4b) はさらにその後に \texttt{run\_wweak\_check.sh}（約20分）で、手順 (4c) はその後に \texttt{run\_tweak\_check.sh}（約15分）で、手順 (5) は (4a) の後に \texttt{run\_zpole\_check.sh}（約3分）で、手順 (6) は (4b) の後に \texttt{run\_mw\_check.sh}（約2分）で、手順 (7) は (4c) の後に \texttt{run\_scheme\_check.sh}（約30秒）で、手順 (8) は最後に \texttt{run\_summary\_check.sh}（約10秒）で、まとめて再実行できる。

\paragraph{手順 (1)：木の幅。}(a) 模型。\texttt{generate\_SM\_CKM.wl} で、写した \texttt{SM\_CKM.fr} から FeynRules の頂点と FeynArts の模型と UFO を書き出した（CKM ノートと同じく UFO を先に書く）。93個の頂点のクラスは CKM ノートの FeynRules の出力と同一で（Mathematica の \texttt{===}）、FeynArts の模型の3個のファイルと UFO の8個のファイルも、日付と場所の行を除いて CKM ノートのものと同一だった。

(b) 振幅の2乗。\texttt{lo\_feyncalc.wl} は、本ノートの FeynArts の模型で2体の崩壊の振幅の2乗（スピン、偏極、色の和）を FeynCalc で記号として求め、\ref{app:lo} の手計算の式と比べる。$Z\to\nu\bar\nu$、$\tau^-\tau^+$、$c\bar c$、$b\bar b$、$H\to b\bar b$、$\tau^-\tau^+$、$c\bar c$、$W^+\to u_i\bar d_j$（6個。$V$ は一般の複素行列）、$W^+\to\nu_\ell\ell^+$（3個）、$t\to d_jW^+$（3個）の19個で、差は記号として0だった（質量はすべて記号のまま）。

(c) 入力。\texttt{nlo\_inputs.py} は、SM ノートの手順 (7) の $M_t$ の値と CKM ノートの手順 (6) の $V(M_t)$ を、CKM ノートの1ループの方程式で走らせる（表~\ref{tab:lopar}）。これは全標準模型の $\MSbar$ の値で（トップを分離しない）、$\mu=M_Z$ の $\hat s^2=0.23242$ は、PDG の $\hat s^2(M_Z)=0.23122$（トップを分離した5フレーバーの処方）とは処方の違う量なので比べない。

(d) 幅。\texttt{lo\_widths.py} は式~\eqref{eq:lo} を $\mu=M$ の結合と極質量で求める（表~\ref{tab:lo}）。MadGraph の UFO は結合を $\alpha^{-1}$、$M_Z$、$M_W$ の木の関係から作るので、同じ結合と運動学を持つ param\_card を四つ作った（$Z$ では $M_Z$ を極質量にして $M_W=M_Z\hat c(M_Z)$、$\alpha^{-1}=\hat\alpha^{-1}(M_Z)$。$W$ と $t$ では $M_W$ を極質量にして $M_Z=M_W/\hat c(\mu)$。$H$ では UFO の湯川結合が運動学の質量で決まるので、式の確かめだけ）。MadGraph の \texttt{compute\_widths} の27個の部分幅は、式~\eqref{eq:lo} と相対 $2\times10^{-15}$ 以内で一致した。MadGraph は「QCD のスケール」より小さい部分幅（$t\to dW^+$、$t\to sW^+$）を0にして書き出すが、その値をログに出すので、それで比べた。

(e) 物理。木の $\Gamma_Z=2.4279$~GeV は測定値~\cite{pdg24}より $2.7\%$ 小さい。差のほとんどはハドロンの幅（1.6762~GeV と 1.7444~GeV、$-3.9\%$）で、式~\eqref{eq:inclusive} の QCD の補正 $\alpha_s/\pi=3.7\%$ にあたる。レプトンの幅（$83.81$~MeV と $83.98$~MeV）と見えない幅（$500.4$~MeV と $499.0$~MeV）は $0.2$〜$0.3\%$ で合う。$\Gamma_W=2.0350$~GeV（測定値 $2.085\pm0.042$~GeV）、$\Gamma_t=1.4536$~GeV（$1.42^{+0.19}_{-0.15}$~GeV）である。$\mu$ を $M/2$、$2M$ にした変化は $\pm1$〜$1.4\%$ で、$\MSbar$ の結合の走りから来る（$H\to b\bar b$ は $y_b$ の QCD の走りで $+11\%$、$-9\%$）。この依存は手順 (2)〜(4) の NLO の補正で打ち消されるはずである。

\begin{table}[tbp]
\centering\small
\caption{手順 (1) の入力：SM ノートの手順 (7) の $\MSbar$ の値（$\mu=M_t$）を CKM ノートの1ループの方程式で走らせた値（\texttt{nlo\_inputs.py}）。$\hat s^2=g'^2/(g^2+g'^2)$、$\hat\alpha=\hat e^2/(4\pi)$。}
\label{tab:lopar}
\begin{tabular}{lllllllll}
\toprule
$\mu$ [GeV] & $g'$ & $g$ & $\hat\alpha^{-1}$ & $\hat s^2$ & $\alpha_s$ & $y_t$ & $y_b$ & $|V_{cb}|$ \\
\midrule
$M_W=80.3692$ & 0.357831 & 0.651440 & 127.7532 & 0.231787 & 0.11902 & 0.97192 & 0.016329 & 0.041669 \\
$M_Z=91.1876$ & 0.358082 & 0.650741 & 127.6795 & 0.232420 & 0.11706 & 0.96512 & 0.016180 & 0.041715 \\
$M_H=125.20$ & 0.358713 & 0.648996 & 127.4945 & 0.234010 & 0.11241 & 0.94871 & 0.015822 & 0.041831 \\
$M_t=172.57$ & 0.359356 & 0.647244 & 127.3073 & 0.235625 & 0.10807 & 0.93298 & 0.015480 & 0.041943 \\
\bottomrule
\end{tabular}
\end{table}

\begin{table}[tbp]
\centering\small
\caption{手順 (1)：木の幅（\texttt{lo\_widths.py}、$\MSbar$ の結合と極質量。GeV）。$\mu$ の変化は $\mu=M/2$、$2M$ の値の $\mu=M$ からの差。測定値は PDG 2024~\cite{pdg24}（参照。NLO の補正を含む量なので、木の値とは手順 (2)〜(4) の補正の分だけ違う）。}
\label{tab:lo}
\begin{tabular}{llllll}
\toprule
量 & 木（$\mu=M$） & $\mu=M/2$ & $\mu=2M$ & 測定値（参照） \\
\midrule
$\Gamma(Z\to e^+e^-)$ & 0.083814 & $+0.95\%$ & $-0.89\%$ & $0.083984\pm0.000086$（$\ell^+\ell^-$） \\
$\Gamma(Z\to\tau^+\tau^-)$ & 0.083625 & $+0.95\%$ & $-0.89\%$ & \\
$\Gamma(Z\to\nu\bar\nu)\times3$ & 0.500412 & $+0.74\%$ & $-0.71\%$ & $0.4990\pm0.0015$ \\
$\Gamma(Z\to u\bar u)$、$c\bar c$ & 0.286377、0.285873 & $+1.36\%$ & $-1.32\%$ & \\
$\Gamma(Z\to d\bar d)$、$b\bar b$ & 0.369366、0.365247 & $+1.17\%$ & $-1.14\%$ & \\
$\Gamma(Z\to\mathrm{hadrons})$ & 1.676229 & $+1.24\%$ & $-1.20\%$ & $1.7444\pm0.0020$ \\
$\Gamma_Z$ & 2.427894 & $+1.11\%$ & $-1.07\%$ & $2.4955\pm0.0023$ \\
\midrule
$\Gamma(W^+\to e^+\nu)$ & 0.226176 & $+1.19\%$ & $-1.17\%$ & \\
$\Gamma(W^+\to u\bar d)$、$c\bar s$ & 0.644166、0.642621 & $+1.19\%$ & $-1.17\%$ & \\
$\Gamma_W$ & 2.034973 & $+1.19\%$ & $-1.17\%$ & $2.085\pm0.042$ \\
\midrule
$\Gamma(t\to bW^+)$ & 1.451064 & $+1.18\%$ & $-1.15\%$ & \\
$\Gamma_t$（$d_jW^+$ の和） & 1.453649 & $+1.18\%$ & $-1.15\%$ & $1.42^{+0.19}_{-0.15}$ \\
\midrule
$\Gamma(H\to b\bar b)$ & 0.001854 & $+10.55\%$ & $-8.83\%$ & \\
$\Gamma(H\to\tau^+\tau^-)$ & 0.000247 & $-1.20\%$ & $+1.04\%$ & \\
\bottomrule
\end{tabular}
\end{table}

\paragraph{手順 (2)：質量0のフェルミオン対の赤外の極。}(a) 実の放射の振幅。\texttt{ir\_massless\_feyncalc.wl} は、$Z\to f\bar f\gamma$ の2個の図の振幅の2乗を FeynCalc で $d$ 次元のまま求める。流れ $\gamma^\mu$ と $\gamma^\mu\gamma_5$ は別々に扱い（$\gamma_5$ は反可換）、二つの偏極の和は $-g$ とする（流れが保存するので）。比 $r$ は \ref{app:ir} の式と記号として一致し、軸性ベクトルでも同じだった。$d=4$ では、本ノートの FeynArts の模型の $Z\to\mu^-\mu^+\gamma$ と $Z\to u\bar ug$ の振幅の2乗が、木の結合と $r$ の積と記号として一致した（グルーオンは色の因子 $C_F=4/3$）。\texttt{ir\_mg5.py} では、MadGraph の木の $|\mathcal M|^2$ を同じ式と比べた（UFO、手順 (1) の $Z$ の param\_card で $m_\mu=0$）。点は Dalitz の平面の7点で、中心から軟らかい端と共線の端の近く（$y=10^{-4}$）までとった。14個の値は相対 $6\times10^{-14}$ 以内で一致した（端に近い点ほど、運動量の丸めで桁を落とす）。MadGraph の Fortran の param\_card の読み込みは値を20文字で切るので、card の値は14桁で書いた（17桁で書くと値が切れて読めない）。

(b) 実の放射の積分。$d$ 次元の3体の位相空間を $y_{ij}$ の Dirichlet の積分にして（\ref{app:ir}）、$\Gamma_R/\Gamma_0$ を $\epsilon$ について厳密に求めた（式~\eqref{eq:irexact}）。展開は $N_\epsilon(2/\epsilon^2+3/\epsilon+19/2-\pi^2)$ だった。

(c) 仮想の補正。頂点の補正を TID で $B_0$ と $C_0$ に帰着させた。質量0の $C_0$ を $B_0$ で書くと、$\mathcal M_1/\mathcal M_0$ は $B_0$ と $\epsilon$ の有理式の積になった（ベクトルと軸性ベクトルで同じ）。$B_0$ の位相の実部をとると式~\eqref{eq:irexact} の $\Gamma_V$ になり、展開は $N_\epsilon(-2/\epsilon^2-3/\epsilon-8+\pi^2)$ だった。

(d) 和。$\Gamma_V+\Gamma_R$ の $1/\epsilon^2$ と $1/\epsilon$ は記号として消えた。残りは $\Gamma_0N_\epsilon\cdot\frac32+O(\epsilon)$ で、式~\eqref{eq:inclusive} の $\frac34\alpha C/\pi$ である（グルーオンでは $C=C_F$ なので $\alpha_s/\pi$）。ここで手計算の誤りが一つ見つかった。計画を立てたときの式~\eqref{eq:masslessVR} は、共通の因子を $(4\pi\mu^2/M^2)^\epsilon H(\epsilon)$（$H=1+O(\epsilon^2)$）とし、有限部分を $-8$ と $19/2$ としていた。しかし有限部分は因子の $\epsilon^2$ の項で符号を逆にしてずれるので、$H$ を定めない有限部分には意味がない。さらに $1/\Gamma(1-\epsilon)$ には $O(\epsilon)$ の項がある。そこで式~\eqref{eq:masslessVR} を $N_\epsilon$ の形に直した。有限部分は $-8+\pi^2$ と $19/2-\pi^2$ で、文献~\cite{esw96}の値と同じである。和は因子の取り方によらないので、式~\eqref{eq:inclusive} は変わらない。

(e) 光子の質量。\texttt{ir\_photonmass.wl} は独立な正則化として、放射する粒子に質量 $\lambda$ を与え（$w=\lambda^2/M^2$。フェルミオンは質量0のまま）、$d=4$ で計算する。軟らかい領域も共線の領域も $\lambda$ で切れるので、$\Gamma_V$ と $\Gamma_R$ は別々に有限になる。仮想の補正は、頂点と場の定数 $\delta Z=\Sigma_V(0)$（$\Sigma(p)=\slashed p\,\Sigma_V(p^2)$）を、TID と Package-X（FeynHelpers）で解析的に求めた。紫外の極は二つの間で記号として消え（Ward の恒等式）、$\mu$ も消えて、
\begin{equation*}
\frac{\Gamma_V}{\Gamma_0}=\frac{\alpha C}{2\pi}\Bigl[-\frac72-2w-(3+2w)\ln w-(1+w)^2\Bigl(\ln^2w+2\re\mathrm{Li}_2(1+w)-\frac{2\pi^2}3\Bigr)\Bigr]
\end{equation*}
になった。実の放射は、質量 $\lambda$ の放射の $|\mathcal M_R|^2$ を Dalitz の平面で数値で積分した（NIntegrate、30桁、精度の目標 $10^{-13}$。$w=0$ で $r$ に戻ることを記号で確かめた）。結果は表~\ref{tab:photonmass} の通りである。二つはそれぞれ $\ln^2w$ で発散するが、和は $w\to0$ で $\frac{\alpha C}{2\pi}\cdot\frac32$ に近づいた。差は $w=10^{-4}$ で $-1.5\times10^{-8}$、$10^{-6}$ で $-1.5\times10^{-12}$（$\approx-1.5w^2$）で、$w\le10^{-8}$ では積分の精度の $10^{-14}$ 程度だった。$\Gamma_V$ は上の式の $w\to0$ の形 $\frac{\alpha C}{2\pi}(-\ln^2w-3\ln w-\frac72+\frac{\pi^2}3)$（$\mathrm{Li}_2(1)=\pi^2/6$）に、$\Gamma_R$ はそれと和が $\frac32$ になる形に、$2w\ln^2w$ 程度の差で近づく。

(f) 物理。\texttt{ir\_inclusive.py} で、式~\eqref{eq:inclusive} を手順 (1) の木の幅に掛けた（$\hat\alpha(M)$、$\alpha_s(M)$。クォークの質量は木の幅にだけ入れる）。結果は表~\ref{tab:inclusive} の通りである。ハドロンの幅は $\alpha_s/\pi=3.73\%$ で $1.7394$~GeV になり、測定値との差は $-3.9\%$ から $-0.29\%$ に縮んだ。$\Gamma_Z$ の差は $-2.7\%$ から $-0.16\%$ に、$R_\ell$ は $-3.7\%$ から $-0.25\%$ に縮んだ。レプトンの幅は QED の $+0.19\%$ で、差は $-0.02\%$ になった。$\Gamma_W$ には QCD の因子だけを掛けた（光子は $W$ にも付くので、QED の部分だけではゲージによらない）。$2.0864$~GeV で、測定値 $2.085\pm0.042$~GeV と合う。残りの差は、$(\alpha_s/\pi)^2=0.14\%$ の大きさの高次の QCD の補正、手順 (3) の質量の補正、手順 (4)、(5) の電弱の補正の大きさである。見えない幅の $+0.28\%$ は QED と QCD の補正を受けないので、電弱の補正だけで決まる。式~\eqref{eq:inclusive} は $\mu$ によらないので、木の幅の $\mu$ の依存は手順 (4) の電弱の補正で打ち消されるはずである。

\begin{table}[tbp]
\centering\small
\caption{手順 (2)：光子の質量 $\lambda$ で正則化した仮想の補正と実の放射（\texttt{ir\_photonmass.wl}。$w=\lambda^2/M^2$、$\Gamma_0\frac{\alpha C}{2\pi}$ を単位とする）。和は $w\to0$ で $\frac32$ になるはずである。}
\label{tab:photonmass}
\begin{tabular}{rrrr}
\toprule
$w$ & $\Gamma_V$ & $\Gamma_R$ & $\Gamma_V+\Gamma_R-\frac32$ \\
\midrule
$10^{-4}$ & $-57.426189745079$ & $58.926189730086$ & $-1.5\times10^{-8}$ \\
$10^{-6}$ & $-149.632311326793$ & $151.132311326792$ & $-1.5\times10^{-12}$ \\
$10^{-8}$ & $-284.269575465696$ & $285.769575465696$ & $5\times10^{-15}$ \\
$10^{-10}$ & $-461.322390230102$ & $462.822390230102$ & $1\times10^{-14}$ \\
\bottomrule
\end{tabular}
\end{table}

\begin{table}[tbp]
\centering\small
\caption{手順 (2)：式~\eqref{eq:inclusive} を木の幅（表~\ref{tab:lo}）に掛けた値（\texttt{ir\_inclusive.py}。GeV）。測定値は PDG 2024~\cite{pdg24}（参照）。括弧は測定値との差。}
\label{tab:inclusive}
\begin{tabular}{llll}
\toprule
量 & 木 & 式~\eqref{eq:inclusive} を掛けた値 & 測定値（参照） \\
\midrule
$\Gamma(Z\to\ell^+\ell^-)$ & 0.083814（$-0.20\%$） & 0.083971（$-0.02\%$） & $0.083984\pm0.000086$ \\
$\Gamma(Z\to\nu\bar\nu)\times3$ & 0.500412（$+0.28\%$） & 0.500412（$+0.28\%$） & $0.4990\pm0.0015$ \\
$\Gamma(Z\to\mathrm{hadrons})$ & 1.676229（$-3.91\%$） & 1.739392（$-0.29\%$） & $1.7444\pm0.0020$ \\
$\Gamma_Z$ & 2.427894（$-2.71\%$） & 2.491527（$-0.16\%$） & $2.4955\pm0.0023$ \\
$R_\ell$ & 19.9993（$-3.70\%$） & 20.7142（$-0.25\%$） & $20.767\pm0.025$ \\
$\Gamma_W$（QCD だけ） & 2.034973（$-2.40\%$） & 2.086368（$+0.07\%$） & $2.085\pm0.042$ \\
\bottomrule
\end{tabular}
\end{table}

\FloatBarrier
\paragraph{手順 (3)：質量のある放射体。}補正を $\Gamma^{(1)}/\Gamma_0=\frac{\alpha C}\pi K$ と書く。対象は $Z\to Q\bar Q$（流れ $\gamma^\mu$ と $\gamma^\mu\gamma_5$ を別々に。$K_V$、$K_A$）と $t\to bW^+$（$\gamma^\mu P_L$）である。

(a) 仮想の補正。\texttt{ir\_massive\_virtual.wl} は、頂点の補正を $d$ 次元で TID により PaVe 関数に帰着させ、Package-X（FeynHelpers）で解析的に求める。質量のある外線の on-shell の場の定数 $\delta Z=\Sigma_V(m^2)+2m^2(\Sigma_V'+\Sigma_S')(m^2)$ も同じ方法で求めた（$p^2$ の微分は被積分関数でとる）。結果は $-\frac{\alpha C}{4\pi}[3(\frac1\epsilon-\gamma_E+\ln\frac{4\pi\mu^2}{m^2})+4]$ で、CKM ノートが NLOCT と比べた定数と同じ形である。紫外の極は頂点と場の定数の間で消えた。残る赤外の単極は、$Z$ では式~\eqref{eq:softpole} と記号として一致した（実部。虚部は Coulomb の位相）。$t\to bW^+$ では、$p_1\cdot p_2=p_t\cdot p_b$ とした式~\eqref{eq:softpole} と、三つの質量の組で一致した。

(b) 軟らかい放射。崩壊する粒子の静止系で放射のエネルギーが $E<\Delta E=\delta E_{\max}$ の領域では、実の放射は木の幅と式~\eqref{eq:eikonal} の積になる。これを $d$ 次元で積分すると
\begin{equation*}
K_S=-\frac1{4\epsilon}\Bigl(\frac{4\pi\mu^2}{\Delta E^2}\Bigr)^\epsilon\frac{\Sigma(\epsilon)}{\Gamma(1-\epsilon)},\qquad
\Sigma=2A_{12}-A_{11}-A_{22},\qquad A_{ij}=\int_{-1}^1\dd c\,(1-c^2)^{-\epsilon}\frac{p_i\cdot p_j}{(E_i-\vec p_i\cdot\hat n)(E_j-\vec p_j\cdot\hat n)}
\end{equation*}
になる（$|\vec k|$ の積分が $1/\epsilon$ を与える）。$O(\epsilon)$ の項は40桁の数値積分で求めた。$K_V+K_S$ の極はすべての過程で $10^{-38}$ 以下まで消え、$\mu$ への依存も $10^{-22}$ 以下だった。

(c) 硬い放射。\texttt{ir\_massive\_real.wl} は、実の放射の振幅の2乗 $R=\sum|\mathcal M_R|^2/(g^2C)$ を4次元で求める（$Z$ と $W$ の偏極の和は $-g+kk/M^2$、グルーオンは $-g$）。$R$ は本ノートの FeynArts の模型の $Z\to b\bar bg$ と $t\to bW^+g$ と記号として一致した。$E\to0$ では木と eikonal 因子の積に近づいた（差は $E^2$ に比例し、$E=10^{-6}E_{\max}$ で $10^{-13}$）。\texttt{ir\_massive\_mg5.py} では、MadGraph の木の $|\mathcal M|^2$（$Z\to b\bar bg$、$Z\to\tau^-\tau^+\gamma$、$t\to bW^+g$）と Dalitz の平面の21点で比べ、相対 $8\times10^{-12}$ 以内で一致した（手の式は PS.input に書いた運動量の不変量で40桁で評価）。$E>\Delta E$ の積分は、$x=M^2-s_{12}=2ME$ と $2p_i\cdot p_3$ の対数を変数にして、Gauss--Legendre の公式で求めた（図~\ref{fig:dalitz}）。初めは積分の精度が出なかった。原因は二つあった。一つは、$R$ を $s_{12}$、$s_{13}$ の式で評価すると、軟らかい端で大きな数の差になって桁を落とすことである。もう一つは、硬い端（$s_{12}\to(m_1+m_2)^2$）で Dalitz の領域の幅が $\sqrt{x_{\max}-x}$ で縮み、その平方根が公式の収束を損なうことである。$R$ を $x$ と $2p_1\cdot p_3$ で厳密に書き直し、硬い側を $x=x_{\max}-\frac{x_{\max}}2t^2$ で写した。パネルの数を16から24にしたときの差は $10^{-14}$ 程度（$m/M=10^{-5}$ の試験で $10^{-8}$）になった。

(d) $\delta$ によらないこと。$\delta$ を $10^{-4}$ から $10^{-8}$ まで変えると、$K_V+K_S$ と $K_H$ はそれぞれ90ほど変わる。和の $K$ の変化は軟らかい近似の誤差の $O(\delta)$ だけで、$\delta$ に比例して縮んだ（表~\ref{tab:massdelta}）。$\delta\to0$ の値は $K(10^{-8})-[K(10^{-7})-K(10^{-8})]/9$ とした。

\begin{figure}[tbp]
\centering
\begin{tikzpicture}[scale=1.0]
  \def\lowerB{(0.000,0.000) (0.289,0.008) (0.571,0.016) (0.848,0.025) (1.118,0.034) (1.382,0.044) (1.641,0.054) (1.893,0.065) (2.139,0.077) (2.379,0.089) (2.613,0.102) (2.841,0.116) (3.062,0.131) (3.278,0.146) (3.488,0.163) (3.691,0.181) (3.889,0.200) (4.080,0.220) (4.266,0.241) (4.445,0.265) (4.618,0.289) (4.785,0.316) (4.946,0.344) (5.101,0.374) (5.250,0.407) (5.393,0.442) (5.530,0.479) (5.660,0.519) (5.785,0.563) (5.903,0.609) (6.016,0.658) (6.122,0.712) (6.222,0.769) (6.316,0.829) (6.405,0.894) (6.487,0.964) (6.563,1.037) (6.632,1.115) (6.696,1.198) (6.754,1.285) (6.806,1.376) (6.851,1.471) (6.891,1.569) (6.924,1.671) (6.951,1.775) (6.973,1.881) (6.988,1.987) (6.997,2.094) (7.000,2.200)}
  \def\upperB{(7.000,2.200) (6.997,2.304) (6.988,2.405) (6.973,2.502) (6.951,2.594) (6.924,2.681) (6.891,2.762) (6.851,2.836) (6.806,2.902) (6.754,2.961) (6.696,3.011) (6.632,3.054) (6.563,3.088) (6.487,3.114) (6.405,3.131) (6.316,3.141) (6.222,3.143) (6.122,3.136) (6.016,3.123) (5.903,3.102) (5.785,3.073) (5.660,3.038) (5.530,2.996) (5.393,2.948) (5.250,2.893) (5.101,2.832) (4.946,2.765) (4.785,2.692) (4.618,2.614) (4.445,2.529) (4.266,2.440) (4.080,2.345) (3.889,2.245) (3.691,2.140) (3.488,2.029) (3.278,1.914) (3.062,1.794) (2.841,1.670) (2.613,1.540) (2.379,1.406) (2.139,1.268) (1.893,1.125) (1.641,0.977) (1.382,0.825) (1.118,0.669) (0.848,0.508) (0.571,0.343) (0.289,0.174) (0.000,0.000)}
  % the hard region and the soft strip
  \fill[black!7] plot coordinates \lowerB -- plot coordinates \upperB -- cycle;
  \begin{scope}
    \clip plot coordinates \lowerB -- plot coordinates \upperB -- cycle;
    \fill[blue!18] (0,0) rectangle (0.8,3.5);
  \end{scope}
  \draw[thick] plot coordinates \lowerB; \draw[thick] plot coordinates \upperB;
  % axes
  \draw[-{Stealth[length=2mm]}] (0,0) -- (7.7,0) node[right,font=\small] {$x$};
  \draw[-{Stealth[length=2mm]}] (0,0) -- (0,3.7) node[above,font=\small] {$w$};
  \node[font=\scriptsize,below] at (7.0,0) {$x_{\max}$}; \draw (7.0,0) -- (7.0,0.08);
  % the cut, the split at x_max/2, the boundary of the two maps in w
  \draw[dashed,blue!60!black] (0.8,0) -- (0.8,0.9);
  \node[font=\scriptsize,below,blue!60!black] at (0.8,0) {$2M\Delta E$};
  \draw[dashed] (3.5,0) -- (3.5,2.08); \node[font=\scriptsize,below] at (3.5,0) {$x_{\max}/2$};
  \draw[densely dotted] (0,0) -- (7.0,2.2);
  % labels
  \node[font=\scriptsize,blue!60!black,align=right,anchor=east] (soft) at (-0.3,0.9) {軟らかい領域\\$E<\Delta E$};
  \draw[-{Stealth[length=1.3mm]},blue!60!black] (soft.east) -- (0.5,0.17);
  \node[font=\scriptsize,align=center] at (5.1,2.2) {硬い領域：\\4次元の数値積分};
  \node[font=\scriptsize] (c1) at (5.3,-0.42) {$p_3\parallel p_1$：$w=w_{\mathrm{lo}}$};
  \draw[-{Stealth[length=1.3mm]}] (c1.north) -- (5.3,0.38);
  \node[font=\scriptsize,anchor=south east] (c2) at (3.3,2.15) {$p_3\parallel p_2$：$x-w=w'_{\mathrm{lo}}$};
  \draw[-{Stealth[length=1.3mm]}] (2.1,2.15) -- (2.2,1.37);
  \node[font=\scriptsize,anchor=west] at (2.75,0.55) {$w=w'$};
  \node[font=\scriptsize,align=left,anchor=west] at (7.15,2.2) {$s_{12}\to(m_1+m_2)^2$：\\幅が $\sqrt{x_{\max}-x}$ で閉じる};
  \draw[{Stealth[length=1.3mm]}-{Stealth[length=1.3mm]}] (0.8,-0.78) -- (3.5,-0.78);
  \draw[{Stealth[length=1.3mm]}-{Stealth[length=1.3mm]}] (3.5,-0.78) -- (7.0,-0.78);
  \node[font=\scriptsize,align=center] at (2.15,-1.05) {$\ln x$ の写像};
  \node[font=\scriptsize,align=center] at (5.25,-1.1) {$x=x_{\max}-\frac{x_{\max}}2t^2$};
\end{tikzpicture}
\caption{手順 (3) の位相空間（$Z\to Q\bar Qg$、$m=0.2M$ で描いた）。横軸は $x=M^2-s_{12}=2ME$（$E$ は放射のエネルギー）、縦軸は $w=2p_1\cdot p_3$。青い帯 $E<\Delta E=\delta E_{\max}$（$\delta$ は誇張した）は、式~\eqref{eq:eikonal} の積を $d$ 次元で解析的に積分する。残りは4次元で数値で積分する。下と上の縁は、放射が $p_1$、$p_2$ に平行な共線の端である。質量があるので、$w$、$w'=2p_2\cdot p_3$ は0にならない。数値の積分では、$x$ と、それぞれの縁から点線までの $w$、$w'$ の対数を変数にする。$x>x_{\max}/2$ では、領域の幅の $\sqrt{x_{\max}-x}$ を $t^2$ の写像で除く。}
\label{fig:dalitz}
\end{figure}

(e) 極限。質量0の極限（$m/M=10^{-5}$）では $K_V-\frac34=-4.5\times10^{-9}$、$K_A-\frac34=1.5\times10^{-8}$ で、手順 (2) の式~\eqref{eq:inclusive} に戻った。閾値（$\beta=0.05$、$0.01$）では $K_V-\pi^2/(2\beta)=-3.803$、$-3.954$ で、$\beta\to0$ で $-4$ に近づく。$\pi^2/(2\beta)$ は Coulomb の Sommerfeld の因子の1次で、$-4$ はベクトルの流れの硬い部分であり、よく知られた閾値の展開と合う。軸性ベクトルでは $K_A-\pi^2/(2\beta)=-1.804$、$-1.954$ だった。$t\to bW^+$ で $m_b$ と $M_W$ を $10^{-3}m_t$、$10^{-4}m_t$ にすると、$K$ と文献の極限 $-(\pi^2/3-5/4)$（$\Gamma^{(1)}/\Gamma_0=-\frac{2\alpha_s}{3\pi}(\frac{2\pi^2}3-\frac52)$~\cite{jk89}）との差は $5.8\times10^{-5}$、$7.9\times10^{-7}$ と縮んだ。

(f) MadLoop。\texttt{ir\_massive\_madloop.py} では、CKM ノートの NLO の UFO の \texttt{[virt=QCD]} で、$Z\to b\bar b$ と $t\to bW^+$ の仮想の補正を求めた。点は CKM ノートの MadLoop の検算の点（$\alpha_s=0.118$、$m_b=2.86$~GeV、$m_t=165$~GeV）、還元は COLLIER である。単極と有限部分（MadLoop の規格化。$(4\pi)^\epsilon/\Gamma(1-\epsilon)$ と $\mu_R^{2\epsilon}$ を除く）は、本ノートの $K_V$ と $\mu_R=91$、$500$~GeV で相対 $7\times10^{-15}$ 以内で一致した。

(g) 物理。$K$ を表~\ref{tab:massive} に、手順 (1) の木の幅に掛けた値（\texttt{ir\_massive\_widths.py}）を表~\ref{tab:massivewidths} に示す。$Z\to b\bar b$ では、$K$ を木の幅のベクトルと軸性ベクトルの割合（0.326 と 0.674）で平均すると 0.858 になる。QCD の補正は $4.26\%$ で、質量0の $\alpha_s/\pi=3.73\%$ より大きい。差のほとんどは軸性ベクトルの $K_A=0.897$ から来る。$\Gamma_{\mathrm{had}}$、$\Gamma_Z$、$R_\ell$ の測定値との差は、手順 (2) の $-0.29\%$、$-0.16\%$、$-0.25\%$ から $-0.16\%$、$-0.07\%$、$-0.12\%$ に縮んだ。一方 $R_b$ は、木の $+0.74\%$ から $+1.11\%$ に離れた。残りは手順 (5) の $Zb\bar b$ の頂点のトップの補正の大きさである。$t\to bW^+$ の QCD の補正は $\frac{\alpha_sC_F}\pi K=-0.7876\alpha_s$ で、$\alpha_s(m_t)=0.10807$ では $-8.51\%$ になる。$\Gamma_t=1.3299$~GeV で、測定値 $1.42^{+0.19}_{-0.15}$~GeV と $0.6\sigma$ で合う。$m_b=M_W=0$ の極限の $-0.866\alpha_s$ より $9\%$ 小さい。文献~\cite{cm98}（参照）は $m_t\approx175$~GeV で $-8.4\%$、物理的な質量では極限より約 $10\%$ 小さいとしており、これと合う。

\begin{table}[tbp]
\centering\small
\caption{手順 (3)：$\delta$ によらないこと（\texttt{ir\_massive\_real.wl}）。$K_V+K_S$ は解析式、$K_H$ は数値の積分で、$\Delta E=\delta E_{\max}$。$\delta\to0$ は本文 (d) の外挿。}
\label{tab:massdelta}
\begin{tabular}{lrrrrrr}
\toprule
 & \multicolumn{3}{c}{$Z\to b\bar b$（ベクトル）} & \multicolumn{3}{c}{$t\to bW^+$} \\
$\delta$ & $K_V+K_S$ & $K_H$ & $K$ & $K_V+K_S$ & $K_H$ & $K$ \\
\midrule
$10^{-4}$ & $-80.015136$ & $80.791936$ & $0.7767996$ & $-40.790963$ & $38.935751$ & $-1.8552118$ \\
$10^{-6}$ & $-125.067777$ & $125.843417$ & $0.7756395$ & $-62.435069$ & $60.579292$ & $-1.8557771$ \\
$10^{-8}$ & $-170.120418$ & $170.896046$ & $0.7756279$ & $-84.079175$ & $82.223392$ & $-1.8557827$ \\
$\to0$ & & & $0.7756278$ & & & $-1.8557828$ \\
\bottomrule
\end{tabular}
\end{table}

\begin{table}[tbp]
\centering\small
\caption{手順 (3)：$\Gamma^{(1)}/\Gamma_0=\frac{\alpha C}\pi K$ の $K$（$\delta\to0$。極質量。QCD では $C=C_F$、QED では $Q_f^2$）。「平均」は木の幅のベクトルと軸性ベクトルの割合での平均。下の段は極限の試験（$M=1$）。}
\label{tab:massive}
\begin{tabular}{llllll}
\toprule
過程 & 質量 & $K_V$ & $K_A$ & 平均 & 極限の値 \\
\midrule
$Z\to b\bar b$ & $m_b=4.78$ & 0.775628 & 0.897293 & 0.857602 & \\
$Z\to c\bar c$ & $m_c=1.67$ & 0.753036 & 0.774179 & 0.771504 & \\
$Z\to\tau^-\tau^+$ & $m_\tau=1.77693$ & 0.753440 & 0.776955 & 0.776839 & \\
$Z\to\mu^-\mu^+$ & $m_\mu=0.10566$ & 0.750012 & 0.750163 & 0.750163 & \\
$t\to bW^+$ & $m_b=4.78$ & \multicolumn{3}{l}{$K=-1.855783$} & \\
\midrule
質量0 & $m=10^{-5}$ & 0.7500000 & 0.7500000 & & $\frac34$ \\
閾値 & $\beta=0.05$ & 94.89277 & 96.89160 & & $\frac{\pi^2}{2\beta}-4=94.69604$（$V$） \\
閾値 & $\beta=0.01$ & 489.52671 & 491.52666 & & $\frac{\pi^2}{2\beta}-4=489.48022$（$V$） \\
$t\to bW^+$ & $m_b=M_W=10^{-4}m_t$ & \multicolumn{3}{l}{$K=-2.039867$} & $-(\frac{\pi^2}3-\frac54)=-2.039868$ \\
\bottomrule
\end{tabular}
\end{table}

\begin{table}[tbp]
\centering\small
\caption{手順 (3)：質量を入れた補正を木の幅（表~\ref{tab:lo}）に掛けた値（\texttt{ir\_massive\_widths.py}。GeV）。手順 (2) は質量0の補正（表~\ref{tab:inclusive}）。測定値は PDG 2024~\cite{pdg24}（参照）。括弧は測定値との差。}
\label{tab:massivewidths}
\begin{tabular}{lllll}
\toprule
量 & 木 & 手順 (2) & 手順 (3) & 測定値（参照） \\
\midrule
$\Gamma(Z\to e^+e^-)$ & 0.083814（$-0.20\%$） & 0.083971（$-0.02\%$） & 0.083971（$-0.02\%$） & $0.083984\pm0.000086$ \\
$\Gamma(Z\to\mathrm{hadrons})$ & 1.676229（$-3.91\%$） & 1.739392（$-0.29\%$） & 1.741668（$-0.16\%$） & $1.7444\pm0.0020$ \\
$\Gamma_Z$ & 2.427894（$-2.71\%$） & 2.491527（$-0.16\%$） & 2.493808（$-0.07\%$） & $2.4955\pm0.0023$ \\
$R_\ell$ & 19.9993（$-3.70\%$） & 20.7142（$-0.25\%$） & 20.7413（$-0.12\%$） & $20.767\pm0.025$ \\
$R_b$ & 0.217898（$+0.74\%$） & — & 0.218696（$+1.11\%$） & $0.21629\pm0.00066$ \\
$R_c$ & 0.170545（$-0.90\%$） & — & 0.170569（$-0.89\%$） & $0.1721\pm0.0030$ \\
$\Gamma_t$ & 1.453649（$+2.37\%$） & — & 1.329919（$-6.34\%$） & $1.42^{+0.19}_{-0.15}$ \\
\bottomrule
\end{tabular}
\end{table}

\FloatBarrier
\paragraph{手順 (4a)：$Z\to f\bar f$ の弱い補正。}手順 (4) は $Z$、$W$、$t$ の崩壊に分け、ここでは (4a) として $Z\to f\bar f$ を扱う。$Z$ の崩壊では、光子の部分（光子の交換、光子による場の定数、実の光子の放射）はそれだけでゲージによらない組をなし、手順 (2)、(3) の $K$（$C=Q_f^2$）に等しい。グルーオンの部分も同じである。残りの弱い部分（$Z$、$W$、$H$ と Goldstone のループ、ボソンの自己エネルギーと $Z$--光子の混合、対抗項）は赤外で有限で、$\Gamma_f=\Gamma_0[1+\delta_{\mathrm{weak}}+\frac{\alpha Q_f^2}\pi K+\frac{\alpha_sC_F}\pi K]$ の $\delta_{\mathrm{weak}}$ を与える。

\begin{figure}[tbp]
\centering
\begin{tikzpicture}[scale=0.85]
% (a) the ZWW vertex
\begin{scope}[shift={(0,0)}]
  \draw[photon] (-1.1,0) -- (0,0);
  \draw[photon] (0,0) -- (0.6,0.45); \draw[photon] (0,0) -- (0.6,-0.45);
  \draw[fermion] (0.6,-0.45) -- (0.6,0.45);
  \draw[fermion] (0.6,0.45) -- (1.2,0.8); \draw[fermion] (1.2,-0.8) -- (0.6,-0.45);
  \node[vtx] at (0,0) {}; \node[vtx] at (0.6,0.45) {}; \node[vtx] at (0.6,-0.45) {};
  \node[font=\scriptsize] at (-0.9,0.25) {$Z$}; \node[font=\scriptsize] at (0.12,0.48) {$W$}; \node[font=\scriptsize] at (0.12,-0.48) {$W$};
  \node[font=\scriptsize] at (0.85,0) {$f'$}; \node[font=\scriptsize] at (1.4,0.85) {$f$}; \node[font=\scriptsize] at (1.4,-0.85) {$\bar f$};
  \node[slab] at (0.1,-1.45) {(a) $ZWW$ の頂点};
\end{scope}
% (b) the W exchange
\begin{scope}[shift={(3.6,0)}]
  \draw[photon] (-1.1,0) -- (0,0);
  \draw[fermion] (0,0) -- (0.6,0.42); \draw[fermion] (0.6,0.42) -- (1.15,0.8);
  \draw[fermion] (1.15,-0.8) -- (0.6,-0.42); \draw[fermion] (0.6,-0.42) -- (0,0);
  \draw[photon] (0.6,0.42) -- (0.6,-0.42);
  \node[vtx] at (0,0) {}; \node[vtx] at (0.6,0.42) {}; \node[vtx] at (0.6,-0.42) {};
  \node[font=\scriptsize] at (-0.9,0.25) {$Z$}; \node[font=\scriptsize] at (0.2,0.45) {$f'$}; \node[font=\scriptsize] at (0.2,-0.45) {$\bar f'$};
  \node[font=\scriptsize] at (0.88,0) {$W$}; \node[font=\scriptsize] at (1.35,0.85) {$f$}; \node[font=\scriptsize] at (1.35,-0.85) {$\bar f$};
  \node[slab] at (0.1,-1.45) {(b) $W$ の交換（$Z$、$H$ も）};
\end{scope}
% (c) the Z-photon mixing
\begin{scope}[shift={(7.6,0)}]
  \draw[photon] (-1.5,0) -- (-0.92,0); \draw[photon] (-0.58,0) -- (0,0);
  \node[blob,minimum size=10pt] at (-0.75,0) {};
  \draw[fermion] (0,0) -- (1.15,0.75); \draw[fermion] (1.15,-0.75) -- (0,0);
  \node[vtx] at (0,0) {};
  \node[font=\scriptsize] at (-1.3,0.25) {$Z$}; \node[font=\scriptsize] at (-0.28,0.25) {$\gamma$};
  \node[font=\scriptsize] at (-0.75,-0.35) {$\Sigma^{AZ}$};
  \node[font=\scriptsize] at (1.35,0.8) {$f$}; \node[font=\scriptsize] at (1.35,-0.8) {$\bar f$};
  \node[slab] at (-0.1,-1.45) {(c) $Z$--光子の混合};
\end{scope}
% (d) the residue of the Z
\begin{scope}[shift={(11.6,0)}]
  \draw[photon] (-1.5,0) -- (-0.92,0); \draw[photon] (-0.58,0) -- (0,0);
  \node[blob,minimum size=10pt] at (-0.75,0) {};
  \draw[fermion] (0,0) -- (1.15,0.75); \draw[fermion] (1.15,-0.75) -- (0,0);
  \node[vtx] at (0,0) {};
  \node[font=\scriptsize] at (-1.3,0.25) {$Z$}; \node[font=\scriptsize] at (-0.28,0.25) {$Z$};
  \node[font=\scriptsize] at (-0.75,-0.35) {$\Sigma^{ZZ}$};
  \node[font=\scriptsize] at (1.35,0.8) {$f$}; \node[font=\scriptsize] at (1.35,-0.8) {$\bar f$};
  \node[slab] at (-0.1,-1.45) {(d) $Z$ の留数};
\end{scope}
% (e) the field constant of a fermion leg
\begin{scope}[shift={(1.2,-3.2)}]
  \draw[photon] (-1.1,0) -- (0,0);
  \draw[fermion] (0,0) -- (0.35,0.245); \draw[fermion] (0.35,0.245) -- (0.9,0.63); \draw[fermion] (0.9,0.63) -- (1.2,0.84);
  \draw[fermion] (1.2,-0.84) -- (0,0);
  \draw[photon] (0.35,0.245) .. controls (0.3,0.85) and (0.75,1.15) .. (0.9,0.63);
  \node[vtx] at (0,0) {}; \node[vtx] at (0.35,0.245) {}; \node[vtx] at (0.9,0.63) {};
  \node[font=\scriptsize] at (-0.9,0.25) {$Z$}; \node[font=\scriptsize] at (0.3,1.05) {$W$, $Z$}; \node[font=\scriptsize] at (0.85,0.32) {$f'$};
  \node[font=\scriptsize] at (1.4,0.9) {$f$}; \node[font=\scriptsize] at (1.4,-0.9) {$\bar f$};
  \node[slab] at (0.1,-1.45) {(e) 脚の場の定数 $\delta Z_f$};
\end{scope}
% (f) the top in the Z b bbar vertex
\begin{scope}[shift={(5.6,-3.2)}]
  \draw[photon] (-1.1,0) -- (0,0);
  \draw[fermion] (0,0) -- (0.6,0.42); \draw[fermion] (0.6,0.42) -- (1.15,0.8);
  \draw[fermion] (1.15,-0.8) -- (0.6,-0.42); \draw[fermion] (0.6,-0.42) -- (0,0);
  \draw[scalar] (0.6,0.42) -- (0.6,-0.42);
  \node[vtx] at (0,0) {}; \node[vtx] at (0.6,0.42) {}; \node[vtx] at (0.6,-0.42) {};
  \node[font=\scriptsize] at (-0.9,0.25) {$Z$}; \node[font=\scriptsize] at (0.2,0.45) {$t$}; \node[font=\scriptsize] at (0.2,-0.45) {$\bar t$};
  \node[font=\scriptsize] at (0.95,0) {$G^\pm$}; \node[font=\scriptsize] at (1.35,0.85) {$b$}; \node[font=\scriptsize] at (1.35,-0.85) {$\bar b$};
  \node[slab] at (0.1,-1.45) {(f) $Zb\bar b$ のトップ（$\propto y_t^2$）};
\end{scope}
% (g) the top with two W/G
\begin{scope}[shift={(10.2,-3.2)}]
  \draw[photon] (-1.1,0) -- (0,0);
  \draw[scalar] (0,0) -- (0.6,0.45); \draw[scalar] (0,0) -- (0.6,-0.45);
  \draw[fermion] (0.6,-0.45) -- (0.6,0.45);
  \draw[fermion] (0.6,0.45) -- (1.2,0.8); \draw[fermion] (1.2,-0.8) -- (0.6,-0.45);
  \node[vtx] at (0,0) {}; \node[vtx] at (0.6,0.45) {}; \node[vtx] at (0.6,-0.45) {};
  \node[font=\scriptsize] at (-0.9,0.25) {$Z$}; \node[font=\scriptsize] at (0.1,0.5) {$G^-$}; \node[font=\scriptsize] at (0.1,-0.5) {$G^+$};
  \node[font=\scriptsize] at (0.85,0) {$t$}; \node[font=\scriptsize] at (1.4,0.85) {$b$}; \node[font=\scriptsize] at (1.4,-0.85) {$\bar b$};
  \node[slab] at (0.1,-1.45) {(g) $Zb\bar b$ のトップ（$G^\pm$ の対）};
\end{scope}
\end{tikzpicture}
\caption{手順 (4a) の $Z\to f\bar f$ の弱い補正の図の例（光子とグルーオンを除く。波線はゲージ場、破線は Goldstone 場）。(a)、(b) は 1PI の頂点（1PI の図はニュートリノで11個、荷電レプトンで25個、クォークで33個）。(c) は式~\eqref{eq:zren} の $Z$--光子の混合で、光子の木を通して $\mathcal V_f$ だけをずらし、手順 (5) の $\hat\kappa_f$ の主な部分になる。(d) の $Z$ の留数は $F_L$、$F_R$ に共通で、$\mathcal V_f/\mathcal A_f$ からは消える。(e) はフェルミオンの脚の場の定数（式~\eqref{eq:zren} の $\delta Z_f^L$）。(f)、(g) は $b$ に固有の、内線のトップと $W$、$G^\pm$ の頂点で（$W$ の線を含む同じ形の図もある）、$y_t^2$ に比例して $Z\to b\bar b$ を $-1.44\%$ ずらし、$R_b$ を測定値に戻す。}
\label{fig:znlo}
\end{figure}

(a) 振幅。\texttt{zff\_feyncalc.wl} は、本ノートの FeynArts の模型で、光子とグルーオンを除いた 1PI の頂点を一般の $R_\xi$ ゲージで求める（図はニュートリノで11個、荷電レプトンで25個、クォークで33個。例を図~\ref{fig:znlo} に示す）。形状因子は $\mathcal M=g_Z\bar u[\slashed\varepsilon(F_LP_L+F_RP_R)+(\varepsilon\cdot p_1)(H_LP_L+H_RP_R)]v$ と書く（木では $F_L=T_3-Q\hat s^2$、$F_R=-Q\hat s^2$）。同じスクリプトで、そのフェルミオンの弱い自己エネルギー（5〜13個の図）も求める。$t$ を除く11個のフェルミオンを並列に計算する（約6分）。

(b) S 行列要素。\texttt{zweak.wl} は、SM ノートと CKM ノートと同じく、裸の切断された頂点に on-shell の LSZ の因子を掛けて S 行列要素にする。
\begin{equation}
F_L=F_L^{(0)}+F_L^{\mathrm{vtx}}+F_L^{(0)}\Bigl[-\tfrac12\Sigma^{ZZ\prime}_T(M_Z^2)+\delta Z^L_f\Bigr]-\frac{\Sigma^{AZ}_T(M_Z^2)}{M_Z^2}F_L^{\gamma}+\delta F_L^{\MSbar}
\label{eq:zren}
\end{equation}
で、$F_R$ も同じ形である。$\delta Z^L_f=-\Sigma^L-m^2(\Sigma^{L\prime}+\Sigma^{R\prime}+2\Sigma^{S\prime})$ は二つの脚の和である。$F^\gamma$ は混合 $Z\to\gamma^*\to f\bar f$ の光子の木の振幅で、同じ単位では $Q\hat s\hat c$ になる（模型の符号ごと求めた）。$\delta F^{\MSbar}$ は $g$、$g'$ の $\MSbar$ の対抗項（$\beta$ 関数の極）である。$\Sigma^{ZZ}$、$\Sigma^{AZ}$ には SM ノートの手順 (3) のものを使った（一般の $R_\xi$、すべてのフェルミオンの質量を含み、1ループでは $V$ によらない）。FJ のタドポールは質量をずらすだけで、この S 行列要素には入らない。

(c) 紫外の有限性。11個のフェルミオンで、$F_L$、$F_R$ の紫外の極は記号として消えた（質量とゲージの母数は記号のまま、$V$ は CKM ノートの有理数のユニタリーな行列）。混合の項を除くか、光子の木の符号を逆にすると、8個の荷電フェルミオンで極が残る。

(d) ゲージによらないこと。$\delta_{\mathrm{weak}}$ は、Package-X（FeynHelpers）のループ関数で、本ノートの $\mu=M_Z$ の $\MSbar$ の結合と極質量を使って求めた（ループの中では木の関係 $M_W=M_Z\hat c$）。$(\xi_Z,\xi_W)=(1,1)$、$(2,3)$、$(1/2,4)$ の間の差は、レプトンで $10^{-47}$、$c$ と $b$ で $10^{-15}$、$u$、$d$、$s$ で $10^{-7}$ 以下だった。$u$、$d$、$s$ で桁が少ないのは、CKM 行列の倍精度の値と、軽いクォークの自己エネルギーの打ち消しのためである。$u$、$d$、$s$ には 0.1~MeV の質量を与えた。保存した自己エネルギーの項の分母に質量があり、0 では 0/0 になるためで、$\delta_{\mathrm{weak}}$ の変化は $m^2/M_Z^2\sim10^{-12}$ である。

(e) MadLoop。\texttt{zweak\_madloop.py} では、CKM ノートの NLO の UFO の \texttt{[virt=QED]} で、$Z\to\nu_\tau\bar\nu_\tau$、$\tau^-\tau^+$、$b\bar b$ の電弱の仮想の補正を、CKM ノートの MadLoop の検算の点で求めた。これは弱い部分と光子の部分の和なので、同じ点の $\delta_{\mathrm{weak}}$ と、手順 (3) の光子の仮想の補正 $\frac{\alpha Q_f^2}\pi K_V$ の和と比べた。有限部分は $\mu_R=91$、$500$~GeV で相対 $4\times10^{-13}$ 以内、単極は $10^{-16}$ 以内で一致した。

(f) $\mu$ への依存。結合を $\mu=M_Z/2$、$2M_Z$ に走らせて $\Gamma_0(1+\delta_{\mathrm{weak}})$ を求めた。木の幅の $\pm0.7$〜$1.2\%$ の変化は、$\Gamma_{\ell\ell}$、$\Gamma_{\mathrm{inv}}$、$\Gamma_{\mathrm{had}}$、$\Gamma_Z$ のすべてで $0.04\%$ 以下になった。QED と QCD の因子まで入れると、ハドロンの幅に $\pm0.4\%$ が残る（表~\ref{tab:zpole}）。これは QCD の因子の $\alpha_s(\mu)$ の走りで、$O(\alpha_s^2)$ の補正で打ち消されるはずの分である。

(g) 物理。$\delta_{\mathrm{weak}}$ を表~\ref{tab:zweak} に、それを入れた $Z$ 極の量を表~\ref{tab:zpole} に示す（\texttt{zweak\_widths.py}）。$b$ の $-1.44\%$ は、$W$、$G^\pm$ とトップのループによる $Zb\bar b$ の頂点の補正である。他のダウン型クォーク（$+0.19\%$）との差は $y_t^2$ に比例する非普遍的な補正で、これで $R_b$ は手順 (3) の $+1.11\%$ から $-0.09\%$（$0.3\sigma$）に戻った。$\sigma^0_{\mathrm{had}}$ は $+0.04\%$（$0.5\sigma$）で合う。見えない幅は $+0.57\%$（$1.9\sigma$）で、測定値が標準模型の値より小さいことは知られている（ニュートリノの数で $N_\nu=2.984\pm0.008$~\cite{lepewwg06}）。$\Gamma_Z$ と $R_\ell$ は $-0.15\%$（$-1.6\sigma$）と $-0.31\%$（$-2.6\sigma$）である。ハドロンの幅の $-0.35\%$ の大部分は $O(\alpha_s^2)$ の QCD の補正の大きさで（よく知られた係数で $1.41(\alpha_s/\pi)^2\approx0.20\%$。参照）、本ノートの NLO の範囲の外にある。

\begin{table}[tbp]
\centering\small
\caption{手順 (4a)：$Z\to f\bar f$ の弱い補正 $\delta_{\mathrm{weak}}$（$\%$。\texttt{zweak.wl}。$\mu$ の結合と極質量。光子とグルーオンの部分は含まない）。ゲージの母数を変えた差は本文 (d) の通り。}
\label{tab:zweak}
\begin{tabular}{lrrr|lrrr}
\toprule
$f$ & $\mu=M_Z/2$ & $M_Z$ & $2M_Z$ & $f$ & $\mu=M_Z/2$ & $M_Z$ & $2M_Z$ \\
\midrule
$\nu_e$、$\nu_\mu$ & $-0.4632$ & $+0.2836$ & $+0.9928$ & $u$ & $-1.2925$ & $+0.0823$ & $+1.4048$ \\
$\nu_\tau$ & $-0.4637$ & $+0.2832$ & $+0.9924$ & $c$ & $-1.2928$ & $+0.0831$ & $+1.4067$ \\
$e$、$\mu$ & $-0.9851$ & $-0.0051$ & $+0.8630$ & $d$ & $-0.9934$ & $+0.1888$ & $+1.3368$ \\
$\tau$ & $-0.9850$ & $-0.0045$ & $+0.8638$ & $s$ & $-0.9964$ & $+0.1858$ & $+1.3337$ \\
 & & & & $b$ & $-2.6322$ & $-1.4382$ & $-0.2786$ \\
\bottomrule
\end{tabular}
\end{table}

\begin{table}[tbp]
\centering\footnotesize\setlength{\tabcolsep}{4pt}
\caption{手順 (4a)：$Z$ 極の量（\texttt{zweak\_widths.py}。GeV、$\sigma^0_{\mathrm{had}}$ は nb）。「QED、QCD」は手順 (3) の値、NLO はさらに $\delta_{\mathrm{weak}}$ を入れた値。括弧は測定値（PDG 2024~\cite{pdg24}、参照）との差。右の2列は $\mu=M_Z/2$、$2M_Z$ にしたときの変化（$\%$）。}
\label{tab:zpole}
\begin{tabular}{lllll|ll}
\toprule
量 & 木 & QED、QCD & NLO & 測定値（参照） & 木 & NLO \\
\midrule
$\Gamma_{\ell\ell}$ & 0.083814（$-0.20\%$） & 0.083971（$-0.02\%$） & 0.083967（$-0.02\%$） & $0.083984(86)$ & $+0.95$、$-0.89$ & $-0.04$、$-0.03$ \\
$\Gamma_{\mathrm{inv}}$ & 0.500412（$+0.28\%$） & 0.500412（$+0.28\%$） & 0.501831（$+0.57\%$） & $0.4990(15)$ & $+0.74$、$-0.71$ & $-0.01$、$-0.01$ \\
$\Gamma_{\mathrm{had}}$ & 1.676229（$-3.91\%$） & 1.741668（$-0.16\%$） & 1.738272（$-0.35\%$） & $1.7444(20)$ & $+1.24$、$-1.20$ & $+0.39$、$-0.36$ \\
$\Gamma_Z$ & 2.427894（$-2.71\%$） & 2.493808（$-0.07\%$） & 2.491818（$-0.15\%$） & $2.4955(23)$ & $+1.11$、$-1.07$ & $+0.27$、$-0.25$ \\
$R_\ell$ & 19.9993（$-3.70\%$） & 20.7413（$-0.12\%$） & 20.7019（$-0.31\%$） & $20.767(25)$ & $+0.29$、$-0.32$ & $+0.43$、$-0.33$ \\
$R_b$ & 0.217898（$+0.74\%$） & 0.218696（$+1.11\%$） & 0.216101（$-0.09\%$） & $0.21629(66)$ & $-0.06$、$+0.06$ & $+0.03$、$-0.02$ \\
$R_c$ & 0.170545（$-0.90\%$） & 0.170569（$-0.89\%$） & 0.171039（$-0.62\%$） & $0.1721(30)$ & $+0.12$、$-0.12$ & $+0.00$、$-0.01$ \\
$\sigma^0_{\mathrm{had}}$ & 42.0751（$+1.43\%$） & 41.5147（$+0.08\%$） & 41.4979（$+0.04\%$） & $41.481(33)$ & $-0.02$、$+0.05$ & $-0.18$、$+0.12$ \\
\bottomrule
\end{tabular}
\end{table}

\FloatBarrier
\paragraph{手順 (4b)：$W^+\to f\bar f'$ の電弱の補正。}$W$ は電荷を持つので、光子は $f$、$\bar f'$ だけでなく $W$ にも付く（図~\ref{fig:wnlo}）。$Z$ のように光子の部分だけを取り出すと、それはゲージによらない組にならない。そこで、グルーオンを除くすべてのループ（光子を含む）の仮想の補正と実の光子の放射を足して、光子について包括的な電弱の補正 $\delta_{\mathrm{EW}}$ を求め、$\Gamma_f=\Gamma_0[1+\delta_{\mathrm{EW}}+\frac{\alpha_s}\pi]$ とする（$\frac{\alpha_s}\pi$ はクォークだけ）。グルーオンは $W$ に付かないので、QCD の因子には手順 (2) の質量0の $\frac34C_F\frac{\alpha_s}\pi=\frac{\alpha_s}\pi$ をそのまま使う。

\begin{figure}[tbp]
\centering
\begin{tikzpicture}[scale=0.85]
% (a) photon from the W
\begin{scope}[shift={(0,0)}]
  \draw[photon] (-1.15,0) -- (-0.55,0); \draw[photon] (-0.55,0) -- (0,0);
  \draw[fermion] (0,0) -- (1.15,0.75); \draw[fermion] (1.15,-0.75) -- (0,0);
  \draw[photon] (-0.55,0) -- (-0.1,0.95);
  \node[vtx] at (0,0) {}; \node[vtx] at (-0.55,0) {};
  \node[font=\scriptsize] at (-1.0,0.24) {$W^+$}; \node[font=\scriptsize] at (1.35,0.8) {$\nu_\ell$}; \node[font=\scriptsize] at (1.4,-0.8) {$\ell^+$};
  \node[font=\scriptsize] at (0.15,1.05) {$\gamma$};
  \node[slab] at (0.1,-1.6) {(a) 実の放射};
\end{scope}
% (b) photon from the lepton
\begin{scope}[shift={(3.4,0)}]
  \draw[photon] (-1.1,0) -- (0,0);
  \draw[fermion] (0,0) -- (1.15,0.75); \draw[fermion] (1.2,-0.55) -- (0.5,-0.35); \draw[fermion] (0.5,-0.35) -- (0,0);
  \draw[photon] (0.5,-0.35) -- (0.95,-1.1);
  \node[vtx] at (0,0) {}; \node[vtx] at (0.5,-0.35) {};
  \node[font=\scriptsize] at (-0.9,0.24) {$W^+$}; \node[font=\scriptsize] at (1.35,0.8) {$\nu_\ell$}; \node[font=\scriptsize] at (1.45,-0.55) {$\ell^+$};
  \node[font=\scriptsize] at (1.2,-1.15) {$\gamma$};
  \node[slab] at (0.1,-1.6) {(b) 実の放射};
\end{scope}
% (c) the Goldstone diagram
\begin{scope}[shift={(6.8,0)}]
  \draw[photon] (-1.15,0) -- (-0.55,0); \draw[scalar] (-0.55,0) -- (0,0);
  \draw[fermion] (0,0) -- (1.15,0.75); \draw[fermion] (1.15,-0.75) -- (0,0);
  \draw[photon] (-0.55,0) -- (-0.1,0.95);
  \node[vtx] at (0,0) {}; \node[vtx] at (-0.55,0) {};
  \node[font=\scriptsize] at (-1.0,0.24) {$W^+$}; \node[font=\scriptsize] at (-0.25,-0.25) {$G^+$};
  \node[font=\scriptsize] at (1.35,0.8) {$\nu_\ell$}; \node[font=\scriptsize] at (1.4,-0.8) {$\ell^+$};
  \node[font=\scriptsize] at (0.15,1.05) {$\gamma$};
  \node[slab] at (0.1,-1.6) {(c) 実の放射（$\propto m_\ell$）};
\end{scope}
% (d) the vertex with the photon between the W and the lepton
\begin{scope}[shift={(10.2,0)}]
  \draw[photon] (-1.15,0) -- (-0.55,0); \draw[photon] (-0.55,0) -- (0,0);
  \draw[fermion] (0,0) -- (1.15,0.75);
  \draw[fermion] (1.15,-0.8) -- (0.6,-0.42); \draw[fermion] (0.6,-0.42) -- (0,0);
  \draw[photon] (-0.55,0) .. controls (-0.45,-0.75) and (0.2,-0.85) .. (0.6,-0.42);
  \node[vtx] at (0,0) {}; \node[vtx] at (-0.55,0) {}; \node[vtx] at (0.6,-0.42) {};
  \node[font=\scriptsize] at (-1.0,0.24) {$W^+$}; \node[font=\scriptsize] at (1.35,0.8) {$\nu_\ell$}; \node[font=\scriptsize] at (1.4,-0.85) {$\ell^+$};
  \node[font=\scriptsize] at (-0.25,0.22) {$W$}; \node[font=\scriptsize] at (-0.25,-0.95) {$\gamma$};
  \node[slab] at (0.1,-1.6) {(d) 頂点の補正};
\end{scope}
% (e) the photon loop of the W field constant
\begin{scope}[shift={(13.6,0)}]
  \draw[photon] (-1.55,0) -- (-1.1,0); \draw[photon] (-0.35,0) -- (0,0);
  \draw[photon] (-1.1,0) arc (180:0:0.375); \draw[photon] (-1.1,0) arc (180:360:0.375);
  \draw[fermion] (0,0) -- (1.15,0.75); \draw[fermion] (1.15,-0.75) -- (0,0);
  \node[vtx] at (0,0) {}; \node[vtx] at (-1.1,0) {}; \node[vtx] at (-0.35,0) {};
  \node[font=\scriptsize] at (-1.45,0.24) {$W^+$}; \node[font=\scriptsize] at (-0.725,0.62) {$\gamma$}; \node[font=\scriptsize] at (-0.725,-0.62) {$W$};
  \node[font=\scriptsize] at (1.35,0.8) {$\nu_\ell$}; \node[font=\scriptsize] at (1.4,-0.8) {$\ell^+$};
  \node[slab] at (0.0,-1.6) {(e) $\delta Z_W$};
\end{scope}
\end{tikzpicture}
\caption{手順 (4b) の図の例（$W^+\to\nu_\ell\ell^+$）。(a)〜(c) は木の実の光子の放射（Feynman ゲージ。クォークでは $u_i$ からの放射が加わって4個）。(c) の $G^+\ell\nu$ の結合は $m_\ell$ に比例し、ユニタリーゲージ（MadGraph）では (a) の $W$ の縦波の部分に入る。(d) は光子の頂点の補正の一つ（グルーオンを除く 1PI の図はレプトンで14個、クォークで46個）、(e) は $W$ の場の定数の光子のループで、$\Sigma^{WW\prime}_T(M_W^2)$ に赤外の極を与える。(a)、(d)、(e) は光子が $W$ に付く図で、そのために光子の部分は弱い部分と分けられない。}
\label{fig:wnlo}
\end{figure}

(a) 振幅。\texttt{wsig\_feyncalc.wl} は、$W$ の自己エネルギー（29個の 1PI の図。光子は質量0）を横の部分に射影して求める。$V=1$、$\xi_A=1$ では、SM ノートの手順 (3) の $\Sigma^{WW}_T$（光子の質量 $M_A$ を入れたもの）の $M_A\to0$ と記号として一致した。\texttt{wff\_feyncalc.wl} は、グルーオンを除いた 1PI の頂点を一般の $R_\xi$ ゲージで求める。形状因子は $\mathcal M=\frac{e}{\sqrt2\hat s}\bar u[\slashed\varepsilon(F_LP_L+F_RP_R)+(\varepsilon\cdot p_1)(H_LP_L+H_RP_R)]v$ と書き（木では $F_L=V_{ij}$、レプトンでは1）、$F_R$、$H_L$、$H_R$ の紫外の極が記号として0であることを確かめた。9個のチャネルを並列に計算する（約11分）。

(b) S 行列要素。\texttt{wweak.wl} は、手順 (4a) と同じく、裸の頂点に on-shell の LSZ の因子を掛ける。
\begin{equation}
F_L=V_{ij}+F_L^{\mathrm{vtx}}-\tfrac12V_{ij}\Sigma^{WW\prime}_T(M_W^2)+\tfrac12\bigl[(\delta Z^{uL})^\dagger V+V\delta Z^{dL}\bigr]_{ij}+V_{ij}\frac{\delta g}g+\delta V_{ij}
\label{eq:wren}
\end{equation}
で、$F_R$、$H_L$、$H_R$ は頂点だけからなる。$\Sigma^{WW\prime}_T(M_W^2)$ は図~\ref{fig:wnlo}(e) の $B_0'(M_W^2;0,M_W^2)$ を含み、赤外で発散する（次元正則化で $-\frac1{2M_W^2}[\frac1\epsilon+\ln\frac{\mu^2}{M_W^2}+2]$。$\MSbar$ の規格化）。$\delta Z^{fL}$ の対角は手順 (4a) の弱い部分と手順 (3) の光子の部分（赤外の極を残す）の和で、非対角は CKM ノートのクォークの自己エネルギーから Denner の形で作った。質量0のニュートリノの脚では $p^2\to0$ の $p^0$ の係数をとる。$\delta g/g$ は $g$ の $\MSbar$ の対抗項（$\beta$ 関数の極）、$\delta V$ は CKM ノートの式 (10) の $\MSbar$ の対抗項である。$u$、$d$、$s$ には CKM ノートの基準の質量（1.27、2.71、54.6~MeV）を与えた。共線の対数を切るためで（包括的な幅では消える）、非対角の場の定数にも異なる質量が要る。木の幅は質量0のままである。

(c) 紫外と赤外の極。紫外と赤外の発散は同じ $\epsilon$ に入るので、仮想の補正の単極が軟らかい光子の極と打ち消すことが、紫外の有限性の確かめも兼ねる。軟らかい光子は手順 (3) の $K_S$ で、$\Sigma$ を $W$ の電荷まで広げて $\Sigma=-\sum_{ij}\eta_i\eta_je_ie_jA_{ij}$ とする（入る $W$ は $\eta=-1$、出る $f$、$\bar f'$ は $+1$。$e_i$ は電荷で、$e_W=e_f+e_{\bar f'}$）。$A_{ij}$ は $W$ の静止系の角度の積分で、$O(\epsilon)$ の項を閉じた形にした（$\int_{-1}^1\dd c\,\ln(1-c^2)/(1-bc)$ と $\ln(1-c^2)/(1-bc)^2$。$1-b=8\times10^{-11}$ まで数値の積分と $10^{-30}$ で一致）。9個のチャネルと $\mu=M_W/2$、$M_W$、$2M_W$ で、極の和は $3\times10^{-15}$ 以下だった。

(d) ゲージによらないこと。$(\xi_Z,\xi_W)=(1+10^{-10},1+10^{-10})$、$(2,3)$、$(1/2,4)$（$\xi_A=1$）で、仮想の補正の有限部分はレプトンで表示した全桁、クォークで $1.3\times10^{-12}$ 以内（CKM 行列の倍精度）で一致した。Feynman ゲージを $\xi=1$ で直接置けないのは、縦波の部分の還元が $1/(\xi_W-1)$ を残すからで、$1+10^{-10}$ を80桁で使う。初めは $1+10^{-20}$ にしていたが、ニュートリノの脚の $p^2\to0$ の展開で80桁が尽き、レプトンの仮想の補正が $2$〜$6\%$ ずれた（Mathematica は警告を出さず、$1+10^{-30}$ では Indeterminate になる）。このずれは $\ln m_\ell^2$ に比例し、下の (g) の KLN の確かめで見つかった。

(e) MadLoop。\texttt{wweak\_madloop.py} では、CKM ノートの NLO の UFO の \texttt{[virt=QED]} で、$W^+\to\nu_\tau\tau^+$、$c\bar b$、$u\bar d$ の電弱の仮想の補正を、CKM ノートの MadLoop の検算の点で求めた。光子を含む仮想の補正そのものなので、そのまま比べられる。有限部分は $\mu_R=91$、$500$~GeV で相対 $6\times10^{-12}$ 以内、単極は $1\times10^{-11}$ 以内で一致した（最も悪いのは $u\bar d$ で、$\nu_\tau\tau^+$ と $c\bar b$ では $3\times10^{-13}$ と $1\times10^{-12}$ 以内）。

(f) 実の光子。\texttt{wreal.wl} は、木の $W^+\to f\bar f'\gamma$（Feynman ゲージ。レプトンで3個、クォークで4個の図）の $R=\sum|\mathcal M_R|^2/(e^2\sum|\mathcal M_0|^2)$ を FeynCalc で記号として求める。光子の偏極を $p_3$ に替えた Ward 恒等式は記号として0で、$E\to0$ の $R$ は3個の電荷の eikonal 因子に近づいた（$E=10^{-4}E_{\max}$、$10^{-6}E_{\max}$ で相対 $2.5\times10^{-9}$、$2.5\times10^{-13}$）。硬い部分 $E>\Delta E=\delta E_{\max}$ は手順 (3) と同じく対数の変数の Gauss--Legendre で積分した（図~\ref{fig:dalitz}。質量0で中性のニュートリノの縁では $w$ を線形の変数にする）。\texttt{wreal\_mg5.py} では、MadGraph の木の $W^+\to\nu_ee^+\gamma$、$\nu_\tau\tau^+\gamma$、$u\bar d\gamma$、$c\bar b\gamma$（ユニタリーゲージ）の振幅の2乗と9個の Dalitz の点で比べた。レプトンでは $3\times10^{-9}$ 以内（共線の端から $10^{-6}$ の点。他の点では $4\times10^{-12}$ 以内）で一致した。クォークでは初め最大 $5\times10^{-6}$ ずれた。原因は MadGraph の側にあった。UFO の $d$ 型のクォークの光子の結合 $-0.3333333333333333\,e$ を、MadGraph は Fortran に $-3.333333\mathrm{D}{-}01\,e$ と7桁で書き出す（模型の Fortran の中で厳密でない数はこれだけ）。そのため $Q_d$ が相対 $10^{-7}$ ずれ、$\bar d$ に沿う光子では $|\mathcal M|^2$ が $2\times10^{-7}$ ずれる（振幅が打ち消す点ではもっと大きい）。この数を $1/3$ に直して模型の Fortran を作り直すと、$8\times10^{-10}$ 以内で一致した。

(g) $\delta$ によらないことと KLN。$\delta=10^{-4}$〜$10^{-8}$ で $\delta_{\mathrm{EW}}$ の変化は $O(\delta)$ だけで（$\delta=10^{-4}$ で $6\times10^{-6}$）、手順 (3) と同じく $\delta\to0$ に外挿した。レプトンの質量は共線の対数を切るだけなので、光子について包括的な幅は $m_\ell\to0$ で有限で、$\delta_{\mathrm{EW}}$ は $m_\ell$ によらないはずである（KLN）。$e$ と $\mu$ では、硬い放射の $K_H$（$\Gamma_H/\Gamma_0=\frac\alpha\pi K_H$）は $188.425$ 違い、共線の分裂関数から予想される $(-\ln\delta-\frac34)\ln(m_\mu^2/m_e^2)=188.426$ と合う（$\delta=10^{-8}$）。仮想と軟らかい部分はその逆に変わり、$\delta_{\mathrm{EW}}$ の差は $-1.2\times10^{-7}$ だった（各部分の変化 $0.47$ の $3\times10^{-7}$）。$\tau$ と $\mu$ の差 $1.9\times10^{-5}$ は $m_\tau^2/M_W^2$ の大きさである。軽いクォークの4個のチャネル（$u\bar d$、$u\bar s$、$c\bar d$、$c\bar s$）の $\delta_{\mathrm{EW}}$ も $1.4\times10^{-5}$ 以内でそろった。

(h) $\mu$ への依存。結合を $\mu=M_W/2$、$2M_W$ に走らせると、木の幅の $+1.19\%$、$-1.17\%$ の変化は、$\Gamma_0(1+\delta_{\mathrm{EW}})$ ではレプトンの幅で $-0.016\%$、$-0.012\%$、ハドロンの幅で $-0.016\%$、$-0.011\%$ になった。QCD の因子まで入れると、ハドロンの幅に $+0.40\%$、$-0.36\%$ が残る。手順 (4a) と同じく、QCD の因子の $\alpha_s(\mu)$ の走りである。

(i) 物理。$\delta_{\mathrm{EW}}$ を表~\ref{tab:wweak} に、幅と分岐比を表~\ref{tab:wwidths} に示す（\texttt{wweak\_widths.py}）。$\mu=M_W$ の $\MSbar$ の結合で書いた木の幅に対して、電弱の補正はレプトンで $+0.011\%$、軽いクォークで $-0.038\%$ と小さい。$\mu$ を $M_W/2$、$2M_W$ にすると $\mp1.2\%$ になるので、$g(M_W)$ が補正の大部分をすでに吸収していることになる。$b$ を含む $u\bar b$、$c\bar b$ では $+0.76\%$ で、他のクォークとの差 $+0.79\%$ はトップのループによる $b$ に固有の補正である（手順 (4a) の $Zb\bar b$ と同じ種類）。NLO で $\Gamma_W=2.0859$~GeV、レプトンの分岐比は $10.844\%$（$e$、$\mu$）と $10.836\%$（$\tau$）、ハドロンの分岐比は $67.475\%$ である。測定値（PDG の2026年の版~\cite{pdg26}、参照）とは、普遍性を仮定したレプトンの分岐比 $(10.86\pm0.09)\%$ と $-0.2\sigma$、ハドロンの分岐比と $+0.2\sigma$ で合う。$\Gamma_W$ は $-1.1\sigma$ である（測定値 $2.14\pm0.05$~GeV には尺度因子 $S=1.7$ が付く）。$\tau$ の分岐比の $-2.6\sigma$ は LEP の $\tau$ の過剰として知られたものである（参照）。本ノートの $\Gamma_W$ には、$O(\alpha_s^2)$ の QCD の補正（手順 (4a) と同じ係数で約 $+0.14\%$、参照）と、QCD の因子のクォークの質量の効果が入っていない。後者は、手順 (3) の $Z\to c\bar c$ の $K$ が $\frac34$ から $0.4$〜$3\%$ ずれることから、$\Gamma_W$ の $10^{-4}$ 程度と見込まれる。

\begin{table}[tbp]
\centering\small
\caption{手順 (4b)：$W^+\to f\bar f'$ の電弱の補正 $\delta_{\mathrm{EW}}$（$\%$。\texttt{wreal.wl}。光子を含む仮想の補正と実の光子の放射の和で、$\delta\to0$。$\mu$ の結合と極質量。グルーオンの部分は含まない）。}
\label{tab:wweak}
\begin{tabular}{lrrr|lrrr}
\toprule
$f\bar f'$ & $\mu=M_W/2$ & $M_W$ & $2M_W$ & $f\bar f'$ & $\mu=M_W/2$ & $M_W$ & $2M_W$ \\
\midrule
$\nu_ee^+$ & $-1.1845$ & $+0.0114$ & $+1.1794$ & $u\bar d$ & $-1.2339$ & $-0.0383$ & $+1.1295$ \\
$\nu_\mu\mu^+$ & $-1.1844$ & $+0.0114$ & $+1.1794$ & $u\bar s$ & $-1.2315$ & $-0.0375$ & $+1.1287$ \\
$\nu_\tau\tau^+$ & $-1.1825$ & $+0.0133$ & $+1.1813$ & $u\bar b$ & $+0.8722$ & $+0.7555$ & $+0.6416$ \\
 & & & & $c\bar d$ & $-1.2331$ & $-0.0376$ & $+1.1300$ \\
 & & & & $c\bar s$ & $-1.2368$ & $-0.0389$ & $+1.1312$ \\
 & & & & $c\bar b$ & $+0.8729$ & $+0.7564$ & $+0.6427$ \\
\bottomrule
\end{tabular}
\end{table}

\begin{table}[tbp]
\centering\footnotesize\setlength{\tabcolsep}{4pt}
\caption{手順 (4b)：$W$ の幅と分岐比（\texttt{wweak\_widths.py}。GeV。$\Gamma_{e\nu}=\Gamma(W^+\to e^+\nu)$、$\mathrm{BR}_{e\nu}=\Gamma_{e\nu}/\Gamma_W$ など）。「QCD」は木に $\frac{\alpha_s}\pi$ を掛けた値、NLO はさらに $\delta_{\mathrm{EW}}$ を入れた値。括弧は測定値（PDG の2026年の版~\cite{pdg26}、参照）との差。$\Gamma_{\mu\nu}$ は $\Gamma_{e\nu}$ と相対 $3\times10^{-6}$ まで同じ。右の2列は $\mu=M_W/2$、$2M_W$ にしたときの変化（$\%$）。}
\label{tab:wwidths}
\begin{tabular}{lllll|ll}
\toprule
量 & 木 & QCD & NLO & 測定値（参照） & 木 & NLO \\
\midrule
$\Gamma_{e\nu}$ & 0.226176 & 0.226176 & 0.226202 & — & $+1.19$、$-1.17$ & $-0.016$、$-0.012$ \\
$\Gamma_{\tau\nu}$ & 0.226010 & 0.226010 & 0.226040 & — & $+1.19$、$-1.17$ & $-0.016$、$-0.012$ \\
$\Gamma_{\mathrm{had}}$ & 1.356611 & 1.408006 & 1.407493 & — & $+1.19$、$-1.17$ & $+0.40$、$-0.36$ \\
$\Gamma_W$ & 2.034973（$-4.91\%$） & 2.086368（$-2.51\%$） & 2.085937（$-2.53\%$） & $2.14(5)$ & $+1.19$、$-1.17$ & $+0.27$、$-0.25$ \\
$\mathrm{BR}_{e\nu}$ & 0.111145（$+3.78\%$） & 0.108407（$+1.22\%$） & 0.108441（$+1.25\%$） & $0.1071(16)$ & $0.00$ & $-0.28$、$+0.23$ \\
$\mathrm{BR}_{\mu\nu}$ & 0.111144（$+4.56\%$） & 0.108406（$+1.98\%$） & 0.108441（$+2.01\%$） & $0.1063(15)$ & $0.00$ & $-0.28$、$+0.23$ \\
$\mathrm{BR}_{\tau\nu}$ & 0.111063（$-2.41\%$） & 0.108327（$-4.81\%$） & 0.108364（$-4.78\%$） & $0.1138(21)$ & $0.00$ & $-0.28$、$+0.23$ \\
$\mathrm{BR}_{\mathrm{had}}$ & 0.666648（$-1.11\%$） & 0.674860（$+0.11\%$） & 0.674753（$+0.10\%$） & $0.6741(27)$ & $0.00$ & $+0.14$、$-0.11$ \\
$\Gamma_{\tau\nu}/\Gamma_{e\nu}$ & 0.999267 & 0.999267 & 0.999286 & — & $0.00$ & $0.00$ \\
\bottomrule
\end{tabular}
\end{table}

\FloatBarrier
\paragraph{手順 (4c)：$t\to d_jW^+$ の電弱の補正。}トップの崩壊では、光子は崩壊するトップ（$Q=\frac23$）、$d_j$（$-\frac13$）、$W^+$（$+1$）の三つに付く（図~\ref{fig:tnlo}）。手順 (4b) と同じく、光子を含む電弱の補正 $\delta_{\mathrm{EW}}$ を実の光子の放射と足して求め、$\Gamma=\Gamma_0[1+\delta_{\mathrm{EW}}+\frac{\alpha_sC_F}\pi K]$ とする。$K=-1.855783$ は手順 (3) の $t\to bW^+$ の値で、$t\to sW^+$、$t\to dW^+$（幅の $0.18\%$）にも同じ値を使う。

\begin{figure}[tbp]
\centering
\begin{tikzpicture}[scale=0.85]
% (a) photon from the top
\begin{scope}[shift={(0,0)}]
  \draw[fermion] (-1.15,0) -- (-0.55,0); \draw[fermion] (-0.55,0) -- (0,0);
  \draw[fermion] (0,0) -- (1.15,0.75); \draw[photon] (0,0) -- (1.15,-0.75);
  \draw[photon] (-0.55,0) -- (-0.1,0.95);
  \node[vtx] at (0,0) {}; \node[vtx] at (-0.55,0) {};
  \node[font=\scriptsize] at (-1.0,0.22) {$t$}; \node[font=\scriptsize] at (1.35,0.8) {$b$}; \node[font=\scriptsize] at (1.4,-0.8) {$W^+$};
  \node[font=\scriptsize] at (0.15,1.05) {$\gamma$};
  \node[slab] at (0.1,-1.6) {(a) 実の放射};
\end{scope}
% (b) photon from the W
\begin{scope}[shift={(3.4,0)}]
  \draw[fermion] (-1.15,0) -- (0,0);
  \draw[fermion] (0,0) -- (1.15,0.75); \draw[photon] (0,0) -- (0.5,-0.35); \draw[photon] (0.5,-0.35) -- (1.2,-0.55);
  \draw[photon] (0.5,-0.35) -- (0.95,-1.1);
  \node[vtx] at (0,0) {}; \node[vtx] at (0.5,-0.35) {};
  \node[font=\scriptsize] at (-1.0,0.22) {$t$}; \node[font=\scriptsize] at (1.35,0.8) {$b$}; \node[font=\scriptsize] at (1.5,-0.55) {$W^+$};
  \node[font=\scriptsize] at (1.2,-1.15) {$\gamma$};
  \node[slab] at (0.1,-1.6) {(b) 実の放射};
\end{scope}
% (c) the Goldstone diagram
\begin{scope}[shift={(6.8,0)}]
  \draw[fermion] (-1.15,0) -- (0,0);
  \draw[fermion] (0,0) -- (1.15,0.75); \draw[scalar] (0,0) -- (0.5,-0.35); \draw[photon] (0.5,-0.35) -- (1.2,-0.55);
  \draw[photon] (0.5,-0.35) -- (0.95,-1.1);
  \node[vtx] at (0,0) {}; \node[vtx] at (0.5,-0.35) {};
  \node[font=\scriptsize] at (-1.0,0.22) {$t$}; \node[font=\scriptsize] at (1.35,0.8) {$b$}; \node[font=\scriptsize] at (1.5,-0.55) {$W^+$};
  \node[font=\scriptsize] at (0.05,-0.42) {$G^+$}; \node[font=\scriptsize] at (1.2,-1.15) {$\gamma$};
  \node[slab] at (0.1,-1.6) {(c) 実の放射（$\propto m_t$）};
\end{scope}
% (d) the vertex with the photon between the top and the W
\begin{scope}[shift={(10.2,0)}]
  \draw[fermion] (-1.15,0) -- (-0.55,0); \draw[fermion] (-0.55,0) -- (0,0);
  \draw[fermion] (0,0) -- (1.15,0.75);
  \draw[photon] (0,0) -- (0.6,-0.42); \draw[photon] (0.6,-0.42) -- (1.15,-0.8);
  \draw[photon] (-0.55,0) .. controls (-0.45,-0.75) and (0.2,-0.85) .. (0.6,-0.42);
  \node[vtx] at (0,0) {}; \node[vtx] at (-0.55,0) {}; \node[vtx] at (0.6,-0.42) {};
  \node[font=\scriptsize] at (-1.0,0.22) {$t$}; \node[font=\scriptsize] at (1.35,0.8) {$b$}; \node[font=\scriptsize] at (1.45,-0.85) {$W^+$};
  \node[font=\scriptsize] at (-0.25,0.22) {$t$}; \node[font=\scriptsize] at (-0.25,-0.95) {$\gamma$};
  \node[slab] at (0.1,-1.6) {(d) 頂点の補正};
\end{scope}
% (e) the photon loop of the top field constant
\begin{scope}[shift={(13.6,0)}]
  \draw[fermion] (-1.55,0) -- (-1.1,0); \draw[fermion] (-1.1,0) -- (-0.35,0); \draw[fermion] (-0.35,0) -- (0,0);
  \draw[fermion] (0,0) -- (1.15,0.75); \draw[photon] (0,0) -- (1.15,-0.75);
  \draw[photon] (-1.1,0) arc (180:0:0.375);
  \node[vtx] at (0,0) {}; \node[vtx] at (-1.1,0) {}; \node[vtx] at (-0.35,0) {};
  \node[font=\scriptsize] at (-1.5,-0.22) {$t$}; \node[font=\scriptsize] at (-0.725,0.6) {$\gamma$};
  \node[font=\scriptsize] at (1.35,0.8) {$b$}; \node[font=\scriptsize] at (1.4,-0.8) {$W^+$};
  \node[slab] at (0.0,-1.6) {(e) $\delta Z_t$};
\end{scope}
\end{tikzpicture}
\caption{手順 (4c) の図の例（$t\to bW^+$）。(a)〜(c) と $b$ からの放射の4個が木の実の光子の放射（Feynman ゲージ）で、(c) の $G^+$ の結合は $m_t$（と $m_b$）に比例する。(d) は光子の頂点の補正の一つ（グルーオンを除く 1PI の図は46個）、(e) はトップの場の定数の光子のループ。光子は崩壊するトップにも付くので、軟らかい光子の式は三つの電荷の和になる。}
\label{fig:tnlo}
\end{figure}

(a) 振幅。\texttt{tff\_feyncalc.wl} は、$t\to d_jW^+$ の 1PI の頂点（グルーオンを除く。46個の図）を一般の $R_\xi$ ゲージで求める。形状因子は $\mathcal M=\frac{e}{\sqrt2\hat s}\bar u(p_1)[\slashed\varepsilon^*(F_LP_L+F_RP_R)+(\varepsilon^*\cdot p_1)(H_LP_L+H_RP_R)]u(p_t)$ と書き、木では $F_L=V_{tj}^*$ になった。$F_R$、$H_L$、$H_R$ の紫外の極は記号として0だった。トップの弱い自己エネルギー（13個の図）は、手順 (4a) の \texttt{zff\_feyncalc.wl} で求めた（$\Sigma^{SL}=\Sigma^{SR}$）。

(b) S 行列要素。\texttt{tweak.wl} は、式~\eqref{eq:wren} のエルミート共役の形
\begin{equation}
F_L=V^*_{tj}+F_L^{\mathrm{vtx}}-\tfrac12V^*_{tj}\Sigma^{WW\prime}_T(M_W^2)+\tfrac12\bigl[(\delta Z^{dL})^\dagger V^\dagger+V^\dagger\delta Z^{uL}\bigr]_{jt}+V^*_{tj}\frac{\delta g}g+\delta V^*_{tj}
\label{eq:tren}
\end{equation}
で S 行列要素を作る（場の定数と対抗項は手順 (4b) と同じもの。トップの場の定数は弱い部分と光子の部分の和）。トップは不安定なので、その自己エネルギーには $t\to bW^+$ の吸収部分がある（一般の $R_\xi$ では、$\sqrt{\xi_W}M_W<m_t$ の非物理的な閾値の吸収部分もある）。幅には、$W$ と同じく木との干渉の実部をとる。

(c) 極とゲージ。軟らかい光子は、静止したトップ（$\eta=-1$）と、背中合わせの $d_j$ と $W^+$ の三つの電荷の $\Sigma$ で、手順 (4b) の $A_{ij}$ をそのまま使う（崩壊する粒子の電荷は $e_{d_j}+e_W=\frac23$）。仮想の補正の極との和は、$t\to bW^+$ で $1\times10^{-16}$、$sW^+$ で $3\times10^{-14}$、$dW^+$ で $3\times10^{-12}$ だった。$(\xi_Z,\xi_W)=(1+10^{-10},1+10^{-10})$、$(2,3)$、$(1/2,4)$ の間の差も、$1\times10^{-17}$、$7\times10^{-14}$、$3\times10^{-12}$ だった（$t\to dW^+$ で桁が少ないのは、$|V_{td}|$ が小さく、CKM 行列の倍精度の値の打ち消しが大きいため）。$\xi_W=3$、$4$ では上の非物理的な閾値が開いているが、それを含めてゲージによらなかった。

(d) MadLoop。\texttt{tweak\_madloop.py} で、CKM ノートの NLO の UFO の \texttt{[virt=QED]} の $t\to bW^+$ を CKM ノートの MadLoop の検算の点で求めた。有限部分は $\mu_R=91$、$500$~GeV で相対 $4\times10^{-14}$、単極は相対 $10^{-14}$ で一致した。

(e) 実の光子。\texttt{treal.wl} は、木の $t\to d_jW^+\gamma$（Feynman ゲージで4個の図）の $R=\sum|\mathcal M_R|^2/(e^2\sum|\mathcal M_0|^2)$ を求める。光子の偏極を $p_3$ に替えた振幅の2乗（$W$ の偏極は和をとる）は記号として0で、$E\to0$ では三つの電荷の eikonal 因子に近づいた（$E=10^{-4}E_{\max}$、$10^{-6}E_{\max}$ で相対 $6\times10^{-9}$、$6\times10^{-13}$）。硬い部分の積分は手順 (4b) と同じである。\texttt{treal\_mg5.py} で MadGraph の木の $t\to bW^+\gamma$ と9個の Dalitz の点で比べると、生成したままでは $2\times10^{-7}$ ずれ（手順 (4b) と同じく $b$ の電荷の7桁）、それを $1/3$ に直すと $2\times10^{-10}$ 以内で一致した（$b$ に沿う共線の端から $10^{-6}$ の点。他の点では $6\times10^{-12}$ 以内）。

(f) $\delta$ と $m_b$ によらないこと。$\delta=10^{-4}$〜$10^{-8}$ で $\delta_{\mathrm{EW}}$ の変化は $O(\delta)$（$t\to bW^+$ で $\delta=10^{-4}$ のとき $5\times10^{-7}$）だった。$b$ の質量は共線の対数を切るだけなので、$\delta_{\mathrm{EW}}$ は $m_b\to0$ で $m_b$ によらない値に近づくはずである（KLN）。$m_b$ を $1/3$、$1/10$ にすると、硬い放射の $K_H$ は $4.345$、$9.069$ 増え（$\frac\alpha\pi K_H$ で $0.011$、$0.023$）、共線の係数 $Q_b^2(-\ln\delta'-\frac34)\ln f^2=4.317$、$9.044$（$\delta'=\Delta E/E_{b\gamma}$、$E_{b\gamma}$ は $b$ と光子の対のエネルギー）と $0.6\%$、$0.3\%$ で合う（差は物理的な $m_b$ での $K_H$ の $m_b^2/m_t^2$ の補正）。仮想と軟らかい部分はその逆に変わり、$\delta_{\mathrm{EW}}$ の変化は $-3.6\times10^{-5}$、$-4.1\times10^{-5}$ にとどまった。$m_b/3$ から $m_b/10$ の間の変化は $4\times10^{-6}$ で、$m_b^2/m_t^2$ のべきで頭打ちになっている。軽い $d$、$s$ の二つのチャネルの $\delta_{\mathrm{EW}}$ も $1.2\times10^{-5}$ 以内でそろった。

(g) $\mu$ への依存。結合を $\mu=m_t/2$、$2m_t$ に走らせると、木の幅の $+1.18\%$、$-1.15\%$ の変化は、$\Gamma_0(1+\delta_{\mathrm{EW}})$ では $+0.06\%$、$-0.08\%$ になった。QCD の因子まで入れると $-0.87\%$、$+0.70\%$ が残る。これは QCD の因子の $\alpha_s(\mu)$ の走りで、$O(\alpha_s^2)$ の補正が打ち消すはずの分である。

(h) 物理。$\delta_{\mathrm{EW}}$ と幅を表~\ref{tab:tweak} に示す（\texttt{tweak\_widths.py}）。$\mu=m_t$ で、電弱の補正は $t\to bW^+$ で $+3.28\%$、$sW^+$、$dW^+$ で $+2.66\%$ である。$b$ に固有の $+0.62\%$ は、手順 (4b) の $u\bar b$、$c\bar b$ の $+0.79\%$ と同じ種類の補正である。$+3.28\%$ のうち $1.73\%$ は木の結合の処方による。$G_F$ で書いた木の結合 $g^2=4\sqrt2G_FM_W^2$ は、$\MSbar$ の $g(m_t)^2$ より $1.73\%$ 大きいので、$G_F$ の木に対する電弱の補正は $+1.52\%$ になる。これは文献~\cite{ds91} の $G_F$ の処方の電弱の補正と同じ $1$〜$2\%$ の大きさである（参照）。QCD の $-8.51\%$（手順 (3)）と合わせて、NLO で $\Gamma_t=1.3776$~GeV、$\mathrm{BR}(t\to bW^+)=99.823\%$ になり、測定値 $1.42^{+0.19}_{-0.15}$~GeV（PDG の2026年の版~\cite{pdg26t}、参照）と $-0.3\sigma$ で合う。$O(\alpha_s^2)$ の QCD の補正（文献~\cite{cm98}、参照）と、$O(\alpha\alpha_s)$ の補正は入っていない。

\begin{table}[tbp]
\centering\footnotesize\setlength{\tabcolsep}{4pt}
\caption{手順 (4c)：$t\to d_jW^+$ の電弱の補正 $\delta_{\mathrm{EW}}$（$\%$。\texttt{treal.wl}。光子を含む仮想の補正と実の光子の放射の和で、$\delta\to0$）と、トップの幅（\texttt{tweak\_widths.py}。GeV）。「QCD」は手順 (3) の値。測定値は PDG の2026年の版~\cite{pdg26t}（参照）。右の2列は $\mu=m_t/2$、$2m_t$ にしたときの変化（$\%$）。}
\label{tab:tweak}
\begin{tabular}{lrrr}
\toprule
$\delta_{\mathrm{EW}}$ & $\mu=m_t/2$ & $m_t$ & $2m_t$ \\
\midrule
$t\to dW^+$ & $+2.7953$ & $+2.6570$ & $+2.5218$ \\
$t\to sW^+$ & $+2.7945$ & $+2.6582$ & $+2.5251$ \\
$t\to bW^+$ & $+2.1370$ & $+3.2811$ & $+4.3990$ \\
$t\to bW^+$（$m_b/3$） & & $+3.2775$ & \\
$t\to bW^+$（$m_b/10$） & & $+3.2771$ & \\
\bottomrule
\end{tabular}

\medskip
\begin{tabular}{lllll|ll}
\toprule
量 & 木 & QCD & NLO & 測定値（参照） & 木 & NLO \\
\midrule
$\Gamma(t\to bW^+)$ & 1.451064 & 1.327554 & 1.375165 & — & $+1.18$、$-1.15$ & $-0.87$、$+0.70$ \\
$\Gamma_t$ & 1.453649（$+2.4\%$） & 1.329919（$-6.3\%$） & 1.377599（$-3.0\%$） & $1.42^{+0.19}_{-0.15}$ & $+1.18$、$-1.15$ & $-0.87$、$+0.70$ \\
$\mathrm{BR}(t\to bW^+)$ & 0.998222 & 0.998222 & 0.998234 & — & $0.00$ & $0.00$ \\
\bottomrule
\end{tabular}
\end{table}

\FloatBarrier
\paragraph{手順 (5)：$Z$ の有効結合と $\sin^2\theta_{\mathrm{eff}}$。}式~\eqref{eq:effcoup} の有効結合は、手順 (4a) の繰り込んだ弱い形状因子（式~\eqref{eq:zren}）から $\mathcal V_f=F_L+F_R$、$\mathcal A_f=F_L-F_R$（$g_Z/2$ の単位。木では $T_3-2Q\hat s^2$、$T_3$）で与えられる。$Z$ の留数 $\frac12\delta Z_{ZZ}$ は $F_L$、$F_R$ に共通なので比 $\mathcal V_f/\mathcal A_f$ からは消える。光子とグルーオンの終状態の補正は、質量0のフェルミオンでは $\mathcal V_f$ と $\mathcal A_f$ に同じ因子を掛けるだけなので含めない。$H_L$、$H_R$ は $O(m_f)$ である。これは LEP の有効結合の定義と同じである。\texttt{zeff.wl} は比を $O(\alpha)$（1ループの部分について線形）でとり、式~\eqref{eq:s2eff} の $\sin^2\theta^f_{\mathrm{eff}}$ を求める。比をそのままとった値との差は $-5\times10^{-5}$ で、$O(\alpha^2)$ の大きさの目安になる。\texttt{zpole.py} は $x_f=1-4|Q_f|\sin^2\theta^f_{\mathrm{eff}}$ から $A_f=2x_f/(1+x_f^2)$、$A^{0,f}_{\mathrm{FB}}=\frac34A_eA_f$ を求める。

(a) ゲージと $\mu$ によらないこと。$(\xi_Z,\xi_W)=(1,1)$、$(2,3)$、$(1/2,4)$ で、8個の荷電フェルミオンの $\sin^2\theta^f_{\mathrm{eff}}$ は $10^{-18}$ 以内で一致した。結合を $\mu=M_Z/2$、$M_Z$、$2M_Z$ に走らせると、$\hat s^2(\mu)=0.228958$、$0.232420$、$0.235903$ は $\pm1.5\%$ 変わる。一方 $\sin^2\theta^{\mathrm{lept}}_{\mathrm{eff}}=0.232448$、$0.232423$、$0.232376$ の変化は $\pm3\times10^{-5}$ にとどまり、$\hat\kappa_f(\mu)$ が $\hat s^2(\mu)$ の走りを打ち消す。

(b) 内訳。$\mu=M_Z$ で $\sin^2\theta_{\mathrm{eff}}-\hat s^2(M_Z)$ を部分に分けると、レプトンでは $Z$--光子の混合（図~\ref{fig:znlo}(c)）が $-5.22\times10^{-4}$、頂点が $+5.69\times10^{-4}$、フェルミオンの脚が $-0.43\times10^{-4}$ で、和は $+0.04\times10^{-4}$（$\hat\kappa_\ell=1.000016$）だった。$Z$ の脚と $\MSbar$ の対抗項の寄与は $10^{-14}$ で、上の通り消える。$b$ では頂点が $+26.6\times10^{-4}$、脚が $-7.2\times10^{-4}$ で、$\hat\kappa_b=1.0061$ になる。手順 (4a) の $Zb\bar b$ と同じく、トップのループによる $b$ に固有の補正である。値を表~\ref{tab:zeff} に示す。

(c) 測定値との比較。本ノートの入力（SM ノートの手順 (7) の $M_W$ の処方）では $\sin^2\theta^{\mathrm{lept}}_{\mathrm{eff}}=0.232423$ で、LEP と SLD の値 $0.23153\pm0.00016$~\cite{lepewwg06}（参照）より $8.9\times10^{-4}$（$+5.6\sigma$）大きい。同じ1ループの計算に、SM ノートの $\alpha$ の処方（$M_W$ の代わりに $\alpha(0)$ を入力にし、$M_W$ は予言になる）の結合を入れると、$\hat s^2(M_Z)=0.231658$、$\sin^2\theta^{\mathrm{lept}}_{\mathrm{eff}}=0.231670$（$+0.9\sigma$）になる（表~\ref{tab:zasym}）。二つの処方は1ループでは同等で、差の $7.5\times10^{-4}$ は1ループで決めた結合に残る高次の項の大きさである。これは測定の誤差の約5倍で、$\sin^2\theta_{\mathrm{eff}}$ の精度は本ノートの NLO の範囲を超えている。$M_W$ の処方では $\hat s^2$ が $M_W/M_Z$ に $\Delta\rho$ を通じて結びつき、$\sin^2\theta_{\mathrm{eff}}$ は $c_W^2\Delta\rho$ だけずれる。$\Delta\rho$ の $O(\alpha\alpha_s)$ の補正 $-\frac23(1+\frac{\pi^2}3)\frac{\alpha_s}\pi\Delta\rho$（文献~\cite{dv87} の係数、参照）を式~\eqref{eq:drho} の $\Delta\rho$（極質量の $m_t$ で $0.00933$）に当てると、$M_W$ の処方で $-7.1\times10^{-4}$、$\alpha$ の処方で $+2.9\times10^{-4}$（$-\hat s^2\hat c^2/(\hat c^2-\hat s^2)\cdot\delta\Delta\rho$）になる。処方の差と同じ大きさである。非対称度のうち、$A^{0,b}_{\mathrm{FB}}$ は $M_W$ の処方で $-0.7\sigma$、SLD の $A_\ell$ は $\alpha$ の処方の方が近い（$-2.6\sigma$）。この二つの測定が $\sin^2\theta_{\mathrm{eff}}$ について約 $3\sigma$ 食い違うことはよく知られている~\cite{lepewwg06}（参照）。$A_b$、$A_c$ は誤差が大きく、どちらの処方でも $0.6\sigma$ 以内で合う。

\begin{table}[tbp]
\centering\small
\caption{手順 (5)：$\mu=M_Z$ の $\sin^2\theta^f_{\mathrm{eff}}$ と $\hat\kappa_f=\sin^2\theta^f_{\mathrm{eff}}/\hat s^2(M_Z)$（\texttt{zeff.wl}。$\hat s^2(M_Z)=0.2324197$。$O(\alpha)$）。右は SM ノートの $\alpha$ の処方の結合での値（$\hat s^2(M_Z)=0.2316581$）。}
\label{tab:zeff}
\begin{tabular}{lll|l}
\toprule
$f$ & $\sin^2\theta^f_{\mathrm{eff}}$ & $\hat\kappa_f$ & $\alpha$ の処方 \\
\midrule
$e$、$\mu$ & 0.2324233 & 1.000016 & 0.2316695 \\
$\tau$ & 0.2324219 & 1.000010 & 0.2316682 \\
$u$、$c$ & 0.2323168、0.2323153 & 0.999557、0.999551 & 0.2315636、0.2315620 \\
$d$、$s$ & 0.2321895、0.2321926 & 0.999010、0.999023 & 0.2314369、0.2314400 \\
$b$ & 0.2338331 & 1.006081 & 0.2330745 \\
\bottomrule
\end{tabular}
\end{table}

\begin{table}[tbp]
\centering\small
\caption{手順 (5)：$Z$ 極の非対称度（\texttt{zpole.py}。$A_f=2x_f/(1+x_f^2)$、$A^{0,f}_{\mathrm{FB}}=\frac34A_eA_f$）。木は $\hat s^2(M_Z)$ で、NLO は本ノートの入力（$M_W$ の処方）と $\alpha$ の処方の値。測定値は LEP と SLD~\cite{lepewwg06}（参照。$A_\ell$ は SLD の値）で、括弧は測定値との差（$\sigma$）。}
\label{tab:zasym}
\begin{tabular}{llllll}
\toprule
量 & 木 & NLO（$M_W$） & NLO（$\alpha$） & 測定値（参照） & $\mu=M_Z/2$、$2M_Z$（$M_W$） \\
\midrule
$\sin^2\theta^{\mathrm{lept}}_{\mathrm{eff}}$ & 0.232420 & 0.232423（$+5.6$） & 0.231670（$+0.9$） & $0.23153(16)$ & 0.232448、0.232376 \\
$A_\ell$ & 0.13995 & 0.13992（$-5.4$） & 0.14586（$-2.6$） & $0.1513(21)$ & 0.13973、0.14029 \\
$A_b$ & 0.93495 & 0.93404（$+0.6$） & 0.93453（$+0.6$） & $0.923(20)$ & 0.93403、0.93406 \\
$A_c$ & 0.66438 & 0.66475（$-0.2$） & 0.66736（$-0.1$） & $0.670(27)$ & 0.66466、0.66492 \\
$A^{0,\ell}_{\mathrm{FB}}$ & 0.01469 & 0.01468（$-2.4$） & 0.01596（$-1.1$） & $0.0171(10)$ & 0.01464、0.01476 \\
$A^{0,b}_{\mathrm{FB}}$ & 0.09813 & 0.09802（$-0.7$） & 0.10223（$+1.9$） & $0.0992(16)$ & 0.09788、0.09828 \\
$A^{0,c}_{\mathrm{FB}}$ & 0.06974 & 0.06976（$-0.3$） & 0.07301（$+0.7$） & $0.0707(35)$ & 0.06965、0.06996 \\
\bottomrule
\end{tabular}
\end{table}

\FloatBarrier
\paragraph{手順 (6)：$\Delta\rho$ と $M_W$ の予言。}\texttt{mw.wl} は、SM ノートの手順 (3) の $\Sigma^{WW}_T$、$\Sigma^{ZZ}_T$（一般の $R_\xi$、すべてのフェルミオンの質量、1PI の部分とタドポールの部分）を使う（図~\ref{fig:drho}）。

\begin{figure}[tbp]
\centering
\begin{tikzpicture}[scale=0.85]
% (a) the top-bottom loop of the W
\begin{scope}[shift={(0,0)}]
  \draw[photon] (-1.5,0) -- (-0.6,0); \draw[photon] (0.6,0) -- (1.5,0);
  \draw[fermion] (-0.6,0) arc (180:0:0.6); \draw[fermion] (0.6,0) arc (0:-180:0.6);
  \node[vtx] at (-0.6,0) {}; \node[vtx] at (0.6,0) {};
  \node[font=\scriptsize] at (-1.2,0.25) {$W^+$}; \node[font=\scriptsize] at (1.2,0.25) {$W^+$};
  \node[font=\scriptsize] at (0,0.82) {$t$}; \node[font=\scriptsize] at (0,-0.82) {$b$};
  \node[slab] at (0,-1.35) {(a) $\Sigma^{WW}_T$ の $t$、$b$};
\end{scope}
% (b) the top loop of the Z
\begin{scope}[shift={(3.7,0)}]
  \draw[photon] (-1.5,0) -- (-0.6,0); \draw[photon] (0.6,0) -- (1.5,0);
  \draw[fermion] (-0.6,0) arc (180:0:0.6); \draw[fermion] (0.6,0) arc (0:-180:0.6);
  \node[vtx] at (-0.6,0) {}; \node[vtx] at (0.6,0) {};
  \node[font=\scriptsize] at (-1.2,0.25) {$Z$}; \node[font=\scriptsize] at (1.2,0.25) {$Z$};
  \node[font=\scriptsize] at (0,0.82) {$t$}; \node[font=\scriptsize] at (0,-0.82) {$t$};
  \node[slab] at (0,-1.35) {(b) $\Sigma^{ZZ}_T$ の $t$（$b$ も）};
\end{scope}
% (c) the tadpole
\begin{scope}[shift={(7.4,0)}]
  \draw[photon] (-1.3,0) -- (0,0); \draw[photon] (0,0) -- (1.3,0);
  \draw[scalar] (0,0) -- (0,0.62);
  \draw[fermion] (0,0.62) arc (-90:270:0.33);
  \node[vtx] at (0,0) {}; \node[vtx] at (0,0.62) {};
  \node[font=\scriptsize] at (-1.0,0.25) {$W$, $Z$}; \node[font=\scriptsize] at (1.0,0.25) {$W$, $Z$};
  \node[font=\scriptsize] at (0.22,0.3) {$H$}; \node[font=\scriptsize] at (0.52,1.12) {$t$};
  \node[slab] at (0,-1.35) {(c) タドポール（FJ）};
\end{scope}
% (d) a bosonic loop
\begin{scope}[shift={(11.1,0)}]
  \draw[photon] (-1.5,0) -- (-0.6,0); \draw[photon] (0.6,0) -- (1.5,0);
  \draw[photon] (-0.6,0) arc (180:0:0.6); \draw[photon] (0.6,0) arc (0:-180:0.6);
  \node[vtx] at (-0.6,0) {}; \node[vtx] at (0.6,0) {};
  \node[font=\scriptsize] at (-1.2,0.25) {$W^+$}; \node[font=\scriptsize] at (1.2,0.25) {$W^+$};
  \node[font=\scriptsize] at (0,0.85) {$W$}; \node[font=\scriptsize] at (0,-0.85) {$Z$, $\gamma$};
  \node[slab] at (0,-1.35) {(d) ボソンのループ};
\end{scope}
\end{tikzpicture}
\caption{手順 (6) の $\Delta\rho=\Sigma^{ZZ}_T(0)/M_Z^2-\Sigma^{WW}_T(0)/M_W^2$ に入る図の例。(a)、(b) のトップとボトムの二重項は紫外で有限、ゲージによらず、$m_t$ と $m_b$ の差に比例する custodial 対称性の破れとして式~\eqref{eq:drho}（$m_b$ を含めた $F_0$）になる。(c) の FJ のタドポールは、$\Sigma^{WW}_T/M_W^2$ と $\Sigma^{ZZ}_T/M_Z^2$ に同じ係数で入り、$\Delta\rho$ と $M_W$ の予言の $dZ-dW$ から消える。(d) のボソンのループ（ほかに $H$、Goldstone、ゴースト）は $g'$ が custodial 対称性を破るので紫外の極を残し、ゲージによる。}
\label{fig:drho}
\end{figure}


(a) タドポール。タドポールの部分は $k^2$ によらず、$\Sigma^{WW}_{\mathrm{tad}}/M_W^2-\Sigma^{ZZ}_{\mathrm{tad}}/M_Z^2$ は記号として0だった。したがってタドポールは、$\Delta\rho$ からも、下の $M_W$ の予言に入る $dZ-dW$（$dV=\re\hat\Sigma^{VV}_T(M_V^2)/M_V^2$）からも消え、1PI の部分だけで足りる。

(b) $\Delta\rho$。トップとボトムの二重項の部分（ループ関数に $m_t$、$m_b$ だけを含む項）の $\Delta\rho$ は、紫外の極もゲージの母数も含まず、$\mu=50$、$300$~GeV で同じ値だった。式 $N_cg^2F_0(m_t^2,m_b^2)/(4M_W^2K)$、$F_0(a,b)=a+b-\frac{2ab}{a-b}\ln\frac ab$（$m_b\to0$ で式~\eqref{eq:drho}）とは、$(m_t,m_b)=(172.57,4.78)$、$(172.57,0.001)$、$(100,50)$~GeV で相対 $10^{-16}$ で一致した（$k^2\to0$ は $k^2=10^{-12}$~GeV$^2$ で60桁）。$G_F$ と極質量では $\Delta\rho_{tb}=0.009237$（$m_b=0$ で $0.009332$）である。一方、すべてのループを含む $\Delta\rho$ は、$\mu=M_Z$ で紫外の極が $(\xi_Z,\xi_W)=(2,2)$、$(3,3)$、$(1/2,4)$ で $0.0027$、$0.0030$、$-0.0004$、有限部分が $0.0086$、$0.0077$、$0.0112$ と、ゲージによって変わった。ボソンのループでは $g'$ が custodial 対称性を破るためで、$\Delta\rho$ そのものは物理量ではない。物理量は次の $M_W$ である。

(c) $M_W$ の予言。入力は SM ノートの $\alpha$ の処方（$\alpha(0)$、$\Delta\alpha$、$G_F$、$M_Z$、$M_H$、$m_t$、$\alpha_s$。$M_W$ は予言）の結合である。SM ノートの当てはめでは、$dW$、$dZ$、$\hat\Delta$ のタドポールの部分（$\mu=M_t$ でそれぞれ約 $0.12$）は $\hat\Delta-dZ$ と $dZ-dW$ の組でしか消えないので、関係は
\begin{equation}
M_W=M_Z\hat c(\mu)\Bigl[1+\tfrac12\bigl(dZ-dW\bigr)\Bigr]
\label{eq:mwpred}
\end{equation}
と線形に書く。\texttt{mw.wl} は、\texttt{nlo\_inputs.py} が1ループの方程式で走らせた $\hat c^2(\mu)$ と、Package-X で求めた $dV$ でこれを求める（ループは SM ノートの当てはめと同じ木の点で、$M_W$ は不動点まで反復する。$M_Z$ は SM ノートの極質量 $91.15387$~GeV）。$\mu=M_t$ で $M_W=80.36795$~GeV になり、三つのゲージで $10^{-36}$ まで一致した。$\mu=M_Z/2$、$M_Z$、$M_t$、$2M_t$ では $80.3883$、$80.3786$、$80.3680$、$80.3546$~GeV で、変化の幅は $34$~MeV である。木の $M_Z\hat c(\mu)$ では同じ範囲で $528$~MeV 変わる（表~\ref{tab:mw}）。

SM ノートの値 $80.38171$~GeV とは $13.8$~MeV 違う。SM ノートの不動点は、$e(\mu)$、$G_F$、$M_Z$ から作った $\hat c^2$ を使う。一方、本ノートの入力の $g$、$g'$ は、その $M_W$ で $M_W$ の処方の式から作り直したもので、その $\hat c^2$ は相対 $3.4\times10^{-4}$（$O(\alpha^2)$）違う。式~\eqref{eq:mwpred} と1ループで同等な $M_Z\hat c\sqrt{1+d}$、$M_Z\hat c/\sqrt{1-d}$（$d=dZ-dW=0.0159$）も $80.3655$、$80.3756$~GeV になる。1ループの $M_W$ には、$d^2\approx2.5\times10^{-4}$ の大きさの、$10$〜$20$~MeV の切り捨ての不定性があり、$\mu$ を変えた幅と同じ程度である。

(d) 測定値との比較。測定の $M_W$ は走る幅の Breit--Wigner の質量なので、$M_{\mathrm{BW}}=M\sqrt{1+\Gamma_W^2/M^2}$（手順 (4b) の $\Gamma_W=2.0859$~GeV）で $27.1$~MeV 大きくして $80.3950$~GeV と比べる。本ノートの入力の PDG 2024 の $80.3692\pm0.0133$~GeV~\cite{pdg24} とは $+25.8$~MeV（$+1.9\sigma$）、PDG の2026年の版の $80.3625\pm0.0077$~GeV~\cite{pdg26} とは $+32.5$~MeV（$+4.2\sigma$）違う（参照）。これは1ループの予言の不定性の中にある。$\Delta\rho$ の $O(\alpha\alpha_s)$ の補正（手順 (5) と同じく文献~\cite{dv87} の係数）は $\delta\Delta\rho=-0.000909$ で、$\alpha$、$G_F$、$M_Z$ を固定すると $\delta M_W=\frac{M_W}2\frac{c^2}{c^2-s^2}\delta\Delta\rho=-51$~MeV になる。これだけで上の差（$+26$〜$+33$~MeV）を超える大きさである。1ループの $M_W$ の予言の不定性は数十 MeV で、測定の誤差より大きい。$\sin^2\theta_{\mathrm{eff}}$（手順 (5)）と同じく、測定の精度に届くには2ループの補正が要る。

\begin{table}[tbp]
\centering\small
\caption{手順 (6)：$M_W$ の予言（\texttt{mw.wl}、\texttt{mw\_summary.py}。GeV）。SM ノートの $\alpha$ の処方の結合、式~\eqref{eq:mwpred}、極質量（BW は手順 (4b) の $\Gamma_W$ で直した Breit--Wigner の質量）。測定値は参照。}
\label{tab:mw}
\begin{tabular}{lll}
\toprule
 & 木 $M_Z\hat c(\mu)$ & NLO の極質量 \\
\midrule
$\mu=M_Z/2$ & 80.0804 & 80.3883 \\
$\mu=M_Z$ & 79.9011 & 80.3786 \\
$\mu=M_t$ & 79.7347 & 80.3680（BW 80.3950） \\
$\mu=2M_t$ & 79.5525 & 80.3546 \\
\midrule
$\mu=M_t$、$M_Z\hat c\sqrt{1+d}$、$M_Z\hat c/\sqrt{1-d}$ & & 80.3655、80.3756 \\
SM ノートの手順 (7)（その当てはめの $\hat c^2$） & & 80.3817 \\
\midrule
測定値（BW）：PDG 2024、PDG 2026 & \multicolumn{2}{l}{$80.3692\pm0.0133$、$80.3625\pm0.0077$} \\
\bottomrule
\end{tabular}
\end{table}

\FloatBarrier
\paragraph{手順 (7)：$\MSbar$ と on-shell の処方。}SM ノートの $M_W$ の処方の換算は、1ループで
\begin{equation}
g(\mu)^2=g_0^2(1+dW-\hat\Delta),\quad g(\mu)^2+g'(\mu)^2=G_0^2(1+dZ-\hat\Delta),\quad \hat c^2(\mu)=c_0^2(1+dW-dZ)
\label{eq:convOS}
\end{equation}
である。$g_0^2=4\sqrt2G_FM_W^2$、$G_0^2=4\sqrt2G_FM_Z^2$、$c_0^2=M_W^2/M_Z^2$ は on-shell の $G_\mu$ の処方の木の結合で、$\hat\Delta$ は SM ノートの $\mu$ 崩壊の補正である（図~\ref{fig:gfdecay}）。これを使うと、$G_\mu$ の処方の1ループの幅は $\Gamma=\Gamma_0(g_0,s_0^2)[1+\delta(\mu)+C(\mu)]$ になる。$\delta(\mu)$ は手順 (4a)〜(4c) の値で、$C$ は木の結合を $(g_0,s_0^2)$ から $(g(\mu),\hat s^2(\mu))$ に替えたときの木の1ループの変化である。$Z\to\nu\bar\nu$ では $C=dZ-\hat\Delta$、$Z\to e^+e^-$ ではさらに $\hat s^2$ の変化 $-c_0^2(dW-dZ)$ の分が付き、$W\to e\nu$、$t\to bW^+$ では $C=dW-\hat\Delta$ である（$V_{tb}$ は両方で同じ値にし、CKM ノートの手順 (6) の換算は含めない）。$Z\to e^+e^-$ は弱い補正だけで比べる（光子の因子は共通）。\texttt{scheme.wl} は $dV-\hat\Delta$ を SM ノートの当てはめの木の点（その極質量 $M_W=80.34217$、$M_Z=91.15387$~GeV）で求め、\texttt{scheme\_compare.py} が幅を比べる。

\begin{figure}[tbp]
\centering
\begin{tikzpicture}[scale=0.85]
% (a) the tree
\begin{scope}[shift={(0,0)}]
  \draw[fermion] (-1.1,0.85) -- (0,0.5); \draw[fermion] (0,0.5) -- (1.1,0.85);
  \draw[photon] (0,0.5) -- (0,-0.5);
  \draw[fermion] (0,-0.5) -- (1.1,-0.85); \draw[fermion] (-1.1,-0.85) -- (0,-0.5);
  \node[vtx] at (0,0.5) {}; \node[vtx] at (0,-0.5) {};
  \node[font=\scriptsize] at (-1.3,0.98) {$\mu^-$}; \node[font=\scriptsize] at (1.3,0.98) {$\nu_\mu$};
  \node[font=\scriptsize] at (1.3,-0.98) {$e^-$}; \node[font=\scriptsize] at (-1.3,-0.98) {$\bar\nu_e$};
  \node[font=\scriptsize] at (0.28,0) {$W$};
  \node[slab] at (0,-1.5) {(a) 木};
\end{scope}
% (b) the W self-energy at zero momentum
\begin{scope}[shift={(3.7,0)}]
  \draw[fermion] (-1.1,0.85) -- (0,0.5); \draw[fermion] (0,0.5) -- (1.1,0.85);
  \draw[photon] (0,0.5) -- (0,0.18); \draw[photon] (0,-0.18) -- (0,-0.5);
  \node[blob,minimum size=10pt] at (0,0) {};
  \draw[fermion] (0,-0.5) -- (1.1,-0.85); \draw[fermion] (-1.1,-0.85) -- (0,-0.5);
  \node[vtx] at (0,0.5) {}; \node[vtx] at (0,-0.5) {};
  \node[font=\scriptsize] at (-1.3,0.98) {$\mu^-$}; \node[font=\scriptsize] at (1.3,0.98) {$\nu_\mu$};
  \node[font=\scriptsize] at (1.3,-0.98) {$e^-$}; \node[font=\scriptsize] at (-1.3,-0.98) {$\bar\nu_e$};
  \node[font=\scriptsize,anchor=west] at (0.2,0) {$\Sigma^{WW}_T(0)$};
  \node[slab] at (0,-1.5) {(b) $W$ の自己エネルギー};
\end{scope}
% (c) a vertex correction
\begin{scope}[shift={(7.4,0)}]
  \draw[fermion] (-1.1,0.85) -- (-0.55,0.675); \draw[fermion] (-0.55,0.675) -- (0,0.5); \draw[fermion] (0,0.5) -- (0.55,0.675); \draw[fermion] (0.55,0.675) -- (1.1,0.85);
  \draw[photon] (-0.55,0.675) .. controls (-0.35,1.3) and (0.35,1.3) .. (0.55,0.675);
  \draw[photon] (0,0.5) -- (0,-0.5);
  \draw[fermion] (0,-0.5) -- (1.1,-0.85); \draw[fermion] (-1.1,-0.85) -- (0,-0.5);
  \node[vtx] at (0,0.5) {}; \node[vtx] at (0,-0.5) {}; \node[vtx] at (-0.55,0.675) {}; \node[vtx] at (0.55,0.675) {};
  \node[font=\scriptsize] at (-1.3,0.98) {$\mu^-$}; \node[font=\scriptsize] at (1.3,0.98) {$\nu_\mu$};
  \node[font=\scriptsize] at (1.3,-0.98) {$e^-$}; \node[font=\scriptsize] at (-1.3,-0.98) {$\bar\nu_e$};
  \node[font=\scriptsize] at (0,1.38) {$Z$}; \node[font=\scriptsize] at (0.28,0) {$W$};
  \node[slab] at (0,-1.5) {(c) 頂点 $\delta_V$};
\end{scope}
% (d) a box
\begin{scope}[shift={(11.1,0)}]
  \draw[fermion] (-1.1,0.85) -- (-0.45,0.5); \draw[fermion] (0.45,0.5) -- (1.1,0.85);
  \draw[fermion] (-0.45,0.5) -- (0.45,0.5);
  \draw[photon] (-0.45,0.5) -- (-0.45,-0.5); \draw[photon] (0.45,0.5) -- (0.45,-0.5);
  \draw[fermion] (0.45,-0.5) -- (1.1,-0.85); \draw[fermion] (-0.45,-0.5) -- (0.45,-0.5); \draw[fermion] (-1.1,-0.85) -- (-0.45,-0.5);
  \node[vtx] at (-0.45,0.5) {}; \node[vtx] at (0.45,0.5) {}; \node[vtx] at (-0.45,-0.5) {}; \node[vtx] at (0.45,-0.5) {};
  \node[font=\scriptsize] at (-1.3,0.98) {$\mu^-$}; \node[font=\scriptsize] at (1.3,0.98) {$\nu_\mu$};
  \node[font=\scriptsize] at (1.3,-0.98) {$e^-$}; \node[font=\scriptsize] at (-1.3,-0.98) {$\bar\nu_e$};
  \node[font=\scriptsize] at (-0.72,0) {$W$}; \node[font=\scriptsize] at (0.72,0) {$Z$};
  \node[font=\scriptsize] at (0,0.7) {$\nu_\mu$}; \node[font=\scriptsize] at (0,-0.72) {$e$};
  \node[slab] at (0,-1.5) {(d) 箱 $\delta_B$};
\end{scope}
\end{tikzpicture}
\caption{手順 (7) の換算 \eqref{eq:convOS} に入る $\mu$ 崩壊の補正 $\hat\Delta$（SM ノートの手順 (6)。外線の運動量と $e$、$\mu$ の質量を0にする）の図の例。(a) の木で $G_F/\sqrt2=1/(2v^2)$ になり、1ループで $G_F/\sqrt2=(1+\hat\Delta)/(2v^2)$、$\hat\Delta=[\Sigma^{WW}_T(0)/M_W^2+\delta_V+\delta_B-\frac12\sum_f\Sigma^L_f(0)]_{\MSbar}$ になる。最後の項は4本の外線の場の定数である。(b) の $\Sigma^{WW}_T(0)$ はタドポールの付いた図を含み（その部分が $\hat\Delta$ で最も大きい）、それが $dW$、$dZ$ のタドポールと打ち消す。比べるのは、標準模型の振幅から $G_F$ を定める Fermi 理論の光子のループを引いたもので、後者は運動量0では0になる。}
\label{fig:gfdecay}
\end{figure}

(a) ゲージ。$dW$、$dZ$（タドポールを含む）と $\hat\Delta$ はそれぞれゲージによらないはずの量で（極質量、$\MSbar$ の結合、$G_F$ に対応する）、$\mu=M_Z$ の $dW-\hat\Delta$、$dZ-\hat\Delta$ は三つのゲージで22桁一致した。タドポールの部分（約 $0.12$）は式~\eqref{eq:convOS} の組で消える。

(b) $\mu$ によらないこと。$G_\mu$ の処方の木 $\Gamma_0(g_0,s_0^2)$ は $\mu$ を含まないので、$\delta(\mu)+C(\mu)$ は $O(\alpha)$ で $\mu$ によらないはずである。$\mu=M/2$〜$2M$ の変化は、$Z\to\nu\bar\nu$ で $0.056\%$、$W\to e\nu$ で $0.016\%$、$t\to bW^+$ で $0.11\%$、$Z\to e^+e^-$ で $0.23\%$ だった。$\MSbar$ の木の幅では $1.5$〜$2.4\%$ 変わるので、残りは $O(\alpha^2)$ である（表~\ref{tab:scheme}）。

(c) 二つの処方の差。$\mu=M$ で $\MSbar$ と $G_\mu$ の NLO の幅の比から1を引くと、$Z\to\nu\bar\nu$ で $-1.1\times10^{-4}$、$W\to e\nu$ で $-1.9\times10^{-4}$、$t\to bW^+$ で $-4.6\times10^{-4}$、$Z\to e^+e^-$ で $+1.2\times10^{-3}$ になった。差は $(C_{\mathrm{exact}}-C)+C_{\mathrm{exact}}\delta$ に分かれ、2〜3桁まで合う。$C_{\mathrm{exact}}$ は木の結合を厳密に替えたときの比で、二つの項はどちらも $O(\alpha^2)$ である。$t\to bW^+$ では補正が大きいので、積 $C_{\mathrm{exact}}\delta=-5.4\times10^{-4}$ が主になる。$Z\to e^+e^-$ では、$v_e=-\frac12+2s^2$ が小さいので木の幅が $s^2$ について2次の形になり、$\Delta s^2=\hat s^2(M_Z)-s_0^2=0.0093$（$\Delta\rho$ で大きくなる1ループの差）の2乗が効く。その2次の項 $\frac12(f''/f)(\Delta s^2)^2=1.4\times10^{-3}$ が差のほとんどである。差は期待通り高次の大きさだが、ベクトル結合の小さい荷電レプトンでは $\Gamma_{\ell\ell}$ の測定の誤差（$0.1\%$）と同じ大きさになる。手順 (5) で $\sin^2\theta_{\mathrm{eff}}$ が入力の処方で大きく変わったのと同じ原因である。$t\to bW^+$ の $G_\mu$ の処方の電弱の補正は、1ループで $\delta+C=+1.64\%$ になった。手順 (4c) の $+1.52\%$ とは、比の形で書いたことによる $O(\alpha^2)$ の差と、$M_W$ を極質量にした差（$0.07\%$）だけ違う。

\begin{table}[tbp]
\centering\small
\caption{手順 (7)：$\MSbar$ と $G_\mu$ の処方（\texttt{scheme\_compare.py}）。$\delta$ と $C$ は $\mu=M$ の値（$\%$）。「$\mu$ の幅」は $\mu=M/2$、$M$、$2M$ の最大と最小の比から1を引いた値（$\%$）。最後の2列は $\mu=M$ の差と、その $O(\alpha^2)$ の部分の和 $(C_{\mathrm{exact}}-C)+C_{\mathrm{exact}}\delta$。}
\label{tab:scheme}
\begin{tabular}{lrrrrrrr}
\toprule
 & $\delta$ & $C$ & \multicolumn{3}{c}{$\mu$ の幅} & $\MSbar/G_\mu-1$ & $O(\alpha^2)$ の和 \\
 & & & 木（$\MSbar$） & NLO（$\MSbar$） & NLO（$G_\mu$） & & \\
\midrule
$Z\to\nu\bar\nu$ & $+0.284$ & $+0.643$ & $1.46$ & $0.013$ & $0.056$ & $-1.07\times10^{-4}$ & $-1.08\times10^{-4}$ \\
$Z\to e^+e^-$（弱い部分） & $-0.005$ & $-0.147$ & $1.86$ & $0.038$ & $0.230$ & $+1.21\times10^{-3}$ & $+1.21\times10^{-3}$ \\
$W\to e\nu$ & $+0.011$ & $-0.338$ & $2.39$ & $0.016$ & $0.016$ & $-1.93\times10^{-4}$ & $-1.92\times10^{-4}$ \\
$t\to bW^+$ & $+3.281$ & $-1.643$ & $2.36$ & $0.143$ & $0.108$ & $-4.62\times10^{-4}$ & $-4.69\times10^{-4}$ \\
\bottomrule
\end{tabular}
\end{table}

\FloatBarrier
\paragraph{手順 (8)：予言と不定性のまとめ。}\texttt{summary.py} は、手順 (4a)〜(6) の $\mu=M$ の予言に、NLO の外の項の大きさの見積もりを三つ付ける。一つ目は $\Delta_\mu$ で、$\mu=M/2$、$M$、$2M$（結合と $\alpha_s$ を $\mu$ に走らせる）での NLO の値の幅の半分である。二つ目は $\Delta_{\mathrm{scheme}}$ で、同じ1ループの予言を別の処方で書いた値との差である。幅と比では手順 (7) の $G_\mu$ の処方（$Z$ と $W$ のすべてのチャネルとトップを組み直す）との差、$\sin^2\theta_{\mathrm{eff}}$ と非対称度では SM ノートの $\alpha$ の処方の結合との差（手順 (5)）、$M_W$ では1ループで同等な形と SM ノートの値との差の最大（手順 (6)）をとる。三つ目は $\Delta_{\mathrm{HO}}$ で、$\mu$ では見えない既知の高次の項である。ここでは $M_W$ の $\Delta\rho$ の $O(\alpha\alpha_s)$ の項（$51$~MeV。文献~\cite{dv87} の係数、参照）だけを入れる（$\sin^2\theta_{\mathrm{eff}}$ ではこれが処方の差と同じ大きさなので、重ねて足さない）。理論の不定性はこれらの2乗和とし、測定値との差を、実験の誤差だけで割ったものと、実験と理論の2乗和で割ったものの両方で示す（表~\ref{tab:summary}）。

理論の不定性を入れると、23個の量のうち20個は測定値と $1.5\sigma$ 以内で合う。$\Gamma_{\mathrm{had}}$ は実験の誤差だけでは $-3.1\sigma$ だが、$\mu$ の幅（ほとんどが QCD の因子の $\alpha_s(\mu)$ の走りで、$O(\alpha_s^2)$ の補正の大きさ。手順 (4a)）を入れると $-0.9\sigma$ になる。$R_\ell$ は $-2.6\sigma$ から $-0.8\sigma$ に、$\sin^2\theta_{\mathrm{eff}}$ は $+5.6\sigma$ から $+1.2\sigma$ に、$M_W$ は $+1.9\sigma$（PDG 2026 では $+4.2\sigma$）から $+0.5\sigma$（$+0.6\sigma$）になる。$\sin^2\theta_{\mathrm{eff}}$ と $M_W$ では、理論の不定性（$7.5\times10^{-4}$、$56$~MeV）が実験の誤差の4〜7倍ある。残るのは、見えない幅の $+1.9\sigma$（測定のニュートリノの数が3より小さいこと~\cite{lepewwg06}）、$\mathrm{BR}(W\to\tau\nu)$ の $-2.6\sigma$（LEP の $\tau$ の過剰。LHC の測定~\cite{atlas21} の $\Gamma(W\to\tau\nu)/\Gamma(W\to\mu\nu)=0.992\pm0.013$ では見えていない）、$A_\ell$ の $-1.8\sigma$（$A^{0,\ell}_{\mathrm{FB}}$ も $-1.5\sigma$。SLD の $A_\ell$ と LEP の $A^{0,b}_{\mathrm{FB}}$ の食い違い）で、どれも測定のよく知られた性質（参照）である。

理論の不定性と実験の誤差の比で、量は三つに分かれる。比が $0.31$ 以下で精度が実験で決まるのは、$\Gamma_{\mathrm{inv}}$、$R_b$、$R_c$、$A_b$、$A_c$、$\Gamma_W$、$W$ の分岐比、$\Gamma_t$ である。これらでは、本ノートの NLO で測定の精度まで標準模型を試せている。比が $2.3$〜$3.3$ の $\Gamma_{\mathrm{had}}$、$\Gamma_Z$、$R_\ell$、$\sigma^0_{\mathrm{had}}$ は、主に QCD の因子の $\alpha_s(\mu)$ の走り、すなわち $O(\alpha_s^2)$ の補正で精度が決まる。比が $2.6$〜$7.2$ の $\sin^2\theta_{\mathrm{eff}}$、$A_\ell$、$A^{0,b}_{\mathrm{FB}}$、$M_W$ は、$\Delta\rho$ の高次と、on-shell と $\MSbar$ の混合角の差で精度が決まる。$\Gamma_{\ell\ell}$、$A^{0,\ell}_{\mathrm{FB}}$、$A^{0,c}_{\mathrm{FB}}$ では比が $0.9$〜$1.3$ で、二つが同じ大きさである。

ただし、理論の不定性は決め方（$\mu$ を2倍と半分に振る、比べる処方を選ぶ、高次の項を一つだけ入れる）による目安で、統計的な誤差ではない。表の23個の量は互いに独立でもない（$\Gamma_Z$ は部分幅の和、$R_\ell$ と $\sigma^0_{\mathrm{had}}$ は部分幅の比と積で、同じ $\sin^2\theta_{\mathrm{eff}}$ が非対称度のすべてに入る）。ずれの $\sigma$ の数は目安として読む。

\begin{table}[tbp]
\centering\footnotesize\setlength{\tabcolsep}{3.5pt}
\caption{手順 (8)：NLO の予言と不定性（\texttt{summary.py}。$\mu=M$。GeV、$\sigma^0_{\mathrm{had}}$ は nb）。$\Delta_\mu$、$\Delta_{\mathrm{scheme}}$、$\Delta_{\mathrm{HO}}$ は本文の定義で、合計はその2乗和。測定値は参照（$Z$ は PDG 2024~\cite{pdg24}、非対称度は LEP と SLD~\cite{lepewwg06}、$W$ とトップの幅は PDG の2026年の版~\cite{pdg26,pdg26t}、$M_W$ は Breit--Wigner の質量どうし）。右の2列は測定値との差を、実験の誤差と、実験と理論の2乗和で割ったもの。}
\label{tab:summary}
\begin{tabular}{lllllllrr}
\toprule
量 & NLO & $\Delta_\mu$ & $\Delta_{\mathrm{scheme}}$ & $\Delta_{\mathrm{HO}}$ & 合計 & 測定値 & 実験 & 実験と理論 \\
\midrule
$\Gamma_{\ell\ell}$ & 0.083967 & $1.5\times10^{-5}$ & $1.0\times10^{-4}$ & & $1.0\times10^{-4}$ & $0.083984(86)$ & $-0.20$ & $-0.13$ \\
$\Gamma_{\mathrm{inv}}$ & 0.50183 & $3.3\times10^{-5}$ & $5.4\times10^{-5}$ & & $6.3\times10^{-5}$ & $0.4990(15)$ & $+1.89$ & $+1.89$ \\
$\Gamma_{\mathrm{had}}$ & 1.73827 & $6.5\times10^{-3}$ & $3.3\times10^{-4}$ & & $6.5\times10^{-3}$ & $1.7444(20)$ & $-3.06$ & $-0.90$ \\
$\Gamma_Z$ & 2.49182 & $6.5\times10^{-3}$ & $8.1\times10^{-5}$ & & $6.5\times10^{-3}$ & $2.4955(23)$ & $-1.60$ & $-0.54$ \\
$R_\ell$ & 20.7019 & $0.078$ & $0.029$ & & $0.084$ & $20.767(25)$ & $-2.61$ & $-0.75$ \\
$R_b$ & 0.216101 & $4.7\times10^{-5}$ & $5.6\times10^{-6}$ & & $4.8\times10^{-5}$ & $0.21629(66)$ & $-0.29$ & $-0.29$ \\
$R_c$ & 0.171039 & $8.5\times10^{-6}$ & $1.9\times10^{-5}$ & & $2.1\times10^{-5}$ & $0.1721(30)$ & $-0.35$ & $-0.35$ \\
$\sigma^0_{\mathrm{had}}$ & 41.4979 & $0.063$ & $0.045$ & & $0.077$ & $41.481(33)$ & $+0.51$ & $+0.20$ \\
$\sin^2\theta^{\mathrm{lept}}_{\mathrm{eff}}$ & 0.232423 & $3.6\times10^{-5}$ & $7.5\times10^{-4}$ & & $7.5\times10^{-4}$ & $0.23153(16)$ & $+5.58$ & $+1.16$ \\
$A_\ell$ & 0.13992 & $2.8\times10^{-4}$ & $5.9\times10^{-3}$ & & $5.9\times10^{-3}$ & $0.1513(21)$ & $-5.42$ & $-1.80$ \\
$A_b$ & 0.93404 & $1.6\times10^{-5}$ & $4.9\times10^{-4}$ & & $4.9\times10^{-4}$ & $0.923(20)$ & $+0.55$ & $+0.55$ \\
$A_c$ & 0.66475 & $1.3\times10^{-4}$ & $2.6\times10^{-3}$ & & $2.6\times10^{-3}$ & $0.670(27)$ & $-0.19$ & $-0.19$ \\
$A^{0,\ell}_{\mathrm{FB}}$ & 0.014684 & $5.9\times10^{-5}$ & $1.3\times10^{-3}$ & & $1.3\times10^{-3}$ & $0.0171(10)$ & $-2.42$ & $-1.49$ \\
$A^{0,b}_{\mathrm{FB}}$ & 0.09802 & $2.0\times10^{-4}$ & $4.2\times10^{-3}$ & & $4.2\times10^{-3}$ & $0.0992(16)$ & $-0.74$ & $-0.26$ \\
$A^{0,c}_{\mathrm{FB}}$ & 0.06976 & $1.5\times10^{-4}$ & $3.2\times10^{-3}$ & & $3.3\times10^{-3}$ & $0.0707(35)$ & $-0.27$ & $-0.20$ \\
$\Gamma_W$ & 2.08594 & $5.3\times10^{-3}$ & $5.7\times10^{-4}$ & & $5.4\times10^{-3}$ & $2.14(5)$ & $-1.08$ & $-1.08$ \\
$\mathrm{BR}(W\to e\nu)$ & 0.108441 & $2.8\times10^{-4}$ & $8.9\times10^{-6}$ & & $2.8\times10^{-4}$ & $0.1071(16)$ & $+0.84$ & $+0.83$ \\
$\mathrm{BR}(W\to\mu\nu)$ & 0.108441 & $2.8\times10^{-4}$ & $8.9\times10^{-6}$ & & $2.8\times10^{-4}$ & $0.1063(15)$ & $+1.43$ & $+1.40$ \\
$\mathrm{BR}(W\to\tau\nu)$ & 0.108364 & $2.8\times10^{-4}$ & $8.9\times10^{-6}$ & & $2.8\times10^{-4}$ & $0.1138(21)$ & $-2.59$ & $-2.57$ \\
$\mathrm{BR}(W\to\mathrm{had})$ & 0.674753 & $8.4\times10^{-4}$ & $2.7\times10^{-5}$ & & $8.4\times10^{-4}$ & $0.6741(27)$ & $+0.24$ & $+0.23$ \\
$\Gamma_t$ & 1.3776 & $0.011$ & $1.4\times10^{-3}$ & & $0.011$ & $1.42^{+0.19}_{-0.15}$ & $-0.28$ & $-0.28$ \\
$M_W$（PDG 2024） & 80.395 & $0.017$ & $0.014$ & $0.051$ & $0.056$ & $80.3692(133)$ & $+1.94$ & $+0.45$ \\
$M_W$（PDG 2026） & 80.395 & $0.017$ & $0.014$ & $0.051$ & $0.056$ & $80.3625(77)$ & $+4.22$ & $+0.58$ \\
\bottomrule
\end{tabular}
\end{table}

\FloatBarrier
\section{議論と今後}\label{sec:discussion}

\paragraph{問いへの答え。}第一の問い（赤外の極の打ち消し）には、手順 (2)〜(4) で答えた。質量0のフェルミオン対では $\epsilon$ について厳密に、質量のある放射体では軟らかい部分の解析的な積分と硬い部分の数値の積分で、極が打ち消すことを確かめた。光子が $W$ と崩壊するトップにも付く場合も含め、結果はゲージによらず、切断 $\delta$ と放射体の質量（KLN）にもよらなかった。MadLoop の仮想の補正とは $10^{-11}$〜$10^{-15}$、MadGraph の木の3体の振幅の2乗とは $10^{-10}$〜$10^{-12}$ で一致した。第二の問い（NLO の幅と補正の大きさ）には手順 (4) で答えた。$\mu=M$ の $\MSbar$ の結合では電弱の補正は小さく（$Z$ で $-0.01$〜$+0.3\%$、$Z\to b\bar b$ は $-1.44\%$。$W$ で $+0.01$〜$-0.04\%$、$b$ を含むチャネルは $+0.76\%$）、トップでは $+3.28\%$（$G_\mu$ の処方で $+1.64\%$）である。QCD の補正はハドロンの幅で約 $\alpha_s/\pi$、トップで $-8.5\%$ である。電弱の部分の $\mu$ による変化は $\pm1.2\%$ から $0.04\%$ 以下になった。第三の問い（$Z$ 極の量と $M_W$）には手順 (5)、(6) で答えた。$R_b$ は $Zb\bar b$ の頂点のトップの補正で測定値と合う。$\sin^2\theta_{\mathrm{eff}}$ と $M_W$ の1ループの予言の不定性は測定の誤差の4〜7倍で、その範囲で測定値と合う。

\paragraph{分かったこと。}$\MSbar$ と FJ のタドポールの処方では、個々の量（$dW$、$dZ$、$\hat\Delta$）にタドポールの大きな部分（約 $12\%$）が入るが、物理量の組では記号として消える。そのため関係式は線形に書く必要があり、積の形の切り捨ての違いは $O(\alpha^2)$ の差として $M_W$ に $10$〜$20$~MeV 残る（手順 (6)）。$\sin^2\theta_{\mathrm{eff}}$、$\Gamma_{\ell\ell}$、$M_W$ が処方に敏感なのは、on-shell の $s_W^2$ と $\MSbar$ の $\hat s^2$ が $\Delta\rho$ のために $4\%$ 違い、その2次や $\Delta\rho$ の高次が測定の精度で見えるためである（手順 (5)〜(7)）。1ループの不定性が測定の誤差より小さいのは、$\Gamma_{\mathrm{inv}}$、$R_b$、$R_c$、$A_b$、$A_c$、$\Gamma_W$、$W$ の分岐比、$\Gamma_t$ である。$\Gamma_{\mathrm{had}}$、$\Gamma_Z$、$R_\ell$、$\sigma^0_{\mathrm{had}}$ では QCD の高次が、$\sin^2\theta_{\mathrm{eff}}$、$A_\ell$、$A^{0,b}_{\mathrm{FB}}$、$M_W$ では電弱の高次が測定の誤差を超える（手順 (8)）。道具の扱いでは、Feynman ゲージの近似 $\xi=1+10^{-20}$ が質量0のニュートリノの脚で80桁を使い切って値を誤らせること（警告は出ない。KLN の確かめで見つかった）と、MadGraph が UFO の $d$ 型の電荷を7桁で書き出すことが分かった（どちらも計算環境の記録~\cite{env} に加えた）。計画の式~\eqref{eq:masslessVR} の有限部分の規格化の誤りも手順 (2) で直した。

\paragraph{今後。}$\sin^2\theta_{\mathrm{eff}}$ と $M_W$ を測定の精度で予言するには、$\Delta\rho$ の $O(\alpha\alpha_s)$ と $O(\alpha^2m_t^4)$ を中心とする2ループの補正と、それに合わせた2ループの入力の換算が要る。ハドロンの幅には $O(\alpha_s^2)$ の QCD の補正（約 $+0.2\%$）と $O(\alpha\alpha_s)$ の補正、$W\to c\bar s$ などには QCD の因子のクォークの質量の効果（$\Gamma_W$ の $10^{-4}$ 程度）、トップの幅には $O(\alpha_s^2)$ と CKM 行列の換算（CKM ノートの手順 (6)）が残っている。また本ノートは運動学の質量に測定の Breit--Wigner の値をそのまま使った（極質量との差は $4\times10^{-4}$）。$M_W$ の比較（手順 (6)）では極質量に直した。幅では $\Gamma/M$ が二つの決め方で同じなので、$M$ に比例する幅に Breit--Wigner の質量を使うと Breit--Wigner の決め方の幅になり、測定値と同じ決め方で比べたことになる。

\appendix

\section{木の崩壊幅の導出}\label{app:lo}

$Z\to f\bar f$ の振幅 $-i\frac{g_Z}2\bar u\gamma^\mu(v-a\gamma_5)v\,\varepsilon_\mu$ の2乗を、$Z$ の偏極で平均し（$-g_{\mu\nu}+p_\mu p_\nu/M_Z^2$。$p_\mu$ の項はベクトルの部分では消え、軸性ベクトルの部分では $m_f^2$ に比例する）、スピンと色で和をとると、$\frac{N_cg_Z^2}{3}\bigl[v^2(M_Z^2+2m_f^2)+a^2(M_Z^2-4m_f^2)\bigr]$ になる。2体の位相空間 $\int\dd\Phi_2=\beta/(8\pi)$ と $1/(2M_Z)$ を掛けると、式~\eqref{eq:lo} の1行目になる（$M_Z^2+2m_f^2=M_Z^2(3-\beta^2)/2$、$M_Z^2-4m_f^2=M_Z^2\beta^2$）。$W^+\to f\bar f'$ は $v=a=1$、$g_Z\to g/\sqrt2$ の左巻きの結合で、異なる質量の2体の位相空間 $\lambda^{1/2}/(8\pi)$ から2行目になる。$t\to bW^+$ は $W$ の偏極の和の $p_\mu p_\nu/M_W^2$ の項が縦波の寄与を与え、3行目の $(1-x_W)^2(1+2x_W)$（$m_b=0$）になる。$H\to f\bar f$ は湯川の頂点 $-iy_f/\sqrt2$ から $\sum|\mathcal M|^2=2N_cy_f^2(M_H^2-4m_f^2)/2$ で、4行目になる。

\section{赤外の極の打ち消し}\label{app:ir}

質量0のフェルミオン対では、$x_i=2E_i/M$（崩壊する粒子の静止系。$x_1+x_2+x_3=2$、3が光子またはグルーオン）として、実の放射の振幅の2乗は $d$ 次元で
\begin{equation*}
r\equiv\frac{M^2\sum|\mathcal M_R|^2}{e^2Q_f^2\sum|\mathcal M_0|^2}=\frac{2\bigl[x_1^2+x_2^2-\epsilon(2-x_1-x_2)^2\bigr]}{(1-x_1)(1-x_2)}
\end{equation*}
になる（ベクトルの流れ。軸性ベクトルでも質量0では同じ。グルーオンでは $e^2Q_f^2\to g_s^2C_F$）。$y_{ij}=s_{ij}/M^2$（$y_{23}=1-x_1$、$y_{13}=1-x_2$、$y_{12}=1-x_3$）で、$d$ 次元の3体と2体の位相空間は
\begin{equation*}
\int\dd\Phi_3=\frac{M^2}{2(4\pi)^3}\Bigl(\frac{4\pi}{M^2}\Bigr)^{2\epsilon}\frac1{\Gamma(2-2\epsilon)}\int\dd y_{13}\dd y_{23}\,(y_{12}y_{13}y_{23})^{-\epsilon},\qquad
\int\dd\Phi_2=\Bigl(\frac{4\pi}{M^2}\Bigr)^\epsilon\frac{\Gamma(1-\epsilon)}{8\pi\Gamma(2-2\epsilon)}
\end{equation*}
である。$r$ を $y_{12}$、$y_{13}$、$y_{23}$ の単項式の和 $2[y_{13}/y_{23}+y_{23}/y_{13}+2y_{12}/(y_{13}y_{23})-\epsilon(y_{13}/y_{23}+y_{23}/y_{13}+2)]$ に書くと、積分は Dirichlet の積分 $\int y_{12}^{a-1}y_{13}^{b-1}y_{23}^{c-1}=\Gamma(a)\Gamma(b)\Gamma(c)/\Gamma(a+b+c)$ になる。$y_{13},y_{23}\to0$ の端が $1/\epsilon^2$ と $1/\epsilon$ を与える。仮想の補正は、質量0のフェルミオンの流れの1ループの形状因子である（光子またはグルーオンの交換。場の定数はスケールのない積分で0）。TID で $B_0(M^2;0,0)$ と $C_0(0,0,M^2;0,0,0)=2(d-3)B_0/((4-d)M^2)$ に帰着し、$\mathcal M_1/\mathcal M_0=\frac{\alpha C}{4\pi}\bigl[-(2-\epsilon+2\epsilon^2)/\epsilon\bigr]B_0$、$B_0=(4\pi\mu^2/M^2)^\epsilon\Gamma(\epsilon)\Gamma(1-\epsilon)^2e^{i\pi\epsilon}/\Gamma(2-2\epsilon)$ になる。結果は $\epsilon$ について厳密に
\begin{equation}
\begin{aligned}
\frac{\Gamma_R}{\Gamma_0}&=\frac{\alpha C}{2\pi}\Bigl(\frac{4\pi\mu^2}{M^2}\Bigr)^\epsilon\frac{6(1-\epsilon)(1-2\epsilon)(2-2\epsilon+\epsilon^2)\Gamma(1-\epsilon)^2}{\epsilon^2\,\Gamma(4-3\epsilon)},\\
\frac{\Gamma_V}{\Gamma_0}&=-\frac{\alpha C}{2\pi}\Bigl(\frac{4\pi\mu^2}{M^2}\Bigr)^\epsilon\frac{(2-\epsilon+2\epsilon^2)\Gamma(1+\epsilon)\Gamma(1-\epsilon)^2\cos\pi\epsilon}{\epsilon^2\,\Gamma(2-2\epsilon)}
\end{aligned}
\label{eq:irexact}
\end{equation}
で（$\cos\pi\epsilon$ は $(-M^2-i0)^{-\epsilon}$ の位相から）、$1/\Gamma(1-\epsilon)$ を除いて展開すると式~\eqref{eq:masslessVR} になる。和は $N_\epsilon\cdot\frac32+O(\epsilon)$ で、式~\eqref{eq:inclusive} の $\frac34\frac{\alpha C}\pi$ である。

\begin{thebibliography}{99}
\bibitem{smnote} 標準模型：$\MSbar$ の電弱の1ループの繰り込みと検算（$V_{\mathrm{CKM}}=1$、FJ のタドポール、一般の $R_\xi$ ゲージ、$d=4$）（本リポジトリ \nolinkurl{models/sm/ResearchNote/note.pdf}、2026年10月17日の版）。式、表、手順の番号はこの版のもの。
\bibitem{ckmnote} 標準模型と CKM 行列：$\MSbar$ の繰り込みと世代を変える過程（FJ のタドポール、一般の $R_\xi$ ゲージ、$d=4$）（本リポジトリ \nolinkurl{models/sm_ckm/ResearchNote/note.pdf}、2026年10月18日の版）。式、表、手順の番号はこの版のもの。
\bibitem{bn37} F. Bloch, A. Nordsieck, ``Note on the radiation field of the electron,'' Phys. Rev. \textbf{52}, 54 (1937)。
\bibitem{kinoshita62} T. Kinoshita, ``Mass singularities of Feynman amplitudes,'' J. Math. Phys. \textbf{3}, 650 (1962)。
\bibitem{ln64} T. D. Lee, M. Nauenberg, ``Degenerate systems and mass singularities,'' Phys. Rev. \textbf{133}, B1549 (1964)。
\bibitem{lepewwg06} ALEPH, DELPHI, L3, OPAL, SLD Collaborations, LEP Electroweak Working Group, SLD Electroweak and Heavy Flavour Groups, ``Precision electroweak measurements on the $Z$ resonance,'' Phys. Rept. \textbf{427}, 257 (2006), arXiv:hep-ex/0509008。
\bibitem{pdg24} S. Navas \textit{et al.} (Particle Data Group), Phys. Rev. D \textbf{110}, 030001 (2024)。
\bibitem{pdg26t} Particle Data Group, pdgLive（2026年の版）のトップクォークの項、\url{https://pdglive.lbl.gov/Particle.action?node=Q007}（2026年10月19日に参照）。
\bibitem{atlas21} ATLAS Collaboration, ``Test of the universality of $\tau$ and $\mu$ lepton couplings in $W$-boson decays with the ATLAS detector,'' Nature Phys. \textbf{17}, 813 (2021), arXiv:2007.14040。
\bibitem{dv87} A. Djouadi, C. Verzegnassi, ``Virtual very heavy top effects in LEP/SLC precision measurements,'' Phys. Lett. B \textbf{195}, 265 (1987)。
\bibitem{ds91} A. Denner, T. Sack, ``The top width,'' Nucl. Phys. B \textbf{358}, 46 (1991)。
\bibitem{pdg26} Particle Data Group, pdgLive（2026年の版）の $W$ の項、\url{https://pdglive.lbl.gov/Particle.action?node=S043}（2026年10月19日に参照）。
\bibitem{mr19} S. P. Martin, D. G. Robertson, ``Standard Model parameters in the tadpole-free pure $\overline{\mathrm{MS}}$ scheme,'' Phys. Rev. D \textbf{100}, 073004 (2019), arXiv:1907.02500。
\bibitem{cdt01} S. Catani, S. Dittmaier, Z. Trocs\'anyi, ``One-loop singular behaviour of QCD and SUSY QCD amplitudes with massive partons,'' Phys. Lett. B \textbf{500}, 149 (2001), arXiv:hep-ph/0011222。
\bibitem{esw96} R. K. Ellis, W. J. Stirling, B. R. Webber, \textit{QCD and Collider Physics} (Cambridge University Press, 1996), Sec. 3.4。
\bibitem{veltman77} M. Veltman, ``Limit on mass differences in the Weinberg model,'' Nucl. Phys. B \textbf{123}, 89 (1977)。
\bibitem{env} ループ計算環境の構築記録（本リポジトリ \nolinkurl{environment/environment.pdf}）。
\bibitem{jk89} M. Je\.zabek, J. H. K\"uhn, ``QCD corrections to semileptonic decays of heavy quarks,'' Nucl. Phys. B \textbf{314}, 1 (1989)。
\bibitem{cm98} A. Czarnecki, K. Melnikov, ``Two-loop QCD corrections to top quark width,'' Nucl. Phys. B \textbf{544}, 520 (1999), arXiv:hep-ph/9806244。
\end{thebibliography}

\end{document}
